% Communication — Physics Upper Sixth — Cours signé OnBuch+
% Official MINESEC syllabus. Compiler avec : tectonic PHY14-communication.tex
\newcommand{\DOCMATIERE}{Physics}
\newcommand{\DOCNIVEAU}{Upper Sixth}
\newcommand{\DOCMODULE}{Upper Sixth — Physics}
\newcommand{\DOCLECON}{14}
\newcommand{\DOCTITRE}{Communication}
% OnBuch+ — préambule « cours signés » ENRICHI (compilé avec tectonic / XeLaTeX)
% Système visuel « Pop » d'OnBuch+ : fond crème (jamais blanc), bordures encre
% épaisses, ombres « dures » décalées, titres Archivo Black en capitales.
% Version MATHÉMATIQUES : graphes (pgfplots/tikz), boîtes pédagogiques
% (définition, propriété, méthode, attention, le savais-tu, exercice résolu,
% exercice, TP, bilan), page de garde et signature OnBuch+ ancrée en bas.
%
% Macros attendues dans main.tex AVANT \input{preamble} :
%   \DOCMATIERE  \DOCNIVEAU  \DOCMODULE  \DOCLECON (numéro)  \DOCTITRE
\documentclass[11pt]{article}

\usepackage{fontspec}
\usepackage[english]{babel}
\usepackage[a4paper,margin=2cm,top=3.4cm,bottom=2.8cm,headheight=1.4cm,footskip=1.3cm]{geometry}
\usepackage{amsmath,amssymb,amsfonts}
\usepackage{mathtools} % \xmapsto, \coloneqq... souvent utilisés par les modèles
\usepackage{cancel} % simplifications barrées (bilans de forces)
% \abs{z} (module, valeur absolue) et \norm{u} (norme), utilisés par les modèles
\providecommand{\abs}{}\let\abs\relax\DeclarePairedDelimiter{\abs}{\lvert}{\rvert}
\providecommand{\norm}{}\let\norm\relax\DeclarePairedDelimiter{\norm}{\lVert}{\rVert}
\usepackage[table]{xcolor} % [table] : fournit aussi \rowcolors (alternance de lignes)
\usepackage{pagecolor}
\usepackage{tikz}
\usetikzlibrary{arrows.meta,positioning,calc,shapes.geometric,shapes.misc,shapes.arrows,decorations.pathmorphing,decorations.pathreplacing,decorations.markings,patterns,shadows,mindmap,trees,fit,backgrounds,matrix,intersections,angles,quotes}
\usepackage{pgfplots}
\usepackage[version=4]{mhchem} % équations nucléaires \ce{^{235}_{92}U}
\usepackage[european,siunitx]{circuitikz} % schémas électriques
\pgfplotsset{compat=1.18}
\usepgfplotslibrary{fillbetween,groupplots} % aires entre courbes (calcul intégral)
\usepackage{enumitem}
\usepackage{fancyhdr}
% [most] charge toutes les bibliothèques tcolorbox (listings, minted, raster,
% documentation...) et gonfle inutilement la mémoire TeX sur un document long
% (~300 boîtes) : on ne charge que ce qui est réellement utilisé.
\usepackage[skins,breakable]{tcolorbox}
\usepackage{graphicx}
\usepackage{colortbl}
\usepackage{booktabs}
\usepackage{tabularx}
\usepackage{diagbox} % case d'en-tête diagonale des tableaux à double entrée
% Colonnes extensibles centrée / gauche / droite (C, L, R), souvent utilisées par les modèles
\newcolumntype{C}{>{\centering\arraybackslash}X}
\newcolumntype{L}{>{\raggedright\arraybackslash}X}
\newcolumntype{R}{>{\raggedleft\arraybackslash}X}
\usepackage{array}
\usepackage{multicol}
\usepackage{siunitx}
% parse-numbers=false : accepte aussi \SI{2,5 \times 10^{-3}}{...}
\sisetup{locale=UK,output-decimal-marker={.},group-minimum-digits=4,parse-numbers=false}
\DeclareSIUnit\ohm{\ensuremath{\symup{\Omega}}} % U+2126 absent de la police
\DeclareSIUnit\division{div} % réglages d'oscilloscope (ms/div, V/div)
\DeclareSIUnit\tr{tr} % tours (vitesse de rotation, ex. \SI{1500}{\tr\per\minute})
\usepackage{qrcode}
% \degree : commande fréquemment utilisée par les modèles pour un angle ou une
% température en dehors de \SI{}{} (non fournie par les paquets chargés).
\providecommand{\degree}{^{\circ}}
% \qed (fin de démonstration), utilisé par les modèles sans amsthm
\providecommand{\qed}{\hfill$\square$}
% \barre : double barre des valeurs interdites dans un tableau de variations (tkz-tab)
\providecommand{\barre}{\ensuremath{\|}}
% environnement proof (amsthm non chargé) utilisé par les modèles
\ifdefined\proof\else\newenvironment{proof}[1][Proof]{\par\noindent\textit{#1.}\ }{\qed\par}\fi
\usepackage{lastpage}
\usepackage{hyperref}
\hypersetup{hidelinks}

% ============================================================
% Palette « Pop » OnBuch+ — mêmes hex que la classe P (plus_ui.dart)
% ============================================================
\definecolor{poporange}{HTML}{F59321}
\definecolor{popdark}{HTML}{E07A0C}
\definecolor{popcream}{HTML}{FFF6E8}
\definecolor{popink}{HTML}{1C1714}
\definecolor{poppurple}{HTML}{7A5AE0}
\definecolor{popgreen}{HTML}{2E9E6A}
\definecolor{poppink}{HTML}{E0457B}
\definecolor{popblue}{HTML}{2F7DE1}
\definecolor{popgold}{HTML}{C98A12}
\definecolor{popmuted}{HTML}{8A7F76}
\definecolor{popline}{HTML}{EADFD2}
\definecolor{popgray}{HTML}{8A7F76}
\colorlet{popgrey}{popgray}
% Alias tolérants (le modèle nomme parfois les couleurs différemment)
\colorlet{popwhite}{white}
\colorlet{popblack}{popink}
\colorlet{popred}{poppink}
\colorlet{popyellow}{popgold}
% Teintes claires pour fonds de tableaux / zones
\colorlet{popblueL}{popblue!10!white}
\colorlet{poporangeL}{poporange!14!white}
\colorlet{popgreenL}{popgreen!12!white}
\colorlet{poppurpleL}{poppurple!12!white}
\colorlet{poppinkL}{poppink!12!white}
\colorlet{popgoldL}{popgold!14!white}

\pagecolor{popcream}

% ============================================================
% Typographie
% ============================================================
\defaultfontfeatures{Ligatures=TeX}
\setmainfont{PlusJakartaSans-Medium.ttf}[Path=../../../fonts/,BoldFont=PlusJakartaSans-Bold.ttf]
% Maths en sans-serif (Fira Math) avec chiffres et lettres droites de Plus
% Jakarta, pour que les formules se fondent dans le texte.
\usepackage[math-style=ISO,bold-style=ISO]{unicode-math}
\setmathfont{FiraMath-Regular.otf}[Path=../../../fonts/]
\setmathfont{PlusJakartaSans-Medium.ttf}[Path=../../../fonts/,range={up/num,up/{Latin,latin},\mathup}]
% (icomma retiré : décimales avec point en anglais)
% Caractères mathématiques Unicode absents des polices de texte (Plus Jakarta,
% Archivo Black) : rendus par leur équivalent mathématique, en texte comme en
% formule — sinon ils s'affichent en carré vide (ex. « Arithmétique dans ℤ »).
\usepackage{newunicodechar}
\newunicodechar{ℝ}{\ensuremath{\mathbb{R}}}
\newunicodechar{ℤ}{\ensuremath{\mathbb{Z}}}
\newunicodechar{ℂ}{\ensuremath{\mathbb{C}}}
\newunicodechar{ℕ}{\ensuremath{\mathbb{N}}}
\newunicodechar{ℚ}{\ensuremath{\mathbb{Q}}}
\newunicodechar{𝔻}{\ensuremath{\mathbb{D}}}
\newunicodechar{α}{\ensuremath{\alpha}}
\newunicodechar{β}{\ensuremath{\beta}}
\newunicodechar{γ}{\ensuremath{\gamma}}
\newunicodechar{δ}{\ensuremath{\delta}}
\newunicodechar{ε}{\ensuremath{\varepsilon}}
\newunicodechar{θ}{\ensuremath{\theta}}
\newunicodechar{λ}{\ensuremath{\lambda}}
\newunicodechar{μ}{\ensuremath{\mu}}
\newunicodechar{σ}{\ensuremath{\sigma}}
\newunicodechar{τ}{\ensuremath{\tau}}
\newunicodechar{φ}{\ensuremath{\varphi}}
\newunicodechar{ψ}{\ensuremath{\psi}}
\newunicodechar{ω}{\ensuremath{\omega}}
\newunicodechar{Σ}{\ensuremath{\Sigma}}
% majuscules produites par \MakeUppercase dans les titres (α-aminés -> Α)
\newunicodechar{Α}{\ensuremath{\alpha}}
\newunicodechar{Β}{\ensuremath{\beta}}
\newunicodechar{Γ}{\ensuremath{\Gamma}}
\newunicodechar{Δ}{\ensuremath{\Delta}}
\newunicodechar{Ε}{\ensuremath{\varepsilon}}
\newunicodechar{Θ}{\ensuremath{\Theta}}
\newunicodechar{Λ}{\ensuremath{\Lambda}}
\newunicodechar{Μ}{\ensuremath{\mu}}
\newunicodechar{Τ}{\ensuremath{\tau}}
\newunicodechar{Φ}{\ensuremath{\Phi}}
\newunicodechar{Ψ}{\ensuremath{\Psi}}
\newunicodechar{Ω}{\ensuremath{\Omega}}
\newunicodechar{↦}{\ensuremath{\mapsto}}
\newunicodechar{∈}{\ensuremath{\in}}
\newunicodechar{∉}{\ensuremath{\notin}}
\newunicodechar{∧}{\ensuremath{\wedge}}
\newunicodechar{⇔}{\ensuremath{\Leftrightarrow}}
\newunicodechar{⟺}{\ensuremath{\Longleftrightarrow}}
\newunicodechar{⇒}{\ensuremath{\Rightarrow}}
\newunicodechar{⟹}{\ensuremath{\Longrightarrow}}
\newunicodechar{∘}{\ensuremath{\circ}}
\newunicodechar{∪}{\ensuremath{\cup}}
\newunicodechar{∩}{\ensuremath{\cap}}
\newunicodechar{≡}{\ensuremath{\equiv}}
\newunicodechar{⊂}{\ensuremath{\subset}}
\newunicodechar{⊥}{\ensuremath{\perp}}
\newunicodechar{∃}{\ensuremath{\exists}}
\newunicodechar{∀}{\ensuremath{\forall}}
\newunicodechar{∛}{\ensuremath{\sqrt[3]{\,}}}
\newunicodechar{∜}{\ensuremath{\sqrt[4]{\,}}}
\newunicodechar{⃗}{\ensuremath{{}^{\to}}}
\newunicodechar{ⁿ}{\ifmmode^{n}\else\textsuperscript{\ensuremath{n}}\fi}
\newunicodechar{ˣ}{\ifmmode^{x}\else\textsuperscript{\ensuremath{x}}\fi}
\newunicodechar{ᵇ}{\ifmmode^{b}\else\textsuperscript{\ensuremath{b}}\fi}
\newunicodechar{ʳ}{\ifmmode^{r}\else\textsuperscript{\ensuremath{r}}\fi}
\newunicodechar{ᵘ}{\ifmmode^{u}\else\textsuperscript{\ensuremath{u}}\fi}
\newunicodechar{ʸ}{\ifmmode^{y}\else\textsuperscript{\ensuremath{y}}\fi}
\newunicodechar{ᵃ}{\ifmmode^{a}\else\textsuperscript{\ensuremath{a}}\fi}
\newunicodechar{ᵉ}{\ifmmode^{e}\else\textsuperscript{\ensuremath{e}}\fi}
\newunicodechar{ᵏ}{\ifmmode^{k}\else\textsuperscript{\ensuremath{k}}\fi}
\newunicodechar{ᵖ}{\ifmmode^{p}\else\textsuperscript{\ensuremath{p}}\fi}
\newunicodechar{ᴺ}{\ifmmode^{N}\else\textsuperscript{\ensuremath{N}}\fi}
\newunicodechar{ᶜ}{\ifmmode^{c}\else\textsuperscript{\ensuremath{c}}\fi}
\newunicodechar{ʰ}{\ifmmode^{h}\else\textsuperscript{\ensuremath{h}}\fi}
\newunicodechar{ᵠ}{\ifmmode^{q}\else\textsuperscript{\ensuremath{q}}\fi}
\newunicodechar{ₙ}{\ifmmode_{n}\else\textsubscript{\ensuremath{n}}\fi}
\newunicodechar{ₐ}{\ifmmode_{a}\else\textsubscript{\ensuremath{a}}\fi}
\newunicodechar{ᵢ}{\ifmmode_{i}\else\textsubscript{\ensuremath{i}}\fi}
\newunicodechar{ᵣ}{\ifmmode_{r}\else\textsubscript{\ensuremath{r}}\fi}
\newunicodechar{ₖ}{\ifmmode_{k}\else\textsubscript{\ensuremath{k}}\fi}
\newunicodechar{ᵦ}{\ifmmode_{\beta}\else\textsubscript{\ensuremath{\beta}}\fi}
\newunicodechar{ₓ}{\ifmmode_{x}\else\textsubscript{\ensuremath{x}}\fi}
\newfontfamily\popdisplay{ArchivoBlack-Regular.ttf}[Path=../../../fonts/]
\newfontfamily\poplogo{SpaceGrotesk-Bold.ttf}[Path=../../../fonts/]
\newfontfamily\popsemi{PlusJakartaSans-SemiBold.ttf}[Path=../../../fonts/]
\newfontfamily\popxbold{PlusJakartaSans-ExtraBold.ttf}[Path=../../../fonts/]
\newfontfamily\popxboldls{PlusJakartaSans-ExtraBold.ttf}[Path=../../../fonts/,LetterSpace=3.0]

\setlist[enumerate]{leftmargin=1.7em,itemsep=5pt,topsep=4pt}
\setlist[itemize]{leftmargin=1.7em,itemsep=4pt,topsep=4pt,label={\color{poporange}\textbullet}}
\setlength{\parindent}{0pt}
\setlength{\parskip}{5pt}
\renewcommand{\arraystretch}{1.25}

% pgfplots : style Pop par défaut
\pgfplotsset{
  every axis/.append style={axis line style={popink,line width=1pt},tick style={popink},
    grid style={popline,line width=0.5pt},label style={font=\small},tick label style={font=\footnotesize}},
  popcurve/.style={poporange,line width=1.8pt,no markers,smooth},
}
% Même style disponible dans un \draw TikZ simple (hors pgfplots)
\tikzset{popcurve/.style={poporange,line width=1.8pt}}
\tikzset{popgrid/.style={popline,very thin}}
\colorlet{popcurve}{poporange} % le nom du style est parfois utilisé comme couleur

% ============================================================
% Pill / squiggle
% ============================================================
\newcommand{\poppill}[3]{%
  \tikz[baseline=-0.55ex]\node[draw=popink,line width=1.2pt,rounded corners=2.6pt,%
    fill=#2,inner xsep=6.5pt,inner ysep=3pt]{%
    \popxboldls\fontsize{7}{7}\selectfont\color{#3}\MakeUppercase{#1}};%
}
\newcommand{\popsquiggle}[1]{%
  \tikz[baseline=-0.4ex]{
    \draw[#1,line width=2.6pt,line cap=round]
      plot [smooth] coordinates {(0,0.09) (0.34,0.01) (0.68,0.21) (1.02,0.01) (1.36,0.10)};
  }%
}
% Mot-clé surligné (marqueur orange sous le texte)
\newcommand{\cle}[1]{{\popxbold\color{popdark}#1}}

% ============================================================
% En-tête / pied
% ============================================================
\pagestyle{fancy}
\fancyhf{}
\renewcommand{\headrulewidth}{0pt}
\renewcommand{\footrulewidth}{0pt}
\lhead{\raisebox{-0.15ex}{\poplogo\fontsize{13}{13}\selectfont\color{popink}OnBuch\color{poporange}+}}
\rhead{\poppill{\DOCMATIERE\ \textbullet\ \DOCNIVEAU\ \textbullet\ Chapter \DOCLECON}{white}{popink}}
\lfoot{\popxbold\fontsize{7.6}{7.6}\selectfont\color{popmuted}\MakeUppercase{Signed course OnBuch+}\ \textbullet\ {\popsemi\DOCTITRE}}
\rfoot{\tikz[baseline=-0.6ex]\node[rounded corners=8.5pt,fill=popink,minimum height=17pt,minimum width=17pt,inner xsep=5pt,inner sep=0]{\popxbold\fontsize{8}{8}\selectfont\color{popcream}\thepage\,/\,\pageref*{LastPage}};}

% ============================================================
% Titres
% ============================================================
\newcommand{\chaptitre}[2]{%
  \begin{tcolorbox}[enhanced,colback=poporange,colframe=popink,boxrule=2.6pt,arc=11pt,
      drop shadow=popink,left=15pt,right=15pt,top=12pt,bottom=11pt,before skip=3pt,after skip=18pt]
    \poppill{#2}{popink}{popcream}\par\vspace{9pt}
    {\raggedright\hyphenpenalty=10000\exhyphenpenalty=10000\popdisplay\fontsize{22}{24}\selectfont\color{popcream}\MakeUppercase{#1}\par}
  \end{tcolorbox}
}
% Section numérotée : pastille orange avec le numéro + titre Archivo
\newcounter{popsec}
\newcommand{\coursec}[1]{%
  \stepcounter{popsec}\setcounter{popsub}{0}%
  \par\vspace{16pt}\noindent
  \begin{minipage}{\linewidth}
  \tikz[baseline=-0.7ex]\node[draw=popink,line width=1.6pt,fill=poporange,rounded corners=4pt,minimum size=22pt,inner sep=0]{\popdisplay\fontsize{12}{12}\selectfont\color{popcream}\thepopsec};%
  \hspace{8pt}{\popdisplay\fontsize{15}{17}\selectfont\color{popink}\MakeUppercase{#1}}\par
  \vspace{4pt}\noindent{\color{poporange}\rule{36pt}{3.6pt}}
  \end{minipage}\par\nopagebreak\vspace{8pt}
}
\newcounter{popsub}
\newcommand{\courssub}[1]{%
  \stepcounter{popsub}%
  \par\vspace{9pt}\noindent{\popxbold\fontsize{11.5}{13}\selectfont\color{popink}{\color{poporange}\thepopsec.\thepopsub}\hspace{6pt}#1}\par\nopagebreak\vspace{4pt}
}

% ============================================================
% Boîtes pédagogiques — même recette : fond blanc, bordure épaisse colorée,
% ombre dure encre, badge plein. Titre optionnel [#1].
% ============================================================
\tcbset{popbase/.style={breakable,enhanced,colback=white,boxrule=2.3pt,arc=9pt,
  shadow={3pt}{-3pt}{0pt}{popink},left=14pt,right=14pt,top=11pt,bottom=11pt,before skip=11pt,after skip=15pt,
  fontupper=\popsemi\fontsize{10.6}{14.5}\selectfont\color{popink},
  fontlower=\popsemi\fontsize{10.6}{14.5}\selectfont\color{popink}}}
% Coche dessinée (le glyphe ✓ n'existe pas dans les polices du document)
\newcommand{\popcheck}[1]{\tikz[baseline=-0.5ex]\draw[#1,line width=1.6pt,line cap=round,line join=round](-0.13,0.0)--(-0.04,-0.09)--(0.13,0.1);}
\providecommand{\coche}{\popcheck{popgreen}}
\providecommand{\resultat}[1]{\par\smallskip\noindent\fcolorbox{popgreen}{popgreenL}{\parbox{\dimexpr\linewidth-2\fboxsep-2\fboxrule\relax}{\textbf{Result.} #1}}\par\smallskip}
\newcommand{\popboxhead}[3]{\poppill{#1}{#2}{white}\ifx\relax#3\relax\else\hspace{6pt}{\popxbold\fontsize{10.6}{12}\selectfont\color{popink}#3}\fi\par\vspace{7pt}}

\newtcolorbox{aretenir}[1][]{popbase,colframe=popgreen,before upper={\popboxhead{Key points}{popgreen}{#1}}}
\newtcolorbox{exemplebox}[1][]{popbase,colframe=poporange,before upper={\popboxhead{Exemple}{poporange}{#1}}}
\newtcolorbox{definition}[1][]{popbase,colframe=popblue,before upper={\popboxhead{Definition}{popblue}{#1}}}
\newtcolorbox{propriete}[1][]{popbase,colframe=poppurple,before upper={\popboxhead{Property}{poppurple}{#1}}}
\newtcolorbox{methode}[1][]{popbase,colframe=popgold,before upper={\popboxhead{Method}{popgold}{#1}}}
\newtcolorbox{attention}[1][]{popbase,colframe=poppink,before upper={\popboxhead{Warning}{poppink}{#1}}}
\newtcolorbox{experience}[1][]{popbase,colframe=popblue,colback=popblueL,before upper={\popboxhead{Experiment}{popblue}{#1}}}
% « Le savais-tu ? » : carte violette pleine (contexte, culture, Cameroun)
\newtcolorbox{savaistu}[1][]{popbase,colback=poppurple,colframe=popink,
  fontupper=\popsemi\fontsize{10.4}{14.2}\selectfont\color{white},
  before upper={\poppill{Did you know?}{white}{poppurple}\ifx\relax#1\relax\else\hspace{6pt}{\popxbold\color{white}#1}\fi\par\vspace{7pt}}}
% Exercice résolu : énoncé en haut, \tcblower puis la solution sur fond crème
\newcounter{popexo}\newcounter{popexr}
\newtcolorbox{exoresolu}[1][]{popbase,colframe=poporange,colbacklower=popcream,
  segmentation style={popink,line width=1.2pt,dashed},
  before upper={\stepcounter{popexr}\popboxhead{Worked example \thepopexr}{poporange}{#1}},
  before lower={\poppill{Solution}{popink}{popcream}\par\vspace{6pt}}}
% Exercice (non résolu en place) — la correction est dans la partie Corrigés
\newtcolorbox{exercice}[1][]{popbase,colframe=popink,
  before upper={\stepcounter{popexo}\popboxhead{Exercise \thepopexo}{popink}{#1}}}
% Corrigé
\newtcolorbox{corrige}[1][]{popbase,colframe=popgreen,colback=popgreenL,
  before upper={\popboxhead{Solution}{popgreen}{#1}}}
% Objectifs / prérequis / bilan
\newtcolorbox{objectifs}{popbase,colframe=popink,colback=poporangeL,
  before upper={\popboxhead{Chapter objectives}{popink}{}}}
\newtcolorbox{prerequis}{popbase,colframe=popink,colback=popblueL,
  before upper={\popboxhead{Prerequisites}{popblue}{}}}
\newtcolorbox{bilan}{popbase,colframe=popink,colback=white,boxrule=2.8pt,
  before upper={\popboxhead{Fiche bilan}{poporange}{L'essentiel à connaître pour l'examen}}}
% Figure encadrée avec légende
\newtcolorbox{popfigure}[1][]{enhanced,breakable=false,colback=white,colframe=popline,boxrule=1.4pt,arc=8pt,
  left=10pt,right=10pt,top=10pt,bottom=8pt,before skip=10pt,after skip=12pt,halign=center,
  lower separated=false,halign lower=center,
  fontlower=\popsemi\fontsize{9}{11.5}\selectfont\color{popmuted},
  #1}
\newcommand{\legende}[1]{\par\vspace{6pt}{\popsemi\fontsize{9}{11.5}\selectfont\color{popmuted}#1}}

% ============================================================
% Signature OnBuch+
% ============================================================
\newcommand{\poplogoinline}[1]{{\poplogo\fontsize{#1}{#1}\selectfont\color{popink}OnBuch\color{poporange}+}}
\newcommand{\signatureonbuch}{%
  \begin{tcolorbox}[enhanced,colback=popink,colframe=popink,boxrule=2.2pt,arc=10pt,
      drop shadow=poporange,left=14pt,right=14pt,top=10pt,bottom=10pt,before skip=0pt,after skip=0pt]
    \tikz[baseline=-0.7ex]\node[circle,fill=popgreen,draw=popcream,line width=1.3pt,minimum size=22pt,inner sep=0]{\popcheck{white}};%
    \hspace{9pt}\begin{minipage}[c]{0.62\linewidth}
      {\popxbold\fontsize{12}{13}\selectfont\color{popcream}Signed\ \textbullet\ {\poplogo OnBuch\color{poporange}+}}\par\vspace{2pt}
      {\raggedright\popsemi\fontsize{8}{10.5}\selectfont\color{popline}\MakeUppercase{OnBuch+ teaching team}\\\MakeUppercase{Follows the official MINESEC syllabus}\par}
    \end{minipage}\hfill
    {\poplogo\fontsize{22}{22}\selectfont\color{popcream}OnBuch\color{poporange}+}
  \end{tcolorbox}%
}

% ============================================================
% Page de garde
% #1 titre · #2 matière · #3 niveau · #4 module · #5 n° leçon · #6 durée ·
% #7 description / objectifs (texte ou itemize)
% ============================================================
\newcommand{\pagegarde}[7]{%
  \thispagestyle{empty}
  \newgeometry{a4paper,margin=1.7cm,top=1.6cm,bottom=1.4cm}%
  \begin{tikzpicture}[remember picture,overlay]
    \fill[poporange!18] ([xshift=-3.2cm,yshift=-3.6cm]current page.north east) circle (4.6cm);
    \fill[poppurple!14] ([xshift=2.4cm,yshift=7.6cm]current page.south west) circle (3.8cm);
  \end{tikzpicture}%
  {\poplogo\fontsize{30}{30}\selectfont\color{popink}OnBuch\color{poporange}+}\hfill
  \raisebox{0.6ex}{\poppill{Signed course \textbullet\ Premium}{popink}{popcream}}\par
  \vspace{1pt}\popsquiggle{poppurple}\par
  \vspace{8pt}
  \begin{tcolorbox}[enhanced,colback=poporange,colframe=popink,boxrule=2.8pt,arc=13pt,
      drop shadow=popink,left=18pt,right=18pt,top=12pt,bottom=13pt,after skip=10pt]
    \poppill{#2 \textbullet\ #3}{popink}{popcream}\hspace{5pt}\poppill{#4}{white}{popink}\par\vspace{8pt}
    {\popdisplay\fontsize{14}{14}\selectfont\color{popink}CHAPTER #5}\par\vspace{4pt}
    {\raggedright\hyphenpenalty=10000\exhyphenpenalty=10000\popdisplay\fontsize{25}{27}\selectfont\color{popcream}\MakeUppercase{#1}\par}
    \vspace{8pt}
    {\popxbold\fontsize{9.5}{9.5}\selectfont\color{popink}\MakeUppercase{Suggested time: #6}}
  \end{tcolorbox}
  \begin{tcolorbox}[enhanced,colback=white,colframe=popink,boxrule=2.2pt,arc=10pt,
      drop shadow=popink,left=16pt,right=16pt,top=10pt,bottom=10pt,after skip=8pt]
    \poppill{In this chapter}{poppurple}{white}\par\vspace{5pt}
    \popsemi\fontsize{9.3}{12.6}\selectfont\color{popink}#7
  \end{tcolorbox}
  \noindent\begin{minipage}[t]{0.487\linewidth}
    \begin{tcolorbox}[enhanced,colback=popgreen,colframe=popink,boxrule=2.2pt,arc=10pt,
        drop shadow=popink,left=11pt,right=11pt,top=10pt,bottom=10pt,height=3.1cm,valign=center]
      \tcbox[colback=white,colframe=popink,boxrule=1.3pt,arc=5pt,boxsep=0pt,nobeforeafter,
        left=3pt,right=3pt,top=3pt,bottom=3pt,valign=center]{\qrcode[height=1.9cm]{https://play.google.com/store/apps/details?id=com.luvvix.onbuch&hl=en}}%
      \hspace{8pt}\begin{minipage}[c]{0.5\linewidth}
        {\popdisplay\fontsize{11}{12.5}\selectfont\color{white}\MakeUppercase{Download OnBuch}}\par\vspace{4pt}
        {\raggedright\popsemi\fontsize{8.4}{11}\selectfont\color{white}Free \textbullet\ Google Play\\AI tutor, past papers, results\par}
      \end{minipage}
    \end{tcolorbox}
  \end{minipage}\hfill
  \begin{minipage}[t]{0.487\linewidth}
    \begin{tcolorbox}[enhanced,colback=poppurple,colframe=popink,boxrule=2.2pt,arc=10pt,
        drop shadow=popink,left=11pt,right=11pt,top=10pt,bottom=10pt,height=3.1cm,valign=center]
      \tikz[baseline=-0.7ex]\node[circle,fill=white,draw=popink,line width=1.3pt,minimum size=1.6cm,inner sep=0]{\popdisplay\fontsize{24}{24}\selectfont\color{poppurple}+};%
      \hspace{8pt}\begin{minipage}[c]{0.6\linewidth}
        {\popdisplay\fontsize{11}{12.5}\selectfont\color{white}\MakeUppercase{Go OnBuch+}}\par\vspace{4pt}
        {\raggedright\popsemi\fontsize{8.4}{11}\selectfont\color{white}All signed courses, unlimited AI tutor, PDF sheets — OnBuch+ tab in the app\par}
      \end{minipage}
    \end{tcolorbox}
  \end{minipage}
  \vspace{8pt}
  \popcontenu
  \vfill
  \signatureonbuch
  \restoregeometry
  \newpage
}

% Rangée « Ce cours contient » (page de garde)
\newcommand{\poptile}[3]{% #1 couleur #2 gros texte #3 légende
  \begin{tcolorbox}[enhanced,colback=#1,colframe=popink,boxrule=2pt,arc=9pt,shadow={2.5pt}{-2.5pt}{0pt}{popink},
      left=6pt,right=6pt,top=9pt,bottom=9pt,halign=center,valign=center,height=1.75cm,width=\linewidth,nobeforeafter]
    {\popdisplay\fontsize{12.5}{13}\selectfont\color{white}\MakeUppercase{#2}}\par\vspace{3pt}
    {\centering\popsemi\fontsize{7.8}{9.5}\selectfont\color{white}#3\par}
  \end{tcolorbox}}
\newcommand{\popcontenu}{%
  \par\noindent{\popxboldls\fontsize{7.5}{7.5}\selectfont\color{popmuted}THIS SIGNED COURSE CONTAINS}\par\vspace{7pt}
  \noindent\begin{minipage}[t]{0.235\linewidth}\poptile{popblue}{Course}{complete, illustrated, step by step}\end{minipage}\hfill
  \begin{minipage}[t]{0.235\linewidth}\poptile{poporange}{Methods}{and worked examples}\end{minipage}\hfill
  \begin{minipage}[t]{0.235\linewidth}\poptile{poppink}{Exercises}{exam style + solutions}\end{minipage}\hfill
  \begin{minipage}[t]{0.235\linewidth}\poptile{popgreen}{Summary sheet}{the essentials to remember}\end{minipage}\par}

% Sommaire
\newtcolorbox{sommaire}{enhanced,colback=white,colframe=popink,boxrule=2.2pt,arc=10pt,
  drop shadow=popink,left=18pt,right=18pt,top=13pt,bottom=13pt,before skip=6pt,after skip=20pt,
  before upper={\poppill{Contents}{poppurple}{white}\par\vspace{9pt}\popsemi\fontsize{10.6}{14.5}\selectfont\color{popink}}}

% Signature de fin de cours — poussée tout en bas de la dernière page
\newcommand{\findecours}{%
  \par\vfill\nopagebreak
  \begin{center}\popsquiggle{poporange}\end{center}\vspace{4pt}
  \signatureonbuch
}
\AtBeginDocument{\def\llbracket{\mathopen{[\mkern-3.5mu[}}\def\rrbracket{\mathclose{]\mkern-3.5mu]}}} % intervalles d'entiers (glyphe ⟦ absent de la police)
% Glyphes absents de Fira Math : redéfinis à partir de glyphes disponibles
\AtBeginDocument{%
  \def\setminus{\mathbin{\backslash}}\let\smallsetminus\setminus
  \def\vdots{\mathord{\vbox{\baselineskip3.2pt\lineskiplimit0pt\kern4pt\hbox{.}\hbox{.}\hbox{.}}}}%
  \def\ddots{\mathinner{\mkern1mu\raise6.5pt\hbox{.}\mkern2mu\raise3.6pt\hbox{.}\mkern2mu\raise0.7pt\hbox{.}\mkern1mu}}%
  \def\triangle{\mathord{\scalebox{0.9}{$\symup{\Delta}$}}}%
  \def\nabla{\mathord{\raisebox{\depth}{\scalebox{1}[-1]{$\symup{\Delta}$}}}}%
  \def\checkmark{\popcheck{popdark}}%
  \def\star{\mathbin{\ast}}\def\bigstar{\mathord{\ast}}%
  \def\diamond{\mathbin{\scalebox{0.6}{$\Diamond$}}}%
  \def\bigcup{\mathop{\mathchoice{\raisebox{-0.3em}{\scalebox{1.7}{$\cup$}}}{\scalebox{1.25}{$\cup$}}{\cup}{\cup}}\displaylimits}%
  \def\bigcap{\mathop{\mathchoice{\raisebox{-0.3em}{\scalebox{1.7}{$\cap$}}}{\scalebox{1.25}{$\cap$}}{\cap}{\cap}}\displaylimits}%
  \def\lesseqqgtr{\mathrel{\lessgtr}}\def\bigoplus{\mathop{\oplus}\displaylimits}%
}
\providecommand{\ssi}{if and only if} % abréviation fréquente
\AtBeginDocument{\providecommand{\wideparen}{\overparen}} % arc de cercle

%% FIGFIX-BEGIN v3 — correctifs globaux des figures (docs/onbuch-generation/tools/figfix)
\usepackage{adjustbox}\usepackage{etoolbox}
% toute figure TikZ/pgfplots plus large que la ligne est réduite à la largeur disponible
\BeforeBeginEnvironment{tikzpicture}{\begin{adjustbox}{max width=\linewidth}}
\AfterEndEnvironment{tikzpicture}{\end{adjustbox}}
% légendes pgfplots lisibles par-dessus les courbes
\pgfplotsset{every axis/.append style={legend style={fill=white,fill opacity=0.92,text opacity=1,draw=black!15,inner sep=3pt}}}
% étiquettes d'arêtes (syntaxe edge["..."]) avec fond blanc
\tikzset{every edge quotes/.append style={fill=white,inner sep=1.2pt}}
% bulles de cartes mentales : texte à la ligne dans la bulle, bulle un peu plus grande
\tikzset{every concept/.append style={text width=2.0cm,align=center,minimum size=2.3cm,inner sep=2pt}}
%% FIGFIX-END
\begin{document}

\pagegarde{Communication}{Physics}{Upper Sixth}{Upper Sixth — Physics}{14}{5 periods}{This chapter explores the physics of modern communication systems: how information is represented in analogue and digital form, how radio waves propagate, how aerials and tuning circuits select signals, and how modulation, digitisation and communication channels (including satellites) make broadcasting and telephony possible.
\vspace{4pt}
\begin{multicols}{2}\small
\begin{itemize}
\item By the end of this chapter I can distinguish analogue from digital methods of representing information and state the advantages and disadvantages of each.
\item By the end of this chapter I can describe the three modes of radio wave propagation (ground, sky and space waves) and match each to its frequency range and typical use.
\item By the end of this chapter I can explain the roles of transmitting and receiving aerials and relate aerial length to wavelength.
\item By the end of this chapter I can explain how an LC tuning circuit selects a station and interpret its resonance curve.
\item By the end of this chapter I can explain the principles of modulation and demodulation and compare AM and FM.
\item By the end of this chapter I can describe the operation of analogue-to-digital and digital-to-analogue converters, including sampling and quantisation.
\end{itemize}\end{multicols}}

\begin{sommaire}
\begin{multicols}{2}
\begin{enumerate}
\item Representing Information and Radio Wave Propagation
\item Aerials and the Tuning Circuit
\item Modulation, Demodulation, AM and FM
\item Digital Conversion and Communication Channels
\item Integration activity
\item Methods and worked examples
\item Exercises
\item Solutions to the exercises
\item Summary sheet
\end{enumerate}
\end{multicols}
\end{sommaire}

\begin{objectifs}
\begin{itemize}
\item By the end of this chapter I can distinguish analogue from digital methods of representing information and state the advantages and disadvantages of each.
\item By the end of this chapter I can describe the three modes of radio wave propagation (ground, sky and space waves) and match each to its frequency range and typical use.
\item By the end of this chapter I can explain the roles of transmitting and receiving aerials and relate aerial length to wavelength.
\item By the end of this chapter I can explain how an LC tuning circuit selects a station and interpret its resonance curve.
\item By the end of this chapter I can explain the principles of modulation and demodulation and compare AM and FM.
\item By the end of this chapter I can describe the operation of analogue-to-digital and digital-to-analogue converters, including sampling and quantisation.
\item By the end of this chapter I can define bandwidth and sidebands and calculate the bandwidth of an AM signal.
\item By the end of this chapter I can describe communication channels, the use of satellites and the role of base stations in mobile telephony.
\end{itemize}
\end{objectifs}

\begin{prerequis}
\begin{itemize}
\item Electromagnetic waves: nature, speed c = 3.0 × 10\^{}8 m/s, relation c = fλ (Lower Sixth)
\item Free oscillations and resonance in an LC circuit; natural frequency f₀ = 1/(2π√(LC)) (Lower Sixth)
\item Alternating currents: peak and r.m.s. values, frequency, period (Lower Sixth)
\item Wave phenomena: reflection, refraction, diffraction and superposition (Lower Sixth)
\item Binary numbers and basic logic levels (useful for digital signals)
\end{itemize}
\end{prerequis}

\newpage

\coursec{Representing Information and Radio Wave Propagation}

Every evening in Yaoundé, Douala or Bamenda, millions of Cameroonians tune in to CRTV, listen to radio stations like Magic FM, or receive calls routed through MTN and Orange base stations --- yet the voice of a presenter in Douala can be heard clearly in Wum, hundreds of kilometres away. How does invisible information travel through air, bend round the curve of the Earth, bounce off the sky, or hop to a satellite $36\,000\ \mathrm{km}$ above the equator? Why does your radio crackle during a storm while a digital TV decoder keeps a stable picture? This chapter uncovers the physics behind these everyday miracles: how information is \cle{represented}, carried by \cle{radio waves}, modulated, digitised and channelled across Cameroon and the world.

In this first section we answer two questions: \emph{how is information represented} (analogue or digital), and \emph{how does a radio wave physically travel} from a transmitting mast to your receiver (ground wave, sky wave, space wave).

\courssub{The Nature of Information in Communication}

\textbf{What is information?}\par
In physics, \cle{information} is any message that can be measured, stored or transmitted: the sound of a voice, a photograph, a text message, a video, or numerical data from a weather station on Mount Cameroon. None of these can travel far in its raw form. Sound waves die out after a few hundred metres; light from a picture cannot pass through walls.

The solution used in all communication systems is a three-stage chain:
\begin{enumerate}
\item a \cle{transducer} converts the message into an electrical signal (a microphone turns sound into a varying voltage);
\item the electrical signal is impressed on a \cle{carrier} --- usually a high-frequency electromagnetic wave --- which can travel enormous distances;
\item the carrier propagates through a \cle{transmission medium} (the atmosphere, a cable, an optical fibre, or free space to a satellite) to a receiver, which extracts the original message.
\end{enumerate}

\begin{aretenir}[The communication chain]
Information $\rightarrow$ electrical signal $\rightarrow$ carrier wave $\rightarrow$ transmission medium $\rightarrow$ receiver $\rightarrow$ original information. The carrier does the travelling; the information only ``rides'' on it.
\end{aretenir}

\courssub{The Analogue Method of Representing Information}

Speak into an old telephone microphone and watch the voltage on an oscilloscope: the voltage rises and falls \emph{continuously}, copying the shape of the sound wave. When the sound pressure doubles, the voltage doubles. The electrical signal is an \emph{analogy} of the original message --- hence the name.

\begin{definition}[Analogue signal]
An \cle{analogue signal} is a signal in which a physical quantity (voltage, current, amplitude) varies \textbf{continuously} with time, taking \textbf{any value} within a range, so that its variation mimics the variation of the original information.
\end{definition}

Everyday examples of the analogue method:
\begin{itemize}
\item the \textbf{vinyl record}: the groove wiggles continuously, copying the sound waveform;
\item the \textbf{analogue telephone}: the microphone voltage varies continuously with the voice;
\item \textbf{AM and FM radio broadcasting}: the amplitude or frequency of the carrier varies continuously with the audio signal.
\end{itemize}

\courssub{The Digital Method of Representing Information}

Now imagine instead that you measure the signal at regular instants and write down each measurement as a number, using only the digits 0 and 1. The smooth curve is replaced by a \emph{staircase} of discrete levels, and each level is coded in \cle{binary}.

\begin{definition}[Digital signal]
A \cle{digital signal} is a signal that can take only a \textbf{finite number of discrete values}, usually two (low $=0$, high $=1$). The original information is represented by a sequence of binary digits (\cle{bits}).
\end{definition}

Converting an analogue signal to digital requires two operations:
\begin{itemize}
\item \cle{sampling}: measuring the value of the signal at regular time intervals (for a CD, $44\,100$ times per second);
\item \cle{quantisation}: rounding each measured value to the nearest allowed level, then coding that level in binary. With $n$ bits per sample, $2^{n}$ different levels can be coded (for example 2 bits give $2^{2}=4$ levels: 00, 01, 10, 11).
\end{itemize}

\begin{popfigure}
\begin{center}
\begin{tikzpicture}[scale=1.0]
% Analogue signal
\begin{scope}
\draw[->,thick] (0,0) -- (6.4,0) node[below left,font=\small]{$t$};
\draw[->,thick] (0,-0.2) -- (0,2.6) node[left,font=\small]{$u$};
\draw[very thick,popblue,domain=0:6,samples=100] plot (\x,{1.2+0.9*sin(120*\x)+0.35*sin(300*\x)});
\node[font=\small\bfseries,popblue] at (3,3.0) {Analogue: continuous};
\end{scope}
% Digital staircase
\begin{scope}[xshift=8cm]
\draw[->,thick] (0,0) -- (6.4,0) node[below left,font=\small]{$t$};
\draw[->,thick] (0,-0.2) -- (0,2.6) node[left,font=\small]{$u$};
\draw[very thick,poporange] (0,1.0) -- (0.75,1.0) -- (0.75,1.8) -- (1.5,1.8) -- (1.5,2.2) -- (2.25,2.2) -- (2.25,1.4) -- (3.0,1.4) -- (3.0,0.6) -- (3.75,0.6) -- (3.75,0.9) -- (4.5,0.9) -- (4.5,1.7) -- (5.25,1.7) -- (5.25,2.1) -- (6.0,2.1);
\foreach \x in {0.75,1.5,2.25,3.0,3.75,4.5,5.25} \draw[dashed,gray] (\x,0) -- (\x,2.4);
\node[font=\small\bfseries,poporange] at (3,3.0) {Digital: sampled levels};
\node[font=\footnotesize,popdark] at (0.37,-0.35) {01};
\node[font=\footnotesize,popdark] at (1.12,-0.35) {10};
\node[font=\footnotesize,popdark] at (1.87,-0.35) {11};
\node[font=\footnotesize,popdark] at (2.62,-0.35) {11};
\node[font=\footnotesize,popdark] at (3.37,-0.35) {01};
\node[font=\footnotesize,popdark] at (4.12,-0.35) {00};
\node[font=\footnotesize,popdark] at (4.87,-0.35) {10};
\node[font=\footnotesize,popdark] at (5.62,-0.35) {11};
\end{scope}
\end{tikzpicture}
\end{center}
\legende{An analogue waveform varies continuously; its digitised version is a staircase of quantised levels, each coded in binary.}
\end{popfigure}

\courssub{Analogue versus Digital: Advantages and Disadvantages}

Why did Cameroon switch television broadcasting from analogue to digital? The answer lies in how each method fights \cle{noise} --- the unwanted random disturbances (storm static, engine ignition, other stations) that add to any signal during transmission.

An analogue signal that picks up noise is damaged \emph{permanently}: the receiver cannot tell which part of the wiggle is the music and which part is the crackle. A digital signal only needs the receiver to decide ``0 or 1?''. Even if noise distorts the pulse, as long as it is still recognisably a 1, the original value is recovered perfectly. Moreover, a digital \cle{regenerator} can receive a weakened, noisy pulse train, decide each bit, and retransmit a \emph{brand-new, clean} copy --- something impossible for analogue amplifiers, which amplify the noise together with the signal.

\begin{center}
\begin{tabularx}{\linewidth}{|l|X|X|}
\hline
\rowcolor{popblueL} \textbf{Criterion} & \textbf{Analogue} & \textbf{Digital} \\
\hline
Noise immunity & Poor: noise adds permanently to the signal & Excellent: bits can be regenerated perfectly \\
\hline
Copying / storage & Each copy degrades (tape hiss) & Perfect copies, easy storage on memory chips \\
\hline
Processing & Difficult & Easy: computers, error correction, encryption \\
\hline
Bandwidth required & Small & Large (many bits per second) \\
\hline
Circuit cost / complexity & Simple, cheap circuits & More complex electronics, needs converters \\
\hline
\end{tabularx}
\end{center}

\begin{aretenir}[Key comparison]
Digital wins on \cle{noise immunity}, \cle{regeneration}, \cle{storage} and \cle{processing}; analogue wins on \cle{simplicity} and \cle{low bandwidth}. This is why mobile phones, digital TV and the internet are digital, while a simple microphone lead remains analogue.
\end{aretenir}

\courssub{Radio Waves and the Electromagnetic Spectrum}

Radio waves are electromagnetic waves --- oscillating electric and magnetic fields travelling through space at the speed of light $c = 3.00\times 10^{8}\ \mathrm{m\,s^{-1}}$. They occupy the low-frequency end of the electromagnetic spectrum, from a few kilohertz up to hundreds of gigahertz. Frequency and wavelength are linked by
\[ c = f\lambda \qquad\Longleftrightarrow\qquad \lambda = \frac{c}{f}. \]

Different frequency bands behave very differently in the atmosphere, so each band is reserved for particular services:

\begin{popfigure}
\begin{center}
\begin{tikzpicture}[scale=1.0]
\draw[->,thick] (0,0) -- (13.5,0) node[below left,font=\small]{frequency (log scale)};
\foreach \x/\lab in {0.5/{30 kHz},2.6/{300 kHz},4.7/{3 MHz},6.8/{30 MHz},8.9/{300 MHz},11.0/{3 GHz}}{
  \draw[thick] (\x,0) -- (\x,-0.15);
  \node[below,font=\footnotesize] at (\x,-0.15) {\lab};
}
% Bands
\fill[popblue!60] (0.5,0.25) rectangle (2.6,1.05); \node[font=\footnotesize,align=center] at (1.55,0.65) {LW};
\fill[popgreen!60] (2.6,0.25) rectangle (4.7,1.05); \node[font=\footnotesize,align=center] at (3.65,0.65) {MW (AM radio)};
\fill[poporange!70] (4.7,0.25) rectangle (6.8,1.05); \node[font=\footnotesize,align=center] at (5.75,0.65) {SW};
\fill[poppink!70] (6.8,0.25) rectangle (8.9,1.05); \node[font=\footnotesize,align=center] at (7.85,0.65) {FM / TV};
\fill[poppurple!50] (8.9,0.25) rectangle (11.0,1.05); \node[font=\footnotesize,align=center,text=white] at (9.95,0.65) {Mobile / TV};
\fill[popgold!70] (11.0,0.25) rectangle (13.2,1.05); \node[font=\footnotesize,align=center] at (12.1,0.65) {Satellite};
% Propagation labels
\node[font=\footnotesize,popblue,align=center] at (1.55,1.45) {ground wave};
\node[font=\footnotesize,popgreen!50!black,align=center] at (3.65,1.45) {ground + sky};
\node[font=\footnotesize,poporange!80!black,align=center] at (5.75,1.45) {sky wave};
\node[font=\footnotesize,poppink!80!black,align=center] at (7.85,1.45) {space wave};
\node[font=\footnotesize,poppurple,align=center] at (9.95,1.45) {space wave};
\node[font=\footnotesize,popgold!80!black,align=center] at (12.1,1.45) {space wave};
\end{tikzpicture}
\end{center}
\legende{The radio spectrum (logarithmic frequency scale) with typical uses and the dominant propagation mode of each band.}
\end{popfigure}

\begin{exemplebox}[Wavelength of an FM station]
A Yaoundé FM station broadcasts at $f = \SI{100}{MHz}$. Its wavelength is
\[ \lambda = \frac{c}{f} = \frac{3.00\times 10^{8}}{1.00\times 10^{8}} = \SI{3.0}{m}. \]
Compare this with a MW station at $f = \SI{1.0}{MHz}$: $\lambda = \SI{300}{m}$. Long waves bend around hills and buildings far better than short ones --- a first clue to why the bands propagate differently.
\end{exemplebox}

\courssub{Surface (Ground) Wave Propagation}

A radio wave launched close to the Earth's surface does not travel only in a straight line: at low frequencies it \cle{diffracts} (bends) around the curvature of the Earth and follows the ground, rather as water waves bend around a headland. This is the \cle{surface wave} or \cle{ground wave}.

\begin{definition}[Surface (ground) wave]
The \cle{surface wave} is a radio wave that travels along the Earth's surface, following its curvature by diffraction. It is effective for frequencies \textbf{below about $\SI{3}{MHz}$} (long and medium waves) and can reach hundreds of kilometres over land and over a thousand kilometres over the sea, which is a good electrical conductor.
\end{definition}

Diffraction is strong only when the wavelength is comparable to or larger than the obstacles and the scale of the curvature --- hence the restriction to low frequencies. The ground wave is steadily absorbed by the soil, so its range is limited and it weakens over rough or poorly conducting terrain such as the rocky highlands of the North-West Region; over sea water, a much better conductor, the losses are smaller and the range is greater.

\courssub{Sky Wave Propagation: Bending Back from the Ionosphere}

High in the atmosphere, between roughly $60\ \mathrm{km}$ and $400\ \mathrm{km}$ above the ground, ultraviolet radiation from the Sun ionises the gases, creating a region of free electrons and ions called the \cle{ionosphere}. Free electrons can be set oscillating by a passing radio wave and re-radiate it; the net effect is that the wave is progressively \emph{refracted} --- bent --- back towards the ground.

\begin{propriete}[Reflection by the ionosphere]
The ionosphere bends back (``reflects'') radio waves in the approximate range $\SI{3}{MHz}$ to $\SI{30}{MHz}$ (medium and short waves). Below this range the waves are largely absorbed; above it (FM, TV, mobile frequencies) the ionosphere is essentially \textbf{transparent} and the waves pass through into space. (The explanation of this frequency dependence is not examinable, but the property itself is.)
\end{propriete}

A sky wave returns to Earth at a considerable distance from the transmitter. The region between the limit of the ground wave and the point where the sky wave first lands is called the \cle{skip distance} (or skip zone): a listener there receives nothing, while a listener much further away receives the station clearly. At night the lower, absorbing layers of the ionosphere weaken, so medium-wave stations can be heard over the whole country and even from abroad --- which is why foreign stations appear on your MW dial after sunset.

\begin{attention}[The ionosphere is not a mirror]
A very common mistake is to say the sky wave ``bounces off the ionosphere like light off a mirror''. In fact the ionosphere is a \emph{gradually varying} medium: the wave is \textbf{refracted} --- bent little by little as it penetrates deeper into the electron layer --- until its path curves back down towards the Earth. There is no hard reflecting surface. In the exam, write ``refraction (gradual bending) in the ionosphere'', not ``reflection off a solid layer''.
\end{attention}

\courssub{Space Wave Propagation: Line of Sight}

Above about $\SI{30}{MHz}$ the ionosphere lets the wave through and the wavelength is too short to diffract usefully around the Earth. The wave therefore travels essentially in a \textbf{straight line} from the transmitting aerial to the receiving aerial: this is the \cle{space wave} (or direct wave), used by FM radio ($88$--$\SI{108}{MHz}$), television and mobile phone networks.

Because the Earth is curved, the straight-line path is blocked by the horizon. The maximum range is roughly the line-of-sight distance between the tops of the two aerials, which is why transmitters are placed on high masts and hills (for example on Mount Febe or Mount Cameroon) and why networks of \cle{relay masts} and base stations are needed to cover a country.

For a transmitting aerial of height $h$ on a spherical Earth of radius $R$, the tangent from the top of the mast touches the Earth at the horizon. Pythagoras' theorem gives $(R+h)^{2} = R^{2} + d^{2}$, and since $h \ll R$,
\[ d^{2} = 2Rh + h^{2} \approx 2Rh \qquad\Longrightarrow\qquad d \approx \sqrt{2Rh}. \]
With $R = 6.4\times 10^{6}\ \mathrm{m}$ and a mast of height $h = \SI{100}{m}$, the horizon is at $d \approx \sqrt{2\times 6.4\times 10^{6}\times 100} \approx \SI{36}{km}$ --- the typical order of magnitude of an FM station's useful range.

\begin{popfigure}
\begin{center}
\begin{tikzpicture}[scale=1.0]
% Earth
\fill[popblue!15] (-6,-4.2) arc[start angle=180,end angle=360,radius=6] -- cycle;
\draw[thick,popblue] (-6,0) arc[start angle=180,end angle=360,radius=6];
\node[font=\small,popblue] at (0,-3.6) {Earth};
% Ionosphere
\draw[thick,poppurple,dashed] (-6.6,1.9) arc[start angle=165,end angle=15,radius=6.6];
\node[font=\small,poppurple] at (4.6,2.6) {Ionosphere};
% Transmitter mast
\draw[thick,popdark] (-4.2,1.55) -- (-4.2,3.1);
\draw[thick,popdark] (-4.35,1.55) -- (-4.05,1.55);
\node[font=\footnotesize,align=center] at (-4.2,3.35) {Tx mast};
% Ground wave
\draw[very thick,popgreen!60!black] (-4.2,1.62) arc[start angle=160,end angle=105,radius=6.05];
\node[font=\footnotesize,popgreen!50!black,align=center] at (-2.6,1.15) {ground wave};
% Sky wave
\draw[very thick,poporange] (-4.2,3.1) -- (-1.2,5.6);
\draw[very thick,poporange,->] (-1.2,5.6) -- (2.6,1.35);
\fill[poppurple] (-1.2,5.6) circle (0.07);
\node[font=\footnotesize,poporange!80!black,align=center] at (-2.4,4.9) {sky wave};
\node[font=\footnotesize,poporange!80!black,align=center] at (2.6,0.9) {skip zone then reception};
% Space wave
\draw[very thick,poppink,->] (-4.2,3.1) -- (-0.9,3.9);
\node[font=\footnotesize,poppink!80!black,align=center] at (-3.4,4.15) {space wave};
% Receiver
\draw[thick,popdark] (2.6,1.35) -- (2.6,2.15);
\node[font=\footnotesize,align=center] at (2.6,2.4) {Rx};
\end{tikzpicture}
\end{center}
\legende{The three propagation modes: the ground wave follows the Earth's curvature, the sky wave is refracted back by the ionosphere beyond the skip zone, and the space wave travels in a straight line limited by the horizon.}
\end{popfigure}

\begin{methode}[Choosing the propagation mode]
Given a frequency $f$ and a target distance $d$:
\begin{enumerate}
\item \textbf{Compare $f$ with the key limits} $\SI{3}{MHz}$ and $\SI{30}{MHz}$.
\item If $f < \SI{3}{MHz}$: use the \textbf{ground wave}; reliable over a few hundred kilometres, day and night.
\item If $\SI{3}{MHz} < f < \SI{30}{MHz}$: use the \textbf{sky wave} for long distances (hundreds to thousands of kilometres, especially at night); remember the skip zone --- nearby receivers use the ground wave instead.
\item If $f > \SI{30}{MHz}$: only the \textbf{space wave} works (line of sight). Check that $d$ is less than the horizon distance; otherwise add relay masts or use a satellite link.
\end{enumerate}
\end{methode}

\begin{exoresolu}[Why FM fades but MW travels far]
\textbf{Statement.} A driver leaves Douala for Bafoussam listening to an FM station at $\SI{98}{MHz}$. About $80\ \mathrm{km}$ from the city the signal fades away, yet a MW station at $\SI{1.2}{MHz}$ from the same city is still received clearly, and at night it can be heard in Garoua. Explain.
\tcblower
\textbf{Solution.}
The FM frequency, $f = \SI{98}{MHz} > \SI{30}{MHz}$, propagates only as a \textbf{space wave}: the ionosphere is transparent at this frequency and the wavelength ($\lambda = c/f \approx \SI{3}{m}$) is far too short to diffract around the Earth. Once the car passes beyond the radio horizon of the transmitting mast --- typically a few tens of kilometres --- the straight-line path is blocked by the Earth's curvature and by hills, so reception is lost.

The MW frequency, $f = \SI{1.2}{MHz} < \SI{3}{MHz}$, travels as a \textbf{ground wave} that diffracts along the Earth's surface, giving coverage over hundreds of kilometres. In addition, at night the ionosphere refracts medium waves back to Earth (the absorbing lower layer disappears after sunset), so a \textbf{sky wave} can reach Garoua, more than $500\ \mathrm{km}$ away. Hence MW covers the whole country at night while FM cannot be received far beyond the horizon.
\end{exoresolu}

\begin{attention}[Exam tip --- GCE Advanced Level]
A favourite GCE question asks: \emph{``Explain why FM radio ($\approx \SI{100}{MHz}$) cannot be received far beyond the horizon while MW radio can cover the whole country at night.''} Structure your answer in three points: (1) state the propagation mode available at each frequency; (2) for FM, say the ionosphere is transparent above $\SI{30}{MHz}$ and diffraction is negligible, so only line-of-sight space waves remain, limited by the horizon; (3) for MW, say the ground wave follows the Earth's curvature \emph{and} the night-time ionosphere returns a sky wave to distant receivers. Quote $c = f\lambda$ if you compare wavelengths.
\end{attention}

\begin{savaistu}[Did you know?]
Short-wave sky waves can hop between the ionosphere and the ground several times, circling the globe. During the twentieth century, Cameroonians listened to international broadcasters such as the BBC on short wave thanks to these multi-hop paths. Today, the same physics limits your mobile phone ($900$ / $\SI{1800}{MHz}$) to a few kilometres from the nearest base station --- which is exactly why MTN and Orange must plant masts in every town and along the highways.
\end{savaistu}

\begin{exercice}[Propagation modes in Cameroon]
A station wants its programmes received (a) throughout the Littoral Region only, (b) throughout the whole country at night, (c) by a satellite uplink. For each case, choose a suitable frequency band and propagation mode, and justify your choice in two or three sentences.
\end{exercice}

\begin{corrige}[Exercise 1]
(a) An FM transmitter around $\SI{100}{MHz}$ on a high mast: the space wave gives clean, interference-free coverage up to the horizon, ideal for regional coverage. (b) A medium-wave transmitter around $\SI{1}{MHz}$: the ground wave covers the region by day, and at night the ionosphere refracts a sky wave back to Earth over the whole country. (c) A microwave frequency of several GHz: far above $\SI{30}{MHz}$, the wave passes straight through the ionosphere as a space wave and can reach a geostationary satellite.
\end{corrige}

\coursec{Aerials and the Tuning Circuit}

In the previous section we saw how information can be represented by analogue or digital signals, and how radio waves travel from place to place as ground waves, sky waves or space waves. But a radio wave does not appear from nowhere, and it does not deliver its message by itself. At the transmitting end, something must launch the electromagnetic wave into space; at the receiving end, something must capture it and something else must pick out \emph{one} station among the hundreds that fill the air simultaneously over Yaoundé or Bamenda. These three tasks are performed by the \cle{transmitting aerial}, the \cle{receiving aerial} and the \cle{tuning circuit}. They are the subject of this section.

\courssub{The Transmitting Aerial}

\textbf{From current to wave.}\par
Inside a transmitter, the information-bearing signal exists as a rapidly alternating current in a wire. A current in a wire, however, keeps its fields close to the wire. To send energy across hundreds of kilometres, the current must be made to \emph{radiate}: part of its electromagnetic field must detach itself and travel away as a self-propagating wave. This is precisely the job of the transmitting aerial.

\begin{definition}[Aerial (antenna)]
An \cle{aerial} (or \cle{antenna}) is a conductor, or system of conductors, designed to convert the energy of a high-frequency alternating current into electromagnetic waves (transmitting aerial), or to convert the energy of an arriving electromagnetic wave back into an alternating e.m.f.\ (receiving aerial).
\end{definition}

The physics is the following. A steady current produces only a static magnetic field, which does not radiate. An \emph{accelerating} charge, however, produces changing electric and magnetic fields that regenerate each other and propagate away at speed $c = 3.00\times 10^{8}\ \mathrm{m\,s^{-1}}$. In an aerial driven by a high-frequency alternating current, the electrons are continually accelerated back and forth, so radiation occurs continuously. The higher the frequency, the more violent the acceleration and the more efficient the radiation.

\begin{propriete}[Radiation efficiency and aerial size]
An aerial radiates efficiently only when its length is a \emph{significant fraction of the wavelength} $\lambda$ of the wave it is meant to launch. A wire much shorter than $\lambda$ behaves almost like a lumped circuit element and radiates very poorly. (Proof not examinable; the result must be known and used.)
\end{propriete}

This explains a fact of everyday life. An audio signal at $f = 1\ \mathrm{kHz}$ has wavelength
\[
\lambda = \frac{c}{f} = \frac{3.00\times 10^{8}}{1.0\times 10^{3}} = 3.0\times 10^{5}\ \mathrm{m} = 300\ \mathrm{km},
\]
so no practical aerial could radiate it directly. This is why we do not transmit the audio signal itself, but a high-frequency \emph{carrier} (for example $1\ \mathrm{MHz}$ in the medium-wave band, $\lambda = 300\ \mathrm{m}$) that is \emph{modulated} by the audio signal — the subject of the next section.

\courssub{Aerial Dimensions and Wavelength}

\textbf{The half-wave dipole.}\par
The simplest efficient aerial is the \cle{half-wave dipole}: a straight conductor, split at its centre where the transmitter is connected, whose total length is approximately half a wavelength:
\[
\ell \approx \frac{\lambda}{2} = \frac{c}{2f}.
\]
At this length the aerial resonates: the current reflected from the open ends returns in phase, and a \emph{standing wave of current} is established, with maximum current at the centre (where the feeder is connected) and zero current at the open ends. Resonance makes the dipole a strong radiator and gives it a convenient, purely resistive impedance at the feed point.

\begin{popfigure}
\begin{center}
\begin{tikzpicture}[scale=1.0]
  % dipole rods
  \draw[line width=2.2pt, popblue] (-3,0) -- (-0.25,0);
  \draw[line width=2.2pt, popblue] (0.25,0) -- (3,0);
  % feed
  \draw[popdark] (-0.25,0) -- (-0.25,-0.9);
  \draw[popdark] (0.25,0) -- (0.25,-0.9);
  \draw[popdark] (-0.25,-0.9) -- (0.25,-0.9) node[midway, below, popdark]{\small to transmitter};
  % standing current distribution
  \draw[domain=-3:3, samples=80, smooth, poporange, line width=1.2pt] plot (\x, {1.1*cos(\x*30)});
  \node[poporange, align=center] at (0,1.55) {\small current $I$};
  \node[popmuted] at (-3,-0.35) {\small $I=0$};
  \node[popmuted] at (3,-0.35) {\small $I=0$};
  % length annotation
  \draw[<->, popdark] (-3,-1.5) -- (3,-1.5) node[midway, below]{\small total length $\ell \approx \lambda/2$};
\end{tikzpicture}
\end{center}
\legende{Half-wave dipole. The standing-wave current distribution (orange) is maximum at the centre feed point and zero at the open ends.}
\end{popfigure}

\textbf{The quarter-wave monopole.}\par
Where one half of the dipole is inconvenient, a \cle{quarter-wave monopole} of length $\ell \approx \lambda/4$ is mounted vertically above a conducting \emph{ground plane} (the Earth, or the metal body of a car). The ground plane acts as a mirror and supplies the ``missing half'' electrically. The telescopic rod of a portable radio and the whip aerial on a taxi in Douala are quarter-wave monopoles.

\begin{attention}[Common mistake: $\lambda$ instead of $\lambda/2$]
When a question asks for the length of a \emph{half-wave dipole} at a given frequency, the answer is $\ell = \lambda/2 = c/(2f)$, \textbf{not} $\lambda = c/f$. Similarly a monopole is $\lambda/4$. Read the question twice and state which aerial you are using before computing.
\end{attention}

\begin{exoresolu}[Length of a dipole for an FM station]
A student in Buea wants to build a simple half-wave dipole to receive an FM station broadcasting at $f = 98.5\ \mathrm{MHz}$. Calculate the required total length of the dipole, and the length of each arm.
\tcblower
The wavelength is
\[
\lambda = \frac{c}{f} = \frac{3.00\times 10^{8}}{98.5\times 10^{6}} = 3.05\ \mathrm{m}.
\]
For a half-wave dipole,
\[
\ell = \frac{\lambda}{2} = \frac{3.05}{2} = 1.52\ \mathrm{m}.
\]
Each of the two arms therefore measures $\ell/2 \approx 0.76\ \mathrm{m}$. (In practice the length is shortened by about $5\%$ to allow for end effects, giving $\approx 1.45\ \mathrm{m}$, but the GCE expects $\lambda/2$.)
\end{exoresolu}

\courssub{The Receiving Aerial}

At the receiver, the process is reversed. The passing electromagnetic wave carries an oscillating electric field $\vec{E}$. This field exerts a force on the free electrons of the receiving conductor and drives them back and forth, inducing a small alternating \cle{e.m.f.} at the same frequency as the wave.

\begin{propriete}[Receiving aerial]
A receiving aerial placed in an electromagnetic wave has an alternating e.m.f.\ induced in it, of the same frequency as the wave. The induced e.m.f.\ is greatest when the aerial is parallel to the electric field of the wave and has a resonant length (e.g.\ $\lambda/2$). The signal is very weak — typically microvolts to millivolts.
\end{propriete}

Here lies the central difficulty of radio reception: the air is filled with waves from \emph{many} stations at once, plus television, mobile-phone and satellite signals. All of them induce e.m.f.s in the same aerial simultaneously. If the aerial were connected directly to an amplifier and loudspeaker, we would hear an unintelligible mixture of every station together. We need a device that \emph{selects} one carrier frequency and rejects the others. That device is the tuning circuit.

\courssub{The Tuning Circuit and Resonance}

\begin{definition}[Tuning circuit]
A \cle{tuning circuit} is a parallel combination of an inductor $L$ and a capacitor $C$ connected to the receiving aerial. It selects one carrier frequency from all the signals induced in the aerial by means of electrical \cle{resonance}.
\end{definition}

\begin{definition}[Resonance of the tuning circuit]
\cle{Resonance} occurs when the frequency $f$ of the incoming signal equals the natural frequency of the $LC$ circuit,
\[
f_0 = \frac{1}{2\pi\sqrt{LC}}.
\]
At resonance the circuit responds to that frequency far more strongly than to any other.
\end{definition}

\textbf{Why the parallel $LC$ selects a frequency.}\par
The inductor and capacitor exchange energy back and forth at the natural frequency $f_0$, just as a swing oscillates at its own natural frequency. Pushing a swing at exactly its natural rhythm builds up a large amplitude; pushing it at any other rhythm achieves little. In the same way, among all the e.m.f.s induced in the aerial, only the one whose frequency matches $f_0$ drives the $LC$ circuit in step and builds up a large voltage across it.

More quantitatively, the inductor presents reactance $X_L = 2\pi f L$ and the capacitor $X_C = 1/(2\pi f C)$. In parallel, at low frequency the inductor short-circuits the output (small $X_L$); at high frequency the capacitor short-circuits it (small $X_C$). Only near $f_0$, where $X_L = X_C$ and the two branch currents cancel, does the combination oppose the signal strongly.

\begin{propriete}[Impedance at resonance]
At resonance the impedance of a (ideal) parallel $LC$ circuit is \emph{maximum} — infinite for an ideal circuit, very large in practice. The weak current delivered by the aerial therefore produces its \emph{maximum voltage} across the circuit precisely at $f_0$. This voltage is passed to the amplifier and demodulator. (The formula $f_0 = 1/(2\pi\sqrt{LC})$ must be known and used; its derivation is not examinable.)
\end{propriete}

\begin{popfigure}
\begin{center}
\begin{circuitikz}[european, scale=1.0]
  % aerial
  \draw[line width=2pt, popblue] (0,3.2) -- (0,4.4);
  \draw[popblue] (-0.5,4.4) -- (0.5,4.4);
  \draw[popblue] (-0.35,4.6) -- (0.35,4.6);
  \draw[popblue] (-0.2,4.8) -- (0.2,4.8);
  \node[popblue] at (0.9,4.5) {\small aerial};
  \draw (0,3.2) to[L=$L$] (0,0.8);
  \draw (0,3.2) -- (2.6,3.2);
  \draw (0,0.8) -- (2.6,0.8);
  \draw (2.6,3.2) to[vC=$C$] (2.6,0.8);
  \draw (2.6,3.2) -- (4.2,3.2);
  \draw (2.6,0.8) -- (4.2,0.8);
  \node[popdark, align=center] at (5.6,2.0) {\small to amplifier\\ \small and detector};
  \draw[->, popdark] (4.2,3.2) -- (4.9,2.6);
  \draw[->, popdark] (4.2,0.8) -- (4.9,1.4);
\end{circuitikz}
\end{center}
\legende{Receiving aerial connected to a parallel $LC$ tuning circuit with variable capacitor $C$. The voltage across the tank circuit is maximum at the resonant frequency.}
\end{popfigure}

\textbf{Tuning to different stations.}\par
Since $f_0 = 1/(2\pi\sqrt{LC})$, changing $C$ changes the selected frequency. In a classical radio the ``tuning knob'' turns a \emph{variable capacitor}: decreasing $C$ increases $f_0$ (tunes to a higher-frequency station), increasing $C$ decreases $f_0$.

\begin{methode}[Finding the capacitance for a given station]
To tune to a station of carrier frequency $f_0$ with an inductor $L$:
\begin{enumerate}
  \item Write the resonance condition $f_0 = \dfrac{1}{2\pi\sqrt{LC}}$.
  \item Square both sides: $f_0^{2} = \dfrac{1}{4\pi^{2}LC}$.
  \item Solve for the capacitance:
  \[
  C = \frac{1}{4\pi^{2} f_0^{2} L}.
  \]
  \item Substitute $f_0$ in hertz and $L$ in henries; $C$ comes out in farads (usually picofarads).
\end{enumerate}
\end{methode}

\begin{exoresolu}[Tuning a medium-wave receiver]
A tuning circuit uses an inductor of $L = 200\ \mu\mathrm{H}$. Calculate the capacitance needed to receive a station broadcasting at $f_0 = 1.00\ \mathrm{MHz}$.
\tcblower
Using the method above:
\[
C = \frac{1}{4\pi^{2} f_0^{2} L}
   = \frac{1}{4\pi^{2}\,(1.00\times 10^{6})^{2}\,(200\times 10^{-6})}.
\]
Computing the denominator: $4\pi^{2} \approx 39.5$; $(1.00\times10^{6})^{2} = 1.00\times10^{12}$; so
\[
C = \frac{1}{39.5 \times 1.00\times10^{12}\times 2.00\times10^{-4}}
  = \frac{1}{7.90\times10^{9}}
  \approx 1.27\times10^{-10}\ \mathrm{F} = 127\ \mathrm{pF}.
\]
A variable capacitor covering roughly $127\ \mathrm{pF}$ at this setting will tune in the station.
\end{exoresolu}

\begin{attention}[GCE exam tip: range of $C$ for a whole band]
Numerical questions frequently give a fixed $L$ and ask for the \emph{range} of $C$ needed to cover a whole band, e.g.\ the MW band $530\ \mathrm{kHz}$ to $1600\ \mathrm{kHz}$. Because $C \propto 1/f_0^{2}$, the \textbf{lowest} frequency needs the \textbf{largest} capacitance and the highest frequency the smallest. Compute $C_{\max}$ at $f_{\min}$ and $C_{\min}$ at $f_{\max}$ — never the other way round.
\end{attention}

\begin{exercice}[Covering the FM band]
A tuning circuit has $L = 0.50\ \mu\mathrm{H}$. Calculate the range of capacitance required to tune across the FM band from $88\ \mathrm{MHz}$ to $108\ \mathrm{MHz}$.
\end{exercice}

\begin{corrige}[Exercise 1]
$C_{\max}$ corresponds to $f_{\min} = 88\ \mathrm{MHz}$:
\[
C_{\max} = \frac{1}{4\pi^{2}(88\times10^{6})^{2}(0.50\times10^{-6})}
= \frac{1}{39.5 \times 7.74\times10^{15}\times 5.0\times10^{-7}}
\approx 6.5\ \mathrm{pF}.
\]
$C_{\min}$ corresponds to $f_{\max} = 108\ \mathrm{MHz}$:
\[
C_{\min} = \frac{1}{4\pi^{2}(108\times10^{6})^{2}(0.50\times10^{-6})}
\approx 4.3\ \mathrm{pF}.
\]
The variable capacitor must cover approximately $4.3\ \mathrm{pF}$ to $6.5\ \mathrm{pF}$.
\end{corrige}

\courssub{The Resonance Curve, Selectivity and Quality Factor}

If we plot the voltage developed across the tuning circuit (its \emph{response}) against the frequency of the incoming signal, we obtain the \cle{resonance curve}: a sharp peak centred on $f_0$.

\begin{popfigure}
\begin{center}
\begin{tikzpicture}
\begin{axis}[
  width=11cm, height=6.2cm,
  xlabel={frequency $f$ (MHz)}, ylabel={response $V$},
  xmin=0.6, xmax=1.4, ymin=0, ymax=1.15,
  xtick={0.6,0.8,1.0,1.2,1.4},
  xticklabels={0.6,0.8,1.0,1.2,1.4},
  ytick=\empty,
  axis lines=left, thick,
]
\addplot[domain=0.6:1.4, samples=200, very thick, popblue]
  {1/(1+40*(x-1)^2)};
% f0 marker
\draw[dashed, popdark] (axis cs:1.0,0) -- (axis cs:1.0,1.0);
\node[popdark, below] at (axis cs:1.0,0.02) {$f_0$};
% half-power level
\draw[dashed, poporange] (axis cs:0.6,0.5) -- (axis cs:1.4,0.5);
\node[poporange, right] at (axis cs:1.4,0.5) {\small $V_{\max}/\sqrt{2}$};
% bandwidth arrows (solve 1/(1+40(x-1)^2)=0.5 -> x=1±0.158)
\draw[<->, popgreen, thick] (axis cs:0.842,0.56) -- (axis cs:1.158,0.56);
\node[popgreen, above] at (axis cs:1.0,0.58) {\small $\Delta f$};
\end{axis}
\end{tikzpicture}
\end{center}
\legende{Resonance curve of a tuning circuit. The bandwidth $\Delta f$ is measured between the frequencies where the response falls to $V_{\max}/\sqrt{2}$ (half-power points).}
\end{popfigure}

\begin{definition}[Bandwidth and quality factor]
The \cle{bandwidth} $\Delta f$ of the tuning circuit is the range of frequencies over which the response is at least $V_{\max}/\sqrt{2}$ of its maximum (the half-power points). The \cle{quality factor} is
\[
Q = \frac{f_0}{\Delta f}.
\]
A high $Q$ means a narrow, sharp peak; a low $Q$ means a broad, flat peak.
\end{definition}

\textbf{The selectivity--fidelity trade-off.}\par
\begin{itemize}
  \item \textbf{High $Q$ (sharp resonance):} excellent \emph{selectivity} — the circuit easily separates two stations with close carrier frequencies. But if the peak is \emph{too} narrow, it clips the sidebands of the wanted station itself (we shall see in the next section that a modulated carrier occupies a band of frequencies, not a single one), and the sound loses its high notes: poor \emph{fidelity}.
  \item \textbf{Low $Q$ (broad resonance):} the whole sideband structure passes, giving good fidelity, but neighbouring stations may leak through: poor selectivity.
\end{itemize}
A practical receiver is designed with $Q$ large enough to reject adjacent stations yet small enough to pass the full bandwidth of the wanted signal.

\begin{exoresolu}[Bandwidth of a tuning circuit]
A tuning circuit resonates at $f_0 = 1.00\ \mathrm{MHz}$ with a quality factor $Q = 80$. Find its bandwidth and comment on its suitability for an AM station whose signal occupies $\pm 5\ \mathrm{kHz}$ around the carrier.
\tcblower
\[
\Delta f = \frac{f_0}{Q} = \frac{1.00\times10^{6}}{80} = 1.25\times10^{4}\ \mathrm{Hz} = 12.5\ \mathrm{kHz}.
\]
The AM signal needs a bandwidth of $2\times 5 = 10\ \mathrm{kHz}$, which is less than $12.5\ \mathrm{kHz}$. The circuit therefore passes the whole signal (acceptable fidelity) while still rejecting stations more than about $6\ \mathrm{kHz}$ away from the carrier. The design is adequate.
\end{exoresolu}

\courssub{The Ferrite Rod Aerial (GCE favourite)}

\begin{savaistu}[Why portable MW radios use a ferrite rod]
Open a small portable medium-wave radio and you will not find a long wire: you will find a coil wound on a black rod of \emph{ferrite}, a magnetic ceramic material. At $1\ \mathrm{MHz}$ the wavelength is $300\ \mathrm{m}$, so a $\lambda/2$ dipole would be $150\ \mathrm{m}$ long — absurd for a pocket radio! Instead, the coil picks up the \emph{magnetic} component of the wave. The ferrite rod, with its very high magnetic permeability, concentrates the magnetic flux of the passing wave through the coil, multiplying the induced e.m.f.\ by a large factor. The coil itself also serves as the inductor $L$ of the tuning circuit, with a variable capacitor connected across it — aerial and tuner in one compact component. Because the rod responds to the magnetic field direction, reception is best when the rod is \emph{perpendicular} to the direction of the transmitter, which is why rotating a portable radio changes the loudness.
\end{savaistu}

\begin{aretenir}[Key points of this section]
\begin{itemize}
  \item A transmitting aerial converts a high-frequency alternating current into an electromagnetic wave; it radiates efficiently only if its length is comparable to $\lambda$.
  \item Half-wave dipole: $\ell = \lambda/2$; quarter-wave monopole: $\ell = \lambda/4$; with $\lambda = c/f$.
  \item A receiving aerial has a small alternating e.m.f.\ induced in it by every passing wave; the tuning circuit selects one.
  \item A parallel $LC$ circuit resonates at $f_0 = \dfrac{1}{2\pi\sqrt{LC}}$; at resonance its impedance is maximum, so the wanted station produces the largest voltage.
  \item Tuning is done with a variable capacitor: $C = \dfrac{1}{4\pi^{2}f_0^{2}L}$; largest $C$ for the lowest frequency.
  \item The resonance curve peaks at $f_0$; bandwidth $\Delta f = f_0/Q$; high $Q$ gives selectivity, low $Q$ gives fidelity — a compromise is needed.
  \item Portable MW radios use a ferrite rod aerial to concentrate the magnetic flux into a small coil.
\end{itemize}
\end{aretenir}

With one station now cleanly selected by the tuning circuit, the next task is to extract the actual information — music or speech — from the carrier wave. That is the work of \emph{demodulation}, which we study in the next section together with the two great modulation schemes, AM and FM.

\coursec{Modulation, Demodulation, AM and FM}

In the previous section we saw how transmitting and receiving aerials launch and capture radio waves, and how the tuning circuit of a receiver selects one station by resonance. One question remains: \emph{how does the sound of a voice or of music actually travel on a radio wave?} A pure radio wave of constant amplitude and constant frequency carries no information at all. The answer is \cle{modulation}: we force a high-frequency radio wave to ``carry'' the low-frequency information signal by varying one of its properties. This section explains why modulation is unavoidable, describes the two classical methods --- amplitude modulation (AM) and frequency modulation (FM) --- and shows how the information is recovered at the receiver by demodulation.

\courssub{Why Modulation is Necessary}

Imagine a radio station in Douala trying to broadcast the audio signal from a microphone directly, without modulation. Three serious problems arise.

\textbf{1. Audio frequencies radiate very poorly.} The human voice and music occupy roughly the band \SI{20}{Hz} to \SI{20}{kHz}. Electromagnetic waves at such low frequencies are radiated extremely inefficiently by any practical aerial: the radiated power is proportional to $(l/\lambda)^2$ for a short dipole, so it falls off dramatically as frequency decreases (wavelength $\lambda$ increases).

\textbf{2. The aerial would be impossibly long.} An efficient transmitting aerial has a length comparable to a quarter or a half of the wavelength, $\lambda = c/f$. For an audio frequency of $f = \SI{1}{kHz}$:
\[
\lambda = \frac{c}{f} = \frac{3.00\times 10^{8}}{1.0\times 10^{3}} = 3.0\times 10^{5}\ \mathrm{m} = \SI{300}{km}.
\]
A quarter-wave aerial would be \SI{75}{km} long --- longer than the distance from Yaound\'e to Bafoussam! For a carrier at \SI{100}{MHz}, $\lambda = \SI{3}{m}$ and a simple aerial of about \SI{75}{cm} suffices.

\textbf{3. All stations would overlap.} If every station transmitted the same audio band (\SI{20}{Hz}--\SI{20}{kHz}), all broadcasts would occupy the same frequencies and no receiver could separate them. By giving each station its own high-frequency \cle{carrier}, the tuning circuit of Section~2 can select one station and reject the others.

\begin{definition}[Modulation, carrier wave, demodulation]
\cle{Modulation} is the process of varying a property (amplitude, frequency or phase) of a high-frequency \cle{carrier wave} in step with the low-frequency information (modulating) signal, so that the carrier ``carries'' the information. The result is the \cle{modulated wave}. \cle{Demodulation} (or detection) is the reverse process, carried out at the receiver, of extracting the information signal from the modulated carrier.
\end{definition}

\courssub{Amplitude Modulation (AM)}

\textbf{Intuition.} Take a carrier $v_c = V_c\sin(2\pi f_c t)$ with $f_c$ in the hundreds of kHz (medium wave) and make its \emph{amplitude} rise and fall in step with the audio signal $v_m = V_m\sin(2\pi f_m t)$, where $f_m \ll f_c$. The ``outline'' (envelope) of the rapid oscillations then reproduces the shape of the audio signal.

\begin{propriete}[AM wave and modulation index]
For a single-tone modulating signal, an amplitude-modulated wave has the form
\[
v(t) = \bigl[V_c + V_m\sin(2\pi f_m t)\bigr]\sin(2\pi f_c t).
\]
The \cle{modulation index} (or depth of modulation) is
\[
m = \frac{V_m}{V_c} = \frac{A_{\max}-A_{\min}}{A_{\max}+A_{\min}},
\]
where $A_{\max}$ and $A_{\min}$ are the maximum and minimum peak-to-peak values of the envelope. For undistorted AM, $0 \le m \le 1$ (often quoted as a percentage, $m\times 100\,\%$). If $m > 1$ the envelope reaches zero and the signal is \cle{overmodulated}: the envelope no longer matches the information signal and the received sound is distorted.
\end{propriete}

\begin{popfigure}
\begin{center}
\begin{tikzpicture}[scale=1.0]
% Modulating signal
\begin{scope}[shift={(0,4.2)}]
\draw[->] (0,0) -- (10.4,0) node[right]{\small $t$};
\draw[->] (0,-1.3) -- (0,1.5);
\draw[domain=0:10,smooth,samples=200,popblue,thick] plot (\x,{sin(360*\x/5)});
\node[popblue,left] at (0,0.9) {\small modulating signal};
\end{scope}
% AM wave
\begin{scope}[shift={(0,0)}]
\draw[->] (0,0) -- (10.4,0) node[right]{\small $t$};
\draw[->] (0,-1.6) -- (0,1.8);
\draw[domain=0:10,smooth,samples=800,poporange,thick] plot (\x,{(1+0.7*sin(360*\x/5))*sin(360*\x*4)});
\draw[domain=0:10,smooth,samples=200,popdark,dashed] plot (\x,{1+0.7*sin(360*\x/5)});
\draw[domain=0:10,smooth,samples=200,popdark,dashed] plot (\x,{-1-0.7*sin(360*\x/5)});
\node[poporange,left] at (0,1.2) {\small AM wave};
\end{scope}
% FM wave
\begin{scope}[shift={(0,-4.2)}]
\draw[->] (0,0) -- (10.4,0) node[right]{\small $t$};
\draw[->] (0,-1.5) -- (0,1.7);
\draw[domain=0:10,smooth,samples=1200,popgreen,thick] plot (\x,{sin(360*\x*4 - 6*cos(360*\x/5))});
\node[popgreen,left] at (0,1.1) {\small FM wave};
\end{scope}
\end{tikzpicture}
\end{center}
\legende{The same modulating (audio) signal impressed on a carrier: in AM the \emph{amplitude} follows the signal (dashed envelope); in FM the \emph{frequency} rises and falls while the amplitude stays constant.}
\end{popfigure}

\begin{propriete}[Frequency spectrum of an AM wave --- sidebands]
Using the identity $\sin A \sin B = \frac{1}{2}[\cos(A-B) - \cos(A+B)]$, the AM wave expands to three frequency components:
\[
v(t) = V_c\sin(2\pi f_c t) + \frac{mV_c}{2}\sin\bigl[2\pi(f_c+f_m)t\bigr] - \frac{mV_c}{2}\sin\bigl[2\pi(f_c-f_m)t\bigr].
\]
The spectrum consists of:
\begin{itemize}
\item The \cle{carrier} at $f_c$ (amplitude $V_c$, unchanged).
\item An \cle{upper sideband} (USB) at $f_c+f_m$ (amplitude $mV_c/2$).
\item A \cle{lower sideband} (LSB) at $f_c-f_m$ (amplitude $mV_c/2$).
\end{itemize}
The \cle{bandwidth} required is $B = 2f_m$. For audio up to \SI{5}{kHz}, $B = \SI{10}{kHz}$.
\end{propriete}

\begin{propriete}[Power in an AM wave]
The total transmitted power $P_t$ (into a load $R$) is the sum of carrier and sideband powers:
\[
P_t = P_c + P_{USB} + P_{LSB} = \frac{V_c^2}{2R} + \frac{(mV_c/2)^2}{2R} + \frac{(mV_c/2)^2}{2R} = P_c\left(1 + \frac{m^2}{2}\right).
\]
At $m=1$ (maximum undistorted modulation), the sidebands carry $1/3$ of the total power and the carrier $2/3$. The carrier itself conveys no information; this is a fundamental inefficiency of AM.
\end{propriete}

\begin{methode}[Reading the modulation index from an oscilloscope trace]
\begin{enumerate}
\item Display the AM wave on an oscilloscope (time base slow enough to see the envelope clearly).
\item Measure $A_{\max}$, the peak-to-peak height of the envelope at its widest, and $A_{\min}$, the peak-to-peak height at its narrowest.
\item Compute $m = \dfrac{A_{\max}-A_{\min}}{A_{\max}+A_{\min}}$.
\item Check: if $A_{\min} = 0$ then $m = 1$; if the trace pinches below zero (envelope crosses itself), the wave is overmodulated ($m>1$).
\end{enumerate}
\end{methode}

\begin{exoresolu}[Modulation index of a broadcast signal]
An oscilloscope connected to an AM transmitter shows an envelope whose maximum height is $A_{\max} = \SI{9.0}{cm}$ and minimum height $A_{\min} = \SI{3.0}{cm}$. (a) Find the modulation index. (b) The carrier amplitude is $V_c = \SI{5.0}{V}$; find the amplitude $V_m$ of the modulating signal.
\tcblower
\textbf{Solution.}
(a) Apply the definition:
\[
m = \frac{A_{\max}-A_{\min}}{A_{\max}+A_{\min}} = \frac{9.0-3.0}{9.0+3.0} = \frac{6.0}{12.0} = 0.50.
\]
The wave is $50\,\%$ modulated, comfortably below overmodulation.

(b) Since $m = V_m/V_c$,
\[
V_m = mV_c = 0.50\times 5.0 = \SI{2.5}{V}.
\]
\end{exoresolu}

\courssub{Frequency Modulation (FM)}

In frequency modulation the \emph{amplitude of the carrier is kept strictly constant}; it is the \emph{instantaneous frequency} that varies in step with the information signal:
\[
f(t) = f_c + k\,v_m(t),
\]
where $k$ is a constant of the modulator (Hz/V). When the audio signal is at its positive peak the frequency is highest; at the negative peak it is lowest. The maximum departure from the carrier frequency is called the \cle{frequency deviation} $\Delta f$ (for FM broadcast radio, $\Delta f = \SI{75}{kHz}$ around a carrier in the band \SI{88}{MHz}--\SI{108}{MHz}). The loudness of the sound is encoded in the \emph{size} of the frequency swing ($\Delta f \propto V_m$), and the pitch of the sound in the \emph{rate} of the swing ($f_m$).

\begin{attention}[Common mistake: do not vary the FM envelope]
A very frequent error in sketches is to draw the FM wave with a varying amplitude like AM. In FM the envelope is \textbf{flat}: every cycle has the same height. Only the spacing between cycles (the instantaneous frequency) changes. Examiners penalise FM waves drawn with a wavy envelope.
\end{attention}

\begin{propriete}[Bandwidth of an FM wave --- Carson's rule]
Unlike AM, an FM wave has an infinite number of sidebands (Bessel functions). However, most power is contained within a bandwidth approximated by \cle{Carson's rule}:
\[
B \approx 2(\Delta f + f_m).
\]
For commercial FM radio ($\Delta f = \SI{75}{kHz}$, $f_m \le \SI{15}{kHz}$), $B \approx \SI{180}{kHz}$. Channels are spaced \SI{200}{kHz} apart. This wide bandwidth is the price paid for FM's noise immunity and high fidelity.
\end{propriete}

\courssub{Demodulation: Recovering the Information}

At the receiver, after the tuning circuit has selected the station and amplifiers have strengthened the signal, the information must be separated from the carrier.

\textbf{AM: the envelope (diode) detector.} The AM signal is passed through a diode, which removes the negative half-cycles (half-wave rectification). The remaining pulses are smoothed by a capacitor $C$ in parallel with a resistor $R$ (the load): the capacitor charges to each peak and discharges slowly through $R$ between peaks, so the output voltage follows the \emph{envelope} --- that is, the original audio signal. The time constant must satisfy
\[
\frac{1}{f_c} \ll RC \ll \frac{1}{f_m},
\]
so that the carrier ripple is removed but the audio variations are preserved.

\begin{popfigure}
\begin{center}
\begin{tikzpicture}[european,scale=1.0]
\draw (0,0) to[sV=$v_{AM}$] (0,2.4);
\draw (0,2.4) to[D,l={$D$}] (3,2.4);
\draw (3,2.4) -- (3.6,2.4);
\draw (3.6,2.4) to[C=$C$] (3.6,0);
\draw (3.6,2.4) -- (5.4,2.4) to[R=$R$] (5.4,0);
\draw (5.4,2.4) -- (6.6,2.4) node[right]{\small audio out};
\draw (0,0) -- (6.6,0);
\end{tikzpicture}
\end{center}
\legende{Simple diode envelope detector: the diode rectifies the AM wave; the $RC$ network smooths it so the output follows the envelope (the audio signal).}
\end{popfigure}

\textbf{FM: principle of detection.} Since the information in FM lies in the frequency variations, an FM detector converts frequency changes into amplitude changes, and then recovers the audio signal. Classic methods include:
\begin{itemize}
\item \cle{Slope detection}: a tuned circuit deliberately detuned so the signal sits on the slope of its resonance curve; frequency changes become amplitude changes.
\item \cle{Foster-Seeley discriminator} / \cle{Ratio detector}: balanced circuits giving a voltage proportional to frequency deviation.
\item \cle{Phase-Locked Loop (PLL)}: modern integrated-circuit method where a VCO locks to the input frequency; the control voltage is the demodulated audio.
\end{itemize}
Crucially, before detection the FM wave passes through a \cle{limiter} stage (a saturated amplifier) that clips the wave to a fixed amplitude, removing all amplitude variations caused by noise.

\begin{popfigure}
\begin{center}
\begin{tikzpicture}[scale=1.0, every node/.style={font=\small},
box/.style={draw=popblue, thick, rounded corners, minimum height=0.8cm, minimum width=1.8cm, align=center, fill=popblueL}]
\node[box] (lim) at (0,0) {Limiter};
\node[box, align=center] (det) at (3,0){Frequency-to-\\Amplitude Converter};
\node[box] (lp) at (6,0) {Low-Pass Filter};
\draw[->, thick, popdark] (lim) -- (det);
\draw[->, thick, popdark] (det) -- (lp);
\node[above, popmuted] at (1.5,0.6) {\scriptsize constant amplitude};
\node[above, popmuted] at (4.5,0.6) {\scriptsize audio voltage};
\end{tikzpicture}
\end{center}
\legende{Block diagram of an FM demodulator: the limiter removes amplitude noise, the converter translates frequency swings into voltage swings, and the low-pass filter removes residual carrier components.}
\end{popfigure}

\courssub{AM versus FM: Advantages and Disadvantages}

\begin{propriete}[Why FM sounds cleaner]
Most natural and man-made interference --- lightning flashes, sparks from car ignition systems, electric motors --- adds unwanted \emph{amplitude} variations to the received wave. An AM receiver interprets these as part of the signal, so they are heard as crackles and hiss. An FM receiver first passes the wave through a limiter that removes all amplitude variations, so this noise is rejected; only the frequency variations, which carry the programme, survive. This is why FM radio sounds noticeably cleaner than AM.
\end{propriete}

\begin{center}
\begin{tabularx}{\linewidth}{|l|X|X|}
\hline
\rowcolor{popblueL} & \textbf{AM} & \textbf{FM} \\
\hline
Property varied & Carrier amplitude & Carrier frequency \\
\hline
Bandwidth & Narrow ($\approx 2f_m$, about \SI{10}{kHz}) & Wide ($\approx 2(\Delta f + f_m)$, about \SI{200}{kHz}) \\
\hline
Noise immunity & Poor; crackle and hiss from amplitude noise & Excellent (limiter rejects amplitude noise) \\
\hline
Sound quality & Lower fidelity (restricted audio bandwidth) & High fidelity (wide audio band, stereo capable) \\
\hline
Range & Long: ground and sky waves, beyond horizon & Shorter: space wave, roughly line of sight \\
\hline
Transmitter efficiency & Low (carrier wastes power) & Higher (constant amplitude, class C amps) \\
\hline
Cost/complexity & Simple, cheap transmitters and receivers & More complex, costlier equipment \\
\hline
\end{tabularx}
\end{center}

\begin{attention}[GCE exam tip]
``Compare AM and FM, giving one advantage of each over the other'' is a frequent structured-question item at Advanced Level. Learn the comparison table above, and always justify: AM's long range comes from ground/sky-wave propagation at lower carrier frequencies (MF/HF); FM's quality comes from noise limiting and its large bandwidth, which is also why FM needs many more channels and is restricted to VHF line-of-sight coverage.
\end{attention}

\courssub{The Complete Radio Communication Chain}

We can now assemble the whole system studied in this chapter into a single block diagram.

\begin{popfigure}
\begin{center}
\begin{tikzpicture}[scale=1.0, every node/.style={font=\small},
box/.style={draw=popblue, thick, rounded corners, minimum height=0.9cm, minimum width=1.9cm, align=center, fill=popblueL}]
\node[box] (mic) at (0,0) {microphone};
\node[box] (mod) at (2.8,0) {modulator};
\node[box] (tx) at (5.6,0) {transmitter + aerial};
\node[box] (ch) at (8.6,0) {channel (radio wave)};
\node[box] (rx) at (11.6,0) {receiver + tuning};
\node[box] (dem) at (14.4,0) {demodulator};
\node[box] (ls) at (17.0,0) {loudspeaker};
\draw[->, thick, popdark] (mic) -- (mod);
\draw[->, thick, popdark] (mod) -- (tx);
\draw[->, thick, popdark] (tx) -- (ch);
\draw[->, thick, popdark] (ch) -- (rx);
\draw[->, thick, popdark] (rx) -- (dem);
\draw[->, thick, popdark] (dem) -- (ls);
\node[below, popmuted] at (2.8,-0.7) {\scriptsize carrier + audio};
\node[below, popmuted] at (8.6,-0.7) {\scriptsize modulated wave};
\node[below, popmuted] at (14.4,-0.7) {\scriptsize audio recovered};
\end{tikzpicture}
\end{center}
\legende{Block diagram of a complete radio communication system: the audio signal modulates a carrier, the modulated wave crosses the channel, and the receiver selects, amplifies and demodulates it to restore the sound.}
\end{popfigure}

\begin{aretenir}[Key points]
\begin{itemize}
\item Modulation is necessary because audio frequencies radiate poorly, would need kilometre-long aerials, and would make all stations overlap.
\item AM: the carrier \emph{amplitude} follows the signal; $m = \dfrac{A_{\max}-A_{\min}}{A_{\max}+A_{\min}}$, with $m\le 1$ to avoid overmodulation. Spectrum: carrier $\pm$ two sidebands; bandwidth $2f_m$. Power $P_t = P_c(1+m^2/2)$.
\item FM: the carrier \emph{frequency} follows the signal; amplitude is constant; maximum frequency swing is deviation $\Delta f$. Bandwidth $\approx 2(\Delta f + f_m)$ (Carson's rule).
\item AM demodulation: diode envelope detector with $RC$ smoothing ($1/f_c \ll RC \ll 1/f_m$).
\item FM demodulation: limiter + frequency-to-amplitude converter (discriminator, ratio detector, or PLL).
\item AM: cheap, long range (ground/sky wave), but noisy and low fidelity. FM: noise-immune, high fidelity, but wide bandwidth and line-of-sight range.
\end{itemize}
\end{aretenir}

\begin{exercice}[AM and FM reasoning]
\begin{enumerate}
\item An AM transmitter has a carrier of amplitude \SI{8.0}{V} modulated by an audio signal of amplitude \SI{6.0}{V}. Calculate the modulation index and state whether the transmission is distorted.
\item Explain, with reference to the limiter stage, why a thunderstorm causes crackles on an AM radio but barely affects an FM receiver.
\item An FM station broadcasts on a carrier of \SI{101.5}{MHz} with a deviation of $\pm\SI{75}{kHz}$. Between which frequencies does the instantaneous frequency vary?
\item (Challenge) An AM broadcast station operates at \SI{1000}{kHz} with a maximum audio frequency of \SI{5}{kHz}. What are the frequencies of the upper and lower sidebands? What is the bandwidth?
\end{enumerate}
\end{exercice}

\begin{corrige}[Exercise 1]
\begin{enumerate}
\item $m = V_m/V_c = 6.0/8.0 = 0.75$. Since $m<1$, the modulation is undistorted.
\item Lightning produces electromagnetic pulses that add spurious \emph{amplitude} variations. The AM receiver's envelope detector reads these as signal, hence crackles. The FM receiver's limiter clips all amplitude variations before detection; only frequency variations carry the programme, so the interference is rejected.
\item $f$ varies between $f_c - \Delta f = 101.5 - 0.075 = \SI{101.425}{MHz}$ and $f_c + \Delta f = \SI{101.575}{MHz}$.
\item USB: $1000 + 5 = \SI{1005}{kHz}$. LSB: $1000 - 5 = \SI{995}{kHz}$. Bandwidth $= 2 \times 5 = \SI{10}{kHz}$.
\end{enumerate}
\end{corrige}

\begin{savaistu}[Radio in Cameroon]
Cameroon's CRTV radio and many private stations broadcast both on medium wave (AM) and on FM. In large towns like Douala and Yaound\'e listeners prefer FM for its clear, stereo sound, but in remote rural areas of the North or East regions medium-wave AM is still valued because its ground and sky waves reach villages far beyond the horizon where no FM transmitter is visible. The next section will show how digital techniques and satellites extend this coverage even further.
\end{savaistu}

\coursec{Digital Conversion and Communication Channels}

In the previous section we saw how a message is carried by a high-frequency carrier through \cle{amplitude modulation} or \cle{frequency modulation}. But before a signal even reaches the transmitter, modern communication systems usually convert it into \cle{digital} form --- a stream of numbers. Your mobile phone in Yaound\'e does not send your voice as a continuously varying wave; it sends numbers, coded in binary, at a rate of thousands of numbers per second. In this final section we study how analogue signals are converted to digital and back, why engineers prefer digital transmission, and how channels --- cables, radio links, optical fibres and satellites --- carry the information to its destination.

\courssub{Analogue and Digital Signals}

\begin{definition}[Analogue and digital signals]
An \cle{analogue signal} varies \emph{continuously} in time and can take \emph{any} value within a range (example: the voltage from a microphone).\\
A \cle{digital signal} can take only a \emph{finite number of discrete values}, usually two (``0'' and ``1'', i.e. low/high voltage). The information is carried as a sequence of \cle{bits} (binary digits).
\end{definition}

An analogue signal carries its information in the exact shape of the wave; any noise added to that shape is indistinguishable from the message. A digital signal carries its information in a \emph{code}; as long as a receiver can still tell a ``0'' from a ``1'', the message is perfectly recovered. This single difference explains almost all the advantages of digital transmission studied below.

\courssub{Analogue-to-Digital Conversion (ADC)}

\textbf{The problem.} An analogue signal varies continuously in time and can take any value within a range. A computer or a digital transmitter can only handle \emph{numbers} with a \emph{finite} number of digits. Conversion therefore requires three steps: \cle{sampling}, \cle{quantisation} and \cle{coding}.

\textbf{Step 1: Sampling.} Instead of describing the signal at every instant, we measure (``sample'') its value at regular time intervals $T_s$, called the \cle{sampling period}.

\begin{definition}[Sampling frequency]
The \cle{sampling frequency} $f_s$ is the number of samples taken per second:
\[ f_s = \frac{1}{T_s} \qquad \text{(unit: hertz, Hz)}. \]
\end{definition}

How fast must we sample? If we sample too slowly, rapid wiggles of the signal are missed and the reconstructed signal is distorted (an effect called \emph{aliasing}). The answer is given by the Nyquist criterion.

\begin{propriete}[Nyquist (sampling) criterion]
To reconstruct a signal faithfully, the sampling frequency must be at least twice the highest frequency present in the signal:
\[ \boxed{\,f_s \geq 2\,f_{\max}\,} \]
The minimum value $f_s = 2f_{\max}$ is called the \emph{Nyquist rate}. (Proof not examinable; the idea is that at least two points per cycle of the fastest component are needed to capture it.)
\end{propriete}

\begin{exemplebox}[Choosing a sampling frequency]
The human ear perceives frequencies up to about $f_{\max} = \SI{20}{kHz}$. To digitise music, we need $f_s \geq 2 \times 20 = \SI{40}{kHz}$. This is why audio CDs use $f_s = \SI{44.1}{kHz}$ --- slightly above the Nyquist rate, leaving a safety margin.
\end{exemplebox}

\textbf{Step 2: Quantisation.} Each sampled value is rounded to the nearest of a finite set of allowed \emph{levels}.

\begin{definition}[Quantisation]
\cle{Quantisation} is the process of approximating each sampled value by the nearest of a set of discrete levels. If each sample is coded on $n$ \cle{bits}, the number of levels is
\[ \boxed{\,N = 2^{n}\,} \]
and the \emph{quantisation step} (resolution) over a voltage range $V_{\max}-V_{\min}$ is
\[ q = \frac{V_{\max}-V_{\min}}{2^{n}}. \]
\end{definition}

\textbf{Step 3: Coding.} Each level is given a binary number (a ``word'') of $n$ bits. The output of the ADC is therefore a sequence of binary words --- a digital signal.

\begin{methode}[Number of quantisation levels, resolution and bit rate]
\begin{enumerate}
\item Identify the number of bits $n$ per sample.
\item Compute the number of levels: $N = 2^{n}$.
\item Compute the step: $q = \dfrac{V_{\max}-V_{\min}}{N}$.
\item The bit rate of the digitised signal is $D = n \times f_s$ (bits per second, $\mathrm{bit\,s^{-1}}$).
\end{enumerate}
\end{methode}

\begin{exoresolu}[Digitising a microphone signal]
A microphone produces a voltage between $0$ and $\SI{8}{V}$. It is digitised with a 3-bit ADC at a sampling frequency $f_s = \SI{8}{kHz}$.\\
(a) How many quantisation levels are there? (b) Find the quantisation step. (c) Find the bit rate. (d) What binary word represents a sample of value $\SI{5.6}{V}$?
\tcblower
\textbf{Solution.}\\
(a) $N = 2^{3} = 8$ levels.\\
(b) $q = \dfrac{8 - 0}{2^{3}} = \SI{1}{V}$. The levels are $0, 1, 2, \dots, 7\ \mathrm{V}$, coded $000, 001, \dots, 111$.\\
(c) Bit rate: $D = n f_s = 3 \times 8000 = \SI{2.4e4}{bit.s^{-1}} = \SI{24}{kbit.s^{-1}}$.\\
(d) $5.6\ \mathrm{V}$ lies between levels 5 and 6; the nearest level is $\SI{6}{V}$, coded $\mathbf{110}$.
\end{exoresolu}

\begin{popfigure}
\begin{center}
\begin{tikzpicture}
\begin{axis}[
  width=12cm, height=6.5cm,
  xlabel={$t$ (ms)}, ylabel={$u$ (V)},
  xmin=0, xmax=8, ymin=0, ymax=8,
  xtick={0,1,2,3,4,5,6,7,8},
  ytick={0,1,2,3,4,5,6,7,8},
  grid=both, grid style={popline, very thin},
  legend style={at={(0.02,0.98)}, anchor=north west, font=\small}
]
\addplot[popblue, thick, domain=0:8, samples=200] {4 + 3*sin(deg(0.9*x))};
\addlegendentry{analogue signal}
\addplot[only marks, mark=*, poporange, mark size=2pt] coordinates {
(0,4.0) (1,6.3) (2,6.9) (3,5.2) (4,2.5) (5,1.1) (6,1.9) (7,4.3) (8,6.0)};
\addlegendentry{samples ($f_s = 1$ kHz)}
\addplot[const plot, poppurple, thick] coordinates {
(0,4) (1,6) (2,7) (3,5) (4,3) (5,1) (6,2) (7,4) (8,6) (9,6)};
\addlegendentry{quantised (3-bit)}
\end{axis}
\end{tikzpicture}
\end{center}
\legende{Sampling and 3-bit quantisation of an analogue signal. The continuous curve is sampled at regular intervals; each sample is rounded to the nearest of the $2^3 = 8$ allowed levels.}
\end{popfigure}

\courssub{Digital-to-Analogue Conversion (DAC)}

At the receiver, the reverse operation is needed: the stream of binary words must be turned back into a continuously varying voltage to drive a loudspeaker or a screen.

\begin{definition}[Digital-to-analogue conversion]
A \cle{DAC} reads each binary word and outputs the corresponding voltage level, holding it until the next word arrives. The raw output is a \emph{staircase} signal; a low-pass \cle{smoothing filter} then removes the sharp steps and restores a smooth analogue waveform.
\end{definition}

The staircase is only an approximation of the original signal. The difference between the staircase and the true signal is the \emph{quantisation error}; it is reduced by using more bits per sample (smaller step $q$) and a higher sampling frequency.

\begin{attention}[Quantisation is irreversible]
Rounding to the nearest level destroys information: two different samples can be coded by the same word. No DAC can recover the lost detail. This is why the number of bits must be chosen \emph{before} transmission.
\end{attention}

\courssub{Digital versus Analogue Transmission}

\begin{center}
\begin{tabularx}{\linewidth}{|l|X|X|}
\hline
\rowcolor{popblueL} & \textbf{Analogue} & \textbf{Digital} \\
\hline
Noise & Noise added along the line is amplified with the signal and accumulates. & \cle{Repeaters} \emph{regenerate} clean pulses: they only need to decide ``0 or 1'', so noise does not accumulate. \\
\hline
\rowcolor{popgreenL} Errors & Impossible to detect or correct. & Extra \emph{parity/check bits} allow \cle{error detection and correction}. \\
\hline
Capacity & One signal per channel. & \cle{Multiplexing}: many digital signals interleaved on one channel; \cle{compression} reduces the bit rate needed. \\
\hline
\rowcolor{popgreenL} Cost/complexity & Simple circuits. & Needs ADC/DAC and more \cle{bandwidth} for the same message. \\
\hline
\end{tabularx}
\end{center}

\begin{aretenir}[Why digital won]
Digital signals can be \textbf{regenerated} (noise removed at each repeater), \textbf{protected} (error-correcting codes), \textbf{compressed} and \textbf{multiplexed}. Their price: greater bandwidth and conversion electronics.
\end{aretenir}

\courssub{Bandwidth and Sidebands}

Every channel (cable, radio link, fibre) can carry only a limited range of frequencies.

\begin{definition}[Bandwidth]
The \cle{bandwidth} $B$ of a channel is the width of the range of frequencies it can transmit:
\[ B = f_{\text{high}} - f_{\text{low}} \qquad \text{(unit: Hz)}. \]
A signal whose spectrum is wider than the channel bandwidth is distorted.
\end{definition}

\textbf{Sidebands in AM.} When a carrier $f_c$ is amplitude-modulated by a sinusoidal signal of frequency $f_m$, the modulated wave is
\[ u(t) = U_c\cos(2\pi f_c t) + \tfrac{mU_c}{2}\cos\bigl(2\pi(f_c - f_m)t\bigr) + \tfrac{mU_c}{2}\cos\bigl(2\pi(f_c + f_m)t\bigr), \]
where $m$ is the modulation index. The spectrum therefore contains \emph{three} frequencies: the \cle{carrier} $f_c$ and two \cle{sidebands} at $f_c - f_m$ and $f_c + f_m$. For a real message occupying frequencies up to $f_{m,\max}$, the sidebands extend from $f_c - f_{m,\max}$ to $f_c + f_{m,\max}$, so the AM signal occupies
\[ \boxed{\,B_{\text{AM}} = 2\,f_{m,\max}\,} \]

\begin{attention}[The classic trap]
The AM bandwidth counts \textbf{both} sidebands: it is $2f_m$, \textbf{not} $f_m$. If speech occupies up to $\SI{4.5}{kHz}$, the AM channel needs $B = 2 \times 4.5 = \SI{9}{kHz}$.
\end{attention}

\begin{popfigure}
\begin{center}
\begin{tikzpicture}[scale=1.0]
\draw[->, thick] (0,0) -- (11,0) node[below left] {frequency};
\draw[->, thick] (0,0) -- (0,3.2) node[above left] {power};
% carrier
\draw[popblue, very thick] (5.5,0) -- (5.5,2.8);
\node[popblue, above] at (5.5,2.8) {carrier $f_c$};
% sidebands
\draw[poporange, very thick] (3.5,0) -- (3.5,1.4);
\draw[poporange, very thick] (7.5,0) -- (7.5,1.4);
\node[poporange, above, align=center] at (3.5,1.4) {lower sideband\\ $f_c - f_m$};
\node[poporange, above, align=center] at (7.5,1.4) {upper sideband\\ $f_c + f_m$};
% bandwidth arrow
\draw[<->, poppurple, thick] (3.5,-0.6) -- (7.5,-0.6);
\node[poppurple, below] at (5.5,-0.6) {bandwidth $B = 2f_m$};
\draw[<->, gray] (3.5,0.5) -- (5.5,0.5);
\node[gray, above, font=\small] at (4.5,0.5) {$f_m$};
\end{tikzpicture}
\end{center}
\legende{Frequency spectrum of an AM signal: carrier at $f_c$ and two sidebands at $f_c \pm f_m$. The total bandwidth is $2f_m$. Most of the power sits in the carrier, which carries no information.}
\end{popfigure}

\textbf{Power distribution.} In AM, the carrier contains most of the transmitted power yet carries \emph{no} information; only the sidebands do (for sinusoidal modulation of index $m$, each sideband carries a fraction $m^2/4$ of the carrier power). This is one weakness of AM. In \cle{FM}, the information is coded in frequency deviations; the spectrum contains many side frequencies, and Carson's rule gives approximately
\[ B_{\text{FM}} \approx 2(\Delta f + f_{m,\max}), \]
where $\Delta f$ is the peak frequency deviation. Since $\Delta f$ is large (e.g. $\SI{75}{kHz}$ for FM radio), FM needs a much \textbf{wider bandwidth} than AM --- the price paid for its superior immunity to noise.

\begin{exoresolu}[AM channel bandwidth]
A radio station broadcasts on a carrier $f_c = \SI{1.0}{MHz}$. The audio signal contains frequencies up to $\SI{5}{kHz}$.\\
(a) Sketch the spectrum and give the frequencies of the extreme sideband components. (b) Find the bandwidth of the channel. (c) How many such stations fit between 530 and $\SI{1600}{kHz}$?
\tcblower
\textbf{Solution.}\\
(a) Sidebands extend from $f_c - f_m = 1000 - 5 = \SI{995}{kHz}$ to $f_c + f_m = \SI{1005}{kHz}$.\\
(b) $B = 2 f_m = 2 \times 5 = \SI{10}{kHz}$.\\
(c) Available band: $1600 - 530 = \SI{1070}{kHz}$. Number of channels: $N = \dfrac{1070}{10} = 107$ stations.
\end{exoresolu}

\courssub{Satellite Communication}

Microwaves travel in straight lines and are not reflected by the ionosphere, so Earth-to-Earth microwave links are limited by the horizon. The solution is to place a relay in the sky: a \cle{communication satellite}.

\begin{propriete}[Geostationary orbit]
A \cle{geostationary satellite} orbits in the \textbf{equatorial plane}, in the same direction as the Earth's rotation, with a period $T = \SI{24}{h}$. Its altitude is about $\SI{36000}{km}$ above the equator (orbital radius $\approx \SI{42000}{km}$ from the Earth's centre). It therefore appears \emph{fixed} above one point of the equator, so ground antennas can point at it permanently. (The altitude follows from Kepler's third law: $T^2 = \dfrac{4\pi^2 r^3}{GM_T}$.)
\end{propriete}

\begin{exemplebox}[Checking the geostationary radius]
With $T = 24 \times 3600 = \SI{8.64e4}{s}$, $G = \SI{6.67e-11}{N.m^2.kg^{-2}}$, $M_T = \SI{6.0e24}{kg}$:
\[ r^3 = \frac{G M_T T^2}{4\pi^2} = \frac{6.67\times10^{-11}\times 6.0\times10^{24}\times (8.64\times10^{4})^2}{39.5} \approx 7.6\times10^{22}\ \mathrm{m^3} \]
so $r \approx \SI{4.2e7}{m} = \SI{42000}{km}$, i.e. an altitude of about $\SI{36000}{km}$ above the surface. $\checkmark$
\end{exemplebox}

\textbf{Uplink and downlink.} The ground station transmits to the satellite on the \cle{uplink} frequency; the satellite amplifies the signal, shifts its frequency, and retransmits to Earth on the \cle{downlink} frequency. Different frequencies (e.g. around $\SI{6}{GHz}$ up / $\SI{4}{GHz}$ down, or 14/12 GHz) are used so that the powerful uplink does not swamp the weak downlink received by the satellite.

\textbf{Advantages:} one satellite covers about a third of the globe --- ideal for continental TV distribution, telephony and internet in regions without cables (vital for rural Cameroon). \textbf{Limitations:} the round trip of $2 \times \SI{36000}{km}$ at $c = \SI{3.00e8}{m.s^{-1}}$ introduces a delay
\[ \Delta t = \frac{2 \times 3.6\times10^{7}}{3.00\times10^{8}} \approx \SI{0.24}{s}, \]
noticeable in telephone conversations; launching and maintaining satellites is costly; signals are weakened by rain at high frequencies.

\courssub{Mobile Telephony: Cells and Base Stations}

A mobile phone is a low-power radio transmitter: it cannot reach a whole country. The solution is the \cle{cellular network}: the territory is divided into small \cle{cells}, each served by a \cle{base station} (the masts you see along roads and on rooftops in Douala or Bamenda).

\begin{definition}[Base station and its role]
A \cle{base station} is a fixed radio transceiver that links the mobile handsets within its cell to the rest of the network (switching centres, other cells, the internet). Its roles are:
\begin{itemize}
\item to receive the weak signals from handsets and relay them into the network;
\item to transmit network signals to the handsets;
\item to manage \cle{frequency reuse}: the same frequencies are reused in distant, non-adjacent cells without interference, multiplying the capacity of the network;
\item to perform \cle{handover}: as a user moves (e.g. in a taxi), the call is transferred seamlessly from one base station to the next.
\end{itemize}
\end{definition}

\begin{popfigure}
\begin{center}
\begin{tikzpicture}[scale=0.9]
% cells (hexagons)
\foreach \x/\y/\l in {0/0/A, 1.8/0/B, 3.6/0/C, 0.9/1.55/D, 2.7/1.55/E, 4.5/1.55/F, 0.9/-1.55/G, 2.7/-1.55/H}{
  \node[draw=popblue, regular polygon, regular polygon sides=6, minimum size=1.7cm, inner sep=1pt, fill=popblueL] at (\x,\y) {};
  \node[popdark, font=\small] at (\x,\y) {\l};
  \fill[poppurple] (\x,\y) ++(0.55,0.32) circle (2.2pt);
}
\node[font=\small, popmuted, align=center] at (2.25,-2.6) {Cellular map: each cell has a base station (purple dot).\\ Frequencies of cell A are reused in cell F, far away.};
% satellite link
\begin{scope}[xshift=8.5cm, yshift=0.3cm]
\draw[popblue, thick] (0,0) circle (1.1);
\draw[popblue, dashed] (0,0) ellipse (1.1 and 0.35);
\node[font=\small, below] at (0,-1.25) {Earth};
\draw[popgreen, dashed, thick] (0,0) circle (2.2);
\fill[poporange] (2.2,0) circle (3pt);
\node[poporange, right, font=\small, align=left] at (2.3,0.1) {geostationary\\ satellite};
\draw[<->, poppurple, thick] (1.15,0.55) -- (2.05,0.15);
\node[poppurple, font=\small, above right, align=left] at (0.9,0.75) {uplink /\\ downlink};
\draw[<->, gray] (0,1.3) -- (2.2,1.3);
\node[gray, font=\small, above] at (1.1,1.3) {$\approx \SI{36000}{km}$};
\end{scope}
\end{tikzpicture}
\end{center}
\legende{Left: a region divided into hexagonal cells, each with a base station; frequencies are reused in distant cells. Right: a geostationary satellite fixed above the equator, exchanging uplink and downlink signals with a ground station.}
\end{popfigure}

\courssub{A High-Bandwidth Channel: Optical Fibres (GCE link to optics)}

For fixed networks, the champion channel is the \cle{optical fibre}: a hair-thin glass fibre guides light pulses by \cle{total internal reflection} at the core--cladding boundary (the phenomenon you met in geometrical optics: light strikes the interface at an angle greater than the critical angle, so it is trapped inside the core). Because light frequencies ($\sim 10^{14}\ \mathrm{Hz}$) are enormous, the available bandwidth is huge --- a single fibre can carry millions of simultaneous phone calls or high-speed internet, with very low attenuation and total immunity to electromagnetic interference. Submarine fibre-optic cables landing on the Cameroonian coast connect the country to the global internet backbone.

\begin{savaistu}[Did you know?]
Cameroon's international internet traffic arrives largely through submarine optical fibre cables landing near Douala, while domestic links combine fibre, microwave relays and satellite. The mobile phone in your pocket constantly negotiates with the nearest base station --- and hands over to the next one as you travel, without you noticing.
\end{savaistu}

\begin{exoresolu}[Bit rate of a digitised telephone channel (GCE style)]
A telephone voice channel is limited to $f_{\max} = \SI{3.4}{kHz}$. It is sampled at the Nyquist rate and coded on 8 bits per sample.\\
(a) Find the sampling frequency. (b) Find the bit rate of the digital channel. (c) How many such channels can a $\SI{2.048}{Mbit.s^{-1}}$ digital link carry?
\tcblower
\textbf{Solution.}\\
(a) Nyquist: $f_s = 2 f_{\max} = 2 \times 3.4 = \SI{6.8}{kHz}$.\\
(b) $D = n f_s = 8 \times 6.8\times10^{3} = \SI{5.44e4}{bit.s^{-1}} = \SI{54.4}{kbit.s^{-1}}$.\\
(c) $N = \dfrac{2.048\times10^{6}}{5.44\times10^{4}} \approx 37.6$, so \textbf{37 channels} (whole number, rounded down).
\end{exoresolu}

\begin{exercice}[Practice]
\begin{enumerate}
\item An audio signal occupies 20 Hz to 15 kHz. It is digitised on 16 bits. (a) Give the minimum sampling frequency. (b) How many quantisation levels are available? (c) Calculate the bit rate.
\item An AM transmitter uses a carrier of $\SI{800}{kHz}$ and modulating frequencies up to $\SI{4}{kHz}$. State the frequencies present in the spectrum and the channel bandwidth.
\item Explain, in three sentences, why digital repeaters give better long-distance quality than analogue amplifiers.
\item A geostationary satellite is $\SI{36000}{km}$ above the Earth's surface. Calculate the total time delay for a signal going Earth $\to$ satellite $\to$ Earth.
\end{enumerate}
\end{exercice}

\begin{corrige}[Exercise 1]
\begin{enumerate}
\item (a) $f_s \geq 2 \times 15 = \SI{30}{kHz}$. (b) $N = 2^{16} = 65536$ levels. (c) $D = 16 \times 30\,000 = \SI{4.8e5}{bit.s^{-1}} = \SI{480}{kbit.s^{-1}}$.
\item Carrier $\SI{800}{kHz}$; sidebands from $800 - 4 = \SI{796}{kHz}$ to $800 + 4 = \SI{804}{kHz}$; bandwidth $B = 2 f_m = \SI{8}{kHz}$.
\item Analogue amplifiers amplify the noise together with the signal, so noise accumulates along the line. Digital repeaters only decide whether each pulse is a 0 or a 1, then retransmit a clean, full-size pulse. Errors can additionally be detected and corrected with check bits, so quality is preserved over any distance.
\item $\Delta t = \dfrac{2 \times 3.6\times10^{7}\ \mathrm{m}}{3.00\times10^{8}\ \mathrm{m\,s^{-1}}} = \SI{0.24}{s}$.
\end{enumerate}
\end{corrige}

\begin{aretenir}[Section summary]
\begin{itemize}
\item An \textbf{analogue} signal is continuous in time and value; a \textbf{digital} signal takes discrete values coded in bits.
\item ADC = \textbf{sampling} ($f_s \geq 2f_{\max}$, Nyquist) + \textbf{quantisation} ($2^n$ levels, step $q = \Delta V/2^n$) + \textbf{binary coding}. Bit rate $D = n f_s$.
\item DAC = binary words $\to$ staircase voltage $\to$ smoothing filter.
\item Digital transmission: regeneration by repeaters, error correction, compression, multiplexing --- at the cost of bandwidth and converters.
\item Bandwidth $B = f_{\text{high}} - f_{\text{low}}$; AM spectrum = carrier + two sidebands at $f_c \pm f_m$, so $B_{\text{AM}} = 2f_{m,\max}$; FM is wider, $B_{\text{FM}} \approx 2(\Delta f + f_{m,\max})$.
\item Geostationary satellite: equatorial orbit, $T = \SI{24}{h}$, altitude $\approx \SI{36000}{km}$; uplink and downlink on different frequencies; delay $\approx \SI{0.24}{s}$.
\item Cellular telephony: cells, base stations (link handsets to network, frequency reuse, handover). Optical fibres: total internal reflection, enormous bandwidth.
\end{itemize}
\end{aretenir}

\newpage

\coursec{Integration activity}

\courssub{Aims of the activity}

By the end of this activity, you should be able to:
\begin{itemize}
\item \cle{analyse} a real communication problem and identify the physical constraints (distance, terrain, frequency allocation);
\item \cle{select} an appropriate propagation mode (ground wave, sky wave, space wave) for a given link and justify the choice;
\item \cle{calculate} the dimensions of a half-wave dipole aerial for a given carrier frequency;
\item \cle{design} a simple $LC$ tuning circuit and determine the range of capacitance needed to cover a given band;
\item \cle{compare} AM and FM and defend an engineering choice using their advantages and disadvantages;
\item \cle{compute} the bandwidth of an AM signal and the number of channels fitting in a given allocation;
\item \cle{propose} a long-distance relay solution (satellite link or network of base stations) and justify it;
\item \cle{communicate} scientific results in the form of a clear, argued one-page engineering brief.
\end{itemize}

\courssub{Situation}

\textbf{Radio Menchum Community Project.}\par
The NGO \emph{Ndawara Development Trust} has obtained a licence to operate a community radio station, \textbf{Radio Menchum}, with its studio and transmitter located on a hill near \textbf{Bamenda} (North-West Region). The station must serve farming villages up to \SI{80}{km} away across the hilly, forested terrain of the region. In addition, recorded programmes must be relayed every evening to a partner station in \textbf{Yaoundé}, about \SI{280}{km} away, for rebroadcast in the Centre Region.

The engineering team (your group) is given the following data and constraints:

\begin{center}
\begin{tabularx}{\linewidth}{|l|X|}
\hline
\rowcolor{popblueL} \textbf{Item} & \textbf{Data / constraint} \\
\hline
Local broadcast range & \SI{80}{km}, hilly terrain, no line of sight for many villages \\
\hline
Available bands & MF band around \SI{1.0}{MHz} (AM), or VHF FM band $88$--$\SI{108}{MHz}$ \\
\hline
Relay link Bamenda--Yaoundé & \SI{280}{km}, mountainous terrain, no line of sight \\
\hline
Highest audio frequency & $f_{m,\max} = \SI{5.0}{kHz}$ \\
\hline
AM allocation in the region & \SI{200}{kHz} of spectrum reserved for community AM stations \\
\hline
Tuning circuit inductor & $L = \SI{0.50}{\micro H}$ available in the workshop \\
\hline
Speed of electromagnetic waves & $c = \SI{3.00e8}{m.s^{-1}}$ \\
\hline
\end{tabularx}
\end{center}

The Trust asks each group to produce a \textbf{one-page engineering brief} containing: the chosen frequency bands with justification of the propagation mode for each link, the design of the transmitting aerial, the design of the receiver tuning circuit, the choice of modulation with a argued defence, the bandwidth calculations, and a diagram of the complete communication chain (transmitter $\rightarrow$ channel $\rightarrow$ receiver, plus the relay to Yaoundé).

\courssub{Tasks}

\begin{enumerate}
\item \textbf{Propagation modes.} For the local Bamenda--villages link (\SI{80}{km}, hilly), explain why a space wave at VHF is problematic without line of sight, and discuss the two candidate solutions: an MF ground wave (AM) or a VHF FM transmitter placed on a high hill. State which propagation mode each uses.
\item \textbf{Transmitting aerial.} The team finally decides to broadcast locally in FM at the carrier frequency $f_c = \SI{98.0}{MHz}$. Calculate the wavelength and the length of a half-wave dipole transmitting aerial.
\item \textbf{Tuning circuit.} A listener's receiver uses a parallel $LC$ circuit with $L = \SI{0.50}{\micro H}$. Using $f_0 = \dfrac{1}{2\pi\sqrt{LC}}$, calculate the capacitance needed to tune to \SI{88}{MHz} and to \SI{108}{MHz}, and hence state the range of the variable capacitor. Sketch the resonance curve and indicate on it the tuning at \SI{98.0}{MHz}.
\item \textbf{Choice of modulation.} For the local broadcast, choose between AM and FM. Give at least two advantages of your choice and one disadvantage, referring to sound quality, noise immunity, bandwidth and cost of receivers.
\item \textbf{Bandwidth.} Suppose instead that the station had used AM with a highest audio frequency of \SI{5.0}{kHz}. 
\begin{enumerate}
\item Calculate the bandwidth of the AM signal.
\item How many such stations can fit, without overlap, in the \SI{200}{kHz} allocation?
\end{enumerate}
\item \textbf{Relay to Yaoundé.} Explain why neither a ground wave nor a direct space wave is suitable for the \SI{280}{km} Bamenda--Yaoundé link. Propose two realistic solutions (a chain of base stations / microwave repeaters on hilltops, or a satellite link) and give one advantage and one disadvantage of each.
\item \textbf{Engineering brief.} Assemble your results into a one-page brief including a labelled block diagram of the complete system: microphone $\rightarrow$ audio amplifier $\rightarrow$ modulator $\rightarrow$ carrier oscillator $\rightarrow$ power amplifier $\rightarrow$ transmitting aerial $\rightarrow$ channel $\rightarrow$ receiving aerial $\rightarrow$ tuning circuit $\rightarrow$ demodulator $\rightarrow$ loudspeaker, plus the relay link.
\end{enumerate}

\courssub{Model answer}

\textbf{1. Propagation modes for the local link.}\par
At VHF ($\sim\SI{100}{MHz}$), radio waves behave as a \cle{space wave}: they travel essentially in straight lines and are blocked by hills. Since many villages have no line of sight to Bamenda, a VHF transmitter in the town itself would leave shadow zones. Two solutions exist:
\begin{itemize}
\item an \textbf{MF ground wave} (around \SI{1}{MHz}): the wave follows the curvature of the Earth and diffracts around obstacles, covering \SI{80}{km} of hilly terrain reliably, but only AM is available there and the audio quality is limited;
\item a \textbf{VHF FM space wave} with the transmitting aerial on a high hill (e.g.\ the Bamenda escarpment): the height of the aerial extends the line-of-sight horizon to several tens of kilometres, so most villages within \SI{80}{km} can be reached, with the superior quality of FM.
\end{itemize}
The team chooses the second solution: \textbf{FM at \SI{98.0}{MHz}, space-wave propagation, hill-top site}.

\textbf{2. Half-wave dipole.}\par
\[
\lambda = \frac{c}{f_c} = \frac{3.00\times 10^{8}}{98.0\times 10^{6}} = \SI{3.06}{m}.
\]
The dipole length is
\[
\ell = \frac{\lambda}{2} = \frac{3.06}{2} \approx \SI{1.53}{m},
\]
i.e.\ two rods of about \SI{0.77}{m} each, mounted horizontally for horizontal polarisation.

\textbf{3. Tuning circuit.}\par
From $f_0 = \dfrac{1}{2\pi\sqrt{LC}}$ we get $C = \dfrac{1}{4\pi^{2} L f_0^{2}}$.
\[
C_{88} = \frac{1}{4\pi^{2}(0.50\times10^{-6})(88\times10^{6})^{2}} \approx \SI{6.5}{pF},
\qquad
C_{108} = \frac{1}{4\pi^{2}(0.50\times10^{-6})(108\times10^{6})^{2}} \approx \SI{4.3}{pF}.
\]
The variable capacitor must cover approximately $\boxed{\SI{4.3}{pF} \text{ to } \SI{6.5}{pF}}$. On the resonance curve (current versus frequency, peaked at $f_0$), tuning to \SI{98.0}{MHz} corresponds to $C \approx \SI{5.3}{pF}$, placing the peak of the curve on the desired carrier so that Radio Menchum is selected and neighbouring stations are rejected.

\textbf{4. Choice of modulation: FM.}\par
\begin{itemize}
\item \emph{Advantages of FM:} much better immunity to noise and electrical interference (lightning, motorbikes, generators), giving clear sound; better audio quality with a wider audio bandwidth, suitable for music.
\item \emph{Disadvantage:} FM requires a much larger bandwidth per station and more complex, slightly more expensive receivers; also its space-wave propagation limits coverage to (quasi) line of sight, which is why the hill-top site is essential.
\end{itemize}
AM would offer simpler, cheaper receivers and ground-wave coverage, but the noisy tropical environment (thunderstorms) would severely degrade reception. \textbf{FM is retained.}

\textbf{5. AM bandwidth (hypothetical AM option).}\par
\begin{enumerate}
\item An AM signal occupies two sidebands: $B = 2 f_{m,\max} = 2 \times 5.0 = \SI{10}{kHz}$.
\item Number of stations in \SI{200}{kHz}: $N = \dfrac{200}{10} = \boxed{20}$ stations (in practice slightly fewer, to leave guard bands).
\end{enumerate}

\textbf{6. Relay to Yaoundé.}\par
At \SI{280}{km} over mountains, the ground wave is far too attenuated and there is no line of sight for a space wave; sky-wave reflection is unreliable for a daily high-quality programme feed. Two solutions:
\begin{itemize}
\item \textbf{Chain of base stations / microwave repeaters} on hilltops every $30$--$\SI{50}{km}$: each base station receives, amplifies and retransmits the signal. \emph{Advantage:} fully under local control, moderate running cost. \emph{Disadvantage:} several sites to build, power and maintain in difficult terrain.
\item \textbf{Satellite link:} the programme is sent up to a geostationary satellite (uplink) and received in Yaoundé by a dish (downlink). \emph{Advantage:} a single hop covers the whole distance with excellent quality. \emph{Disadvantage:} high cost of satellite capacity and of the uplink equipment.
\end{itemize}
The team proposes the satellite link for reliability, with a single hill-top microwave repeater as a backup.

\textbf{7. Engineering brief (summary diagram).}

\begin{popfigure}
\begin{center}
\begin{tikzpicture}[font=\small, >=latex]
\tikzset{blk/.style={draw=popblue, thick, rounded corners, fill=popblueL, minimum height=0.8cm, align=center, inner sep=2pt}}
\node[blk] (mic) at (0,0) {Mic};
\node[blk, align=center] (mod) at (2.2,0){FM\\modulator};
\node[blk, align=center] (pa) at (4.6,0){Power\\amp};
\node[blk, align=center] (tx) at (6.8,0){Dipole\\1.53 m};
\node[blk, align=center] (rx) at (9.2,0){Receiving\\aerial};
\node[blk] (tun) at (11.4,0) {LC tuner};
\node[blk] (dem) at (13.4,0) {Demod.};
\draw[->, popdark, thick] (mic) -- (mod);
\draw[->, popdark, thick] (mod) -- (pa);
\draw[->, popdark, thick] (pa) -- (tx);
\draw[->, poporange, thick, dashed] (tx) -- node[above, text=popdark]{space wave 80 km} (rx);
\draw[->, popdark, thick] (rx) -- (tun);
\draw[->, popdark, thick] (tun) -- (dem);
\node[blk, fill=popgreenL, draw=popgreen] (sat) at (4.6,-1.8) {Satellite / repeater relay};
\node[blk, fill=popgreenL, draw=popgreen] (yde) at (9.2,-1.8) {Yaound\'e station};
\draw[->, popgreen, thick] (pa.south) |- (sat.west);
\draw[->, popgreen, thick] (sat) -- node[above, text=popdark]{280 km} (yde);
\end{tikzpicture}
\end{center}
\legende{Block diagram of the Radio Menchum system: FM broadcast chain and relay to Yaoundé.}
\end{popfigure}

\begin{attention}[Marking tips]
Every numerical answer must carry its unit and a justification; a choice (band, modulation, relay) earns full marks only when defended by at least one physical argument drawn from the course (propagation mode, noise immunity, bandwidth, line of sight). A diagram without labels, or calculations without a concluding sentence, loses marks.
\end{attention}

\newpage

\coursec{Methods and worked examples}

\coursec{Communication}

\courssub{Representing Information and Radio Wave Propagation}

\begin{definition}[Analogue and Digital Signals]
\textbf{Analogue Signal:} A continuous signal in both time and amplitude. It can take any value within a range (e.g., voice, temperature, natural light intensity).
\textbf{Digital Signal:} A discrete signal in both time and amplitude. It is represented by a sequence of binary digits (bits), typically two voltage levels representing 0 and 1.
\end{definition}

\begin{aretenir}[Comparison: Analogue vs Digital Representation]
\begin{center}
\begin{tabularx}{\linewidth}{|l|X|X|}
\hline
\rowcolor{popblueL}
\textbf{Characteristic} & \textbf{Analogue} & \textbf{Digital} \\ \hline
\textbf{Resolution} & Infinite (continuous) & Finite (quantisation steps) \\ \hline
\textbf{Noise Immunity} & Poor (noise adds directly) & High (regeneration, error correction) \\ \hline
\textbf{Processing} & Simple (amplifiers, filters) & Complex (requires ADC/DAC, DSP) \\ \hline
\textbf{Storage/Transmission} & Degrades with copies/distance & Perfect copies, encryption possible \\ \hline
\textbf{Bandwidth Efficiency} & High (baseband) & Lower (requires higher BW for same info) \\ \hline
\end{tabularx}
\end{center}
\end{aretenir}

\begin{propriete}[Radio Wave Propagation Modes]
The propagation mode depends primarily on \textbf{frequency} and \textbf{distance}.
\begin{enumerate}
\item \textbf{Ground (Surface) Wave:} \textbf{LF/MF ($< 3\ \mathrm{MHz}$).} Wave follows Earth's curvature via diffraction. Attenuation increases with frequency and ground conductivity (sea $>$ land). Range $\propto$ Transmitter Power. \emph{Uses: AM Radio, Maritime/Naval comms, Time signals.}
\item \textbf{Sky Wave (Ionospheric):} \textbf{HF ($3\text{--}30\ \mathrm{MHz}$).} Wave refracted (bent) back to Earth by ionised layers (D, E, F1, F2). Enables \textbf{intercontinental} communication.
\begin{itemize}
\item \textbf{Day:} D-layer absorbs lower HF; F-layers refract higher HF.
\item \textbf{Night:} D-layer vanishes; lower HF refract well off E/F layers $\to$ longer range but fading/multipath.
\item \textbf{Critical Frequency ($f_c$):} Highest frequency refracted vertically. $f_c = 9\sqrt{N_{\max}}$ ($N_{\max}$ = max electron density $\mathrm{m}^{-3}$).
\item \textbf{Maximum Usable Frequency (MUF):} $\mathrm{MUF} = f_c \sec \phi_i$ ($\phi_i$ = angle of incidence).
\item \textbf{Skip Zone:} Region between ground wave limit and first sky wave return $\to$ no reception.
\end{itemize}
\item \textbf{Space Wave (Line-of-Sight - LOS):} \textbf{VHF/UHF/SHF ($> 30\ \mathrm{MHz}$).} Travels in straight lines. Range limited by radio horizon.
\begin{itemize}
\item \textbf{Radio Horizon Distance:} $d_{\mathrm{horizon}} \approx \sqrt{2 k R_E h} \approx 3.57\ \mathrm{km} \sqrt{h(\mathrm{m})}$ (standard $k=1$) or $\approx 4.12\ \mathrm{km} \sqrt{h(\mathrm{m})}$ (effective Earth radius factor $k=4/3$ accounting for standard atmospheric refraction). GCE typically uses \textbf{$3.57\ \mathrm{km} \sqrt{h}$}.
\item \textbf{Max LOS Range:} $d_{\max} \approx 3.57\ \mathrm{km} \left( \sqrt{h_T} + \sqrt{h_R} \right)$.
\item \emph{Uses:} FM/TV Broadcast, Mobile Phones, Satellite, Radar, Microwave Links.
\end{itemize}
\end{enumerate}
\end{propriete}

\begin{methode}[Selecting Propagation Mode for a Link]
\begin{enumerate}
\item \textbf{Identify Frequency Band:} LF/MF $\to$ Ground; HF $\to$ Sky; VHF+ $\to$ Space.
\item \textbf{Calculate LOS Horizon:} $d_{\max} = 3.57(\sqrt{h_T} + \sqrt{h_R})\ \mathrm{km}$ ($h$ in m).
\item \textbf{Check Distance vs Horizon:} If $d > d_{\max}$ at VHF+, link impossible without relay/satellite.
\item \textbf{Assess Ionospheric Conditions (HF):} Time of day, solar cycle, MUF vs Operating Frequency.
\item \textbf{Ground Wave Range (MF):} Estimate via empirical curves (Millington) or $E \propto P_t / d^2$ (simplified); conductivity crucial.
\end{enumerate}
\end{methode}

\begin{exoresolu}[Radio Horizon and Propagation Mode Selection]
A transmitting aerial is mounted on a mast $150\ \mathrm{m}$ high. The receiving aerial is $20\ \mathrm{m}$ above ground. The carrier frequency is $100\ \mathrm{MHz}$.
\begin{enumerate}
\item Calculate the maximum line-of-sight range for space wave propagation.
\item Explain why ground wave propagation is ineffective at this frequency.
\item If the frequency were changed to $1\ \mathrm{MHz}$ (Medium Wave), which propagation mode would dominate for a distance of $200\ \mathrm{km}$? Justify.
\end{enumerate}
\tcblower
\textbf{Detailed Solution}
\begin{enumerate}
\item \textbf{Space Wave Range (Line-of-Sight):}
The maximum distance $d_{\max}$ is the sum of the horizon distances of transmitter ($h_T$) and receiver ($h_R$).
\[
d_{\max} \approx 3.57 \left( \sqrt{h_T} + \sqrt{h_R} \right) \quad \text{with } h \text{ in metres, } d \text{ in km}.
\]
$h_T = 150\ \mathrm{m}$, $h_R = 20\ \mathrm{m}$.
\[
\sqrt{150} \approx 12.25,\quad \sqrt{20} \approx 4.47
\]
\[
d_{\max} \approx 3.57 \times (12.25 + 4.47) = 3.57 \times 16.72 \approx 59.7\ \mathrm{km}.
\]
\textbf{Answer:} Maximum LOS range $\approx 60\ \mathrm{km}$.

\item \textbf{Ground Wave at 100 MHz:}
Ground wave attenuation increases sharply with frequency due to dielectric losses in the Earth and diffraction losses. At VHF ($100\ \mathrm{MHz}$), the ground wave is heavily attenuated within a few kilometres. It is negligible beyond the horizon. Space wave dominates at VHF/UHF.

\item \textbf{Propagation at 1 MHz (MF) over 200 km:}
At $1\ \mathrm{MHz}$ (Medium Frequency), ground wave propagation is effective (range up to few hundred km over land/sea depending on power/conductivity). However, for $200\ \mathrm{km}$:
\begin{itemize}
\item \textbf{Daytime:} D-layer absorbs sky wave; \textbf{Ground Wave} is the reliable mode.
\item \textbf{Nighttime:} D-layer disappears; Sky wave (refraction from E/F layers) becomes very strong and often dominates, but suffers from fading/multipath interference.
\end{itemize}
\textbf{Conclusion:} For a stable $200\ \mathrm{km}$ link, \textbf{Ground Wave} is the primary intended mode (especially daytime), though Sky Wave may interfere at night. Space wave is impossible (horizon $\approx 60\ \mathrm{km}$).
\end{enumerate}
\end{exoresolu}

\courssub{Aerials and the Tuning Circuit}

\begin{definition}[Aerial (Antenna) Parameters]
\begin{itemize}
\item \textbf{Radiation Pattern:} 3D plot of radiated power vs direction. \emph{Isotropic} (theoretical, equal all directions), \emph{Omni-directional} (e.g., vertical monopole), \emph{Directional} (e.g., Yagi, Dish).
\item \textbf{Gain ($G$):} Ratio of power density in preferred direction to that of an isotropic radiator with same input power. Unit: \textbf{dBi} (ref: isotropic) or \textbf{dBd} (ref: half-wave dipole, $0\ \mathrm{dBd} = 2.15\ \mathrm{dBi}$).
\item \textbf{Effective Length ($h_{\mathrm{eff}}$):} $V_{\mathrm{oc}} = E \cdot h_{\mathrm{eff}}$ (Open-circuit voltage induced by field strength $E$).
\item \textbf{Aperture ($A_e$):} Effective capture area. $A_e = \dfrac{G \lambda^2}{4\pi}$.
\item \textbf{Polarisation:} Orientation of $\vec{E}$-field (Vertical, Horizontal, Circular). Tx and Rx must match for max signal (cross-pol loss $\to \infty$ ideally).
\item \textbf{Common Types:} \emph{Half-wave Dipole} ($L=\lambda/2$, $G\approx 2.15\ \mathrm{dBi}$), \emph{Quarter-wave Monopole} (needs ground plane), \emph{Yagi-Uda} (directional, high gain), \emph{Parabolic Dish} (very high gain, SHF/Satellite), \emph{Loop Aerial} (direction finding).
\end{itemize}
\end{definition}

\begin{propriete}[Parallel LC Tuning Circuit (Resonance)]
Used in RF amplifiers and receivers to select desired frequency $f_0$.
\begin{itemize}
\item \textbf{Circuit:} Inductor $L$ (with series loss resistance $r$) in parallel with Capacitor $C$ (assumed ideal).
\item \textbf{Resonance Frequency:} $f_0 = \dfrac{1}{2\pi\sqrt{LC}}$. \emph{Derivation (Examinable):} Admittance $Y = \dfrac{1}{r+j\omega L} + j\omega C$. Set $\Im(Y)=0 \implies \omega_0 C = \dfrac{\omega_0 L}{r^2+\omega_0^2 L^2} \approx \dfrac{1}{\omega_0 L}$ (for $Q\gg 1$).
\item \textbf{Dynamic Resistance ($R_{\mathrm{dyn}}$):} Impedance at resonance (purely resistive).
\[
R_{\mathrm{dyn}} = \dfrac{L}{rC} = Q \omega_0 L = \dfrac{Q}{\omega_0 C}
\]
\item \textbf{Q-Factor (Quality Factor):}
\[
Q = \dfrac{\omega_0 L}{r} = \dfrac{1}{\omega_0 C r} = \dfrac{f_0}{\Delta f}
\]
where $\Delta f$ is the $3\ \mathrm{dB}$ bandwidth (frequency difference between half-power points).
\item \textbf{Resonance Curve:} Plot of $|Z|$ or $I$ vs $f$. Peak at $f_0$, width $\Delta f$ at $1/\sqrt{2}$ of peak. High $Q \to$ Sharp peak $\to$ High Selectivity.
\item \textbf{Selectivity:} Ability to reject adjacent channels. $\frac{V_{\mathrm{out}}(f_0 \pm \delta f)}{V_{\mathrm{out}}(f_0)} \approx \dfrac{1}{\sqrt{1 + Q^2 \left( \frac{2\delta f}{f_0} \right)^2}}$ (for $\delta f \ll f_0$).
\end{itemize}
\end{propriete}

\begin{methode}[Tuning Circuit Analysis: Resonance, Q-Factor and Selectivity]
\begin{enumerate}
\item \textbf{Identify Configuration:} Parallel LC (practical: $L$ with series $r$, $C$ ideal). At resonance, $Z_{\max} = R_{\mathrm{dyn}}$.
\item \textbf{Calculate $f_0$:} $f_0 = 1/(2\pi\sqrt{LC})$.
\item \textbf{Calculate $Q$:} $Q = \omega_0 L / r$ (series $r$ given) or $Q = R_p / (\omega_0 L)$ (parallel $R_p$ given). \textbf{Check which resistance is given!}
\item \textbf{Find Bandwidth:} $\Delta f = f_0 / Q$.
\item \textbf{Assess Selectivity:} Compute attenuation at adjacent channel offset $\delta f$ using approximation formula or exact impedance formula.
\end{enumerate}
\end{methode}

\begin{exoresolu}[Tuning Circuit Design and Selectivity]
A radio receiver uses a parallel LC tuning circuit. The inductor has inductance $L = 200\ \mu\mathrm{H}$ and series resistance $r = 10\ \Omega$. The capacitor is variable.
\begin{enumerate}
\item Calculate the capacitance required to tune to a station at $95\ \mathrm{MHz}$ (FM band).
\item Determine the Q-factor of the circuit at this frequency.
\item Calculate the $3\ \mathrm{dB}$ bandwidth $\Delta f$.
\item If the adjacent channel is $200\ \mathrm{kHz}$ away (standard FM spacing), will this circuit provide adequate selectivity? (Assume adequate if adjacent signal is attenuated by $> 20\ \mathrm{dB}$ relative to resonance).
\end{enumerate}
\tcblower
\textbf{Detailed Solution}
\begin{enumerate}
\item \textbf{Resonance Capacitance:}
$f_0 = 95\ \mathrm{MHz} = 95 \times 10^6\ \mathrm{Hz}$.
\[
f_0 = \frac{1}{2\pi\sqrt{LC}} \implies C = \frac{1}{(2\pi f_0)^2 L}
\]
\[
(2\pi f_0)^2 = (2\pi \times 95 \times 10^6)^2 \approx (5.969 \times 10^8)^2 \approx 3.563 \times 10^{17}\ \mathrm{s^{-2}}
\]
\[
C = \frac{1}{3.563 \times 10^{17} \times 200 \times 10^{-6}} = \frac{1}{7.126 \times 10^{13}} \approx 1.40 \times 10^{-14}\ \mathrm{F} = 14.0\ \mathrm{pF}.
\]

\item \textbf{Q-Factor:}
$\omega_0 = 2\pi f_0 = 2\pi \times 95 \times 10^6 \approx 5.97 \times 10^8\ \mathrm{rad/s}$.
\[
Q = \frac{\omega_0 L}{r} = \frac{5.97 \times 10^8 \times 200 \times 10^{-6}}{10} = \frac{1.194 \times 10^5}{10} = 11940.
\]
\textit{Note: This Q is extremely high for a discrete component at VHF; typical values are lower due to stray capacitance/skin effect. But we use given data.}

\item \textbf{Bandwidth:}
\[
\Delta f = \frac{f_0}{Q} = \frac{95 \times 10^6}{11940} \approx 7950\ \mathrm{Hz} \approx 8.0\ \mathrm{kHz}.
\]

\item \textbf{Selectivity Check:}
Channel spacing = $200\ \mathrm{kHz}$. Offset from resonance $f_{\mathrm{adj}} = f_0 \pm 200\ \mathrm{kHz}$.
Relative frequency deviation: $\frac{\Delta f_{\mathrm{offset}}}{f_0} = \frac{200 \times 10^3}{95 \times 10^6} \approx 2.105 \times 10^{-3}$.
For a parallel resonant circuit, the impedance ratio (voltage gain ratio) at offset $\delta f$ is approx:
\[
\frac{V_{\mathrm{out}}(\delta f)}{V_{\mathrm{out}}(f_0)} \approx \frac{1}{\sqrt{1 + Q^2 \left( \frac{2\delta f}{f_0} \right)^2}}
\]
(Using approximation for small $\delta f$: $\frac{\omega L - 1/\omega C}{r} \approx Q \frac{2\delta f}{f_0}$).
\[
Q \frac{2\delta f}{f_0} = 11940 \times 2 \times 2.105 \times 10^{-3} \approx 50.3
\]
Attenuation $\approx 20 \log_{10}(50.3) \approx 34\ \mathrm{dB}$.
Since $34\ \mathrm{dB} > 20\ \mathrm{dB}$, the selectivity is \textbf{more than adequate}.
\textit{Comment:} In practice, FM receivers use IF filters (ceramic/SAW) for channel selection, not just the RF tuning circuit, because tracking a high-Q circuit over the band is difficult.
\end{enumerate}
\end{exoresolu}

\courssub{Modulation, Demodulation, AM and FM}

\begin{definition}[Modulation]
Process of varying a parameter of a high-frequency \textbf{carrier wave} ($c(t) = E_c \sin \omega_c t$) in accordance with a low-frequency \textbf{modulating signal} ($m(t) = E_m \sin \omega_m t$).
\begin{itemize}
\item \textbf{Amplitude Modulation (AM):} Carrier amplitude varies. $v_{\mathrm{AM}}(t) = E_c [1 + m \sin \omega_m t] \sin \omega_c t$.
\item \textbf{Frequency Modulation (FM):} Carrier frequency varies. $v_{\mathrm{FM}}(t) = E_c \sin [\omega_c t + \beta \sin \omega_m t]$.
\item \textbf{Phase Modulation (PM):} Carrier phase varies. (Related to FM).
\end{itemize}
\textbf{Demodulation:} Reverse process: recovering $m(t)$ from modulated wave.
\end{definition}

\begin{propriete}[Amplitude Modulation (AM) -- Parameters, Spectrum, Power]
\begin{itemize}
\item \textbf{Modulation Index (Depth):} $m = \dfrac{E_m}{E_c} = \dfrac{V_{\max} - V_{\min}}{V_{\max} + V_{\min}}$. Condition for no envelope distortion: $0 \le m \le 1$. Over-modulation ($m>1$) causes phase reversal/distortion.
\item \textbf{AM Wave Equation (Single Tone):}
\[
v(t) = E_c \sin \omega_c t + \frac{mE_c}{2} \cos(\omega_c - \omega_m)t - \frac{mE_c}{2} \cos(\omega_c + \omega_m)t
\]
\item \textbf{Frequency Spectrum (DSB-FC -- Double Sideband Full Carrier):}
\begin{center}
\begin{tabular}{|c|c|c|}
\hline
\textbf{Component} & \textbf{Frequency} & \textbf{Amplitude (Peak)} \\ \hline
Carrier & $f_c$ & $E_c$ \\ \hline
Lower Sideband (LSB) & $f_c - f_m$ & $mE_c/2$ \\ \hline
Upper Sideband (USB) & $f_c + f_m$ & $mE_c/2$ \\ \hline
\end{tabular}
\end{center}
\item \textbf{Bandwidth:} $BW = 2 f_m(\max)$.
\item \textbf{Power (Load $R$):}
$P_c = \dfrac{E_c^2}{2R}$ (Carrier Power).
$P_{\mathrm{SB}} = \dfrac{(mE_c/2)^2}{2R} = \dfrac{m^2}{4} P_c$ (Power in \textbf{each} sideband).
$P_{\mathrm{total}} = P_c + 2P_{\mathrm{SB}} = P_c \left(1 + \dfrac{m^2}{2}\right)$.
\item \textbf{Efficiency:} $\eta = \dfrac{P_{\mathrm{sidebands}}}{P_{\mathrm{total}}} = \dfrac{m^2/2}{1 + m^2/2}$. Max $\eta = 33.3\%$ at $m=1$.
\item \textbf{Variants:} DSB-SC (Suppressed Carrier, $50\%$ efficiency, needs coherent detection), SSB (Single Sideband, $100\%$ info in half BW, complex).
\end{itemize}
\end{propriete}

\begin{exoresolu}[AM Power and Bandwidth Calculations]
An AM transmitter radiates $10\ \mathrm{kW}$ when the carrier is unmodulated. The carrier frequency is $1\ \mathrm{MHz}$. The modulating signal is a single tone at $5\ \mathrm{kHz}$.
\begin{enumerate}
\item Determine the carrier power $P_c$.
\item If the modulation depth is $80\%$, calculate the total transmitted power.
\item Calculate the power in each sideband.
\item Determine the transmission bandwidth.
\item If the transmitter is limited to $15\ \mathrm{kW}$ total power, what is the maximum modulation index possible?
\end{enumerate}
\tcblower
\textbf{Detailed Solution}
\begin{enumerate}
\item \textbf{Carrier Power:} Unmodulated power = Carrier power $P_c = 10\ \mathrm{kW}$.

\item \textbf{Total Power at $m=0.8$:}
$P_T = P_c \left(1 + \frac{m^2}{2}\right) = 10 \left(1 + \frac{0.8^2}{2}\right) = 10 \left(1 + \frac{0.64}{2}\right) = 10 (1 + 0.32) = 13.2\ \mathrm{kW}$.

\item \textbf{Sideband Power (each):}
$P_{\mathrm{SB}} = \frac{m^2}{4} P_c = \frac{0.64}{4} \times 10 = 0.16 \times 10 = 1.6\ \mathrm{kW}$.
Check: Total SB power = $2 \times 1.6 = 3.2\ \mathrm{kW}$. $P_T = 10 + 3.2 = 13.2\ \mathrm{kW}$. Correct.

\item \textbf{Bandwidth:}
$BW = 2 f_m = 2 \times 5\ \mathrm{kHz} = 10\ \mathrm{kHz}$.
Spectrum occupies $995\ \mathrm{kHz}$ to $1005\ \mathrm{kHz}$.

\item \textbf{Maximum Modulation Index for $P_{T,\max} = 15\ \mathrm{kW}$:}
$15 = 10 \left(1 + \frac{m_{\max}^2}{2}\right) \implies 1.5 = 1 + \frac{m_{\max}^2}{2} \implies 0.5 = \frac{m_{\max}^2}{2} \implies m_{\max}^2 = 1 \implies m_{\max} = 1$ ($100\%$).
The transmitter can handle full modulation without exceeding $15\ \mathrm{kW}$.
\end{enumerate}
\end{exoresolu}

\begin{propriete}[Frequency Modulation (FM) -- Parameters, Spectrum, Bandwidth]
\begin{itemize}
\item \textbf{Instantaneous Frequency:} $f_i(t) = f_c + \Delta f \sin \omega_m t$. $\Delta f$ = Peak Frequency Deviation (Hz).
\item \textbf{Modulation Index:} $\beta = \dfrac{\Delta f}{f_m}$. $\beta$ can be $>1$ (unlike AM).
\item \textbf{FM Wave (Single Tone):} $v(t) = E_c \sin(\omega_c t + \beta \sin \omega_m t)$.
\item \textbf{Spectrum:} Infinite sidebands at $f_c \pm n f_m$ ($n=0,1,2...$). Amplitudes $E_c J_n(\beta)$ where $J_n$ are Bessel functions of first kind. $J_0(\beta)$ = Carrier amplitude factor.
\item \textbf{Carson's Rule (Approximate Bandwidth for $>98\%$ power):}
\[
BW \approx 2(\Delta f + f_m) = 2f_m(\beta + 1)
\]
\item \textbf{Narrowband FM (NBFM):} $\beta \ll 1$ (typically $<0.3$). $BW \approx 2f_m$ (similar to AM). Used in 2-way radio (walkie-talkies, $\Delta f = 2.5\ \mathrm{kHz}$ or $5\ \mathrm{kHz}$).
\item \textbf{Wideband FM (WBFM):} $\beta > 1$. Broadcast FM: $\Delta f = 75\ \mathrm{kHz}$, $f_m(\max)=15\ \mathrm{kHz} \Rightarrow \beta=5$, $BW \approx 180\ \mathrm{kHz}$. High fidelity, noise immunity.
\item \textbf{Pre-emphasis / De-emphasis:} High frequencies boosted at Tx (Pre-emphasis, time constant $\tau = 50\ \mu\mathrm{s}$ or $75\ \mu\mathrm{s}$) and cut at Rx (De-emphasis) to improve SNR for high audio frequencies (triangular noise spectrum in FM).
\end{itemize}
\end{propriete}

\begin{exoresolu}[FM Bandwidth and Modulation Index Analysis]
An FM broadcast transmitter operates at $98\ \mathrm{MHz}$ with a maximum frequency deviation $\Delta f = 75\ \mathrm{kHz}$. The audio modulating signal ranges from $50\ \mathrm{Hz}$ to $15\ \mathrm{kHz}$.
\begin{enumerate}
\item Calculate the modulation index $\beta$ for the maximum modulating frequency ($15\ \mathrm{kHz}$) and for the minimum ($50\ \mathrm{Hz}$).
\item Estimate the transmission bandwidth using Carson's Rule.
\item The transmitter is received by a receiver with an IF bandwidth of $200\ \mathrm{kHz}$. Is this sufficient? Explain.
\item Compare the bandwidth required for this FM signal with an AM signal carrying the same audio bandwidth ($15\ \mathrm{kHz}$).
\end{enumerate}
\tcblower
\textbf{Detailed Solution}
\begin{enumerate}
\item \textbf{Modulation Index $\beta$:}
$\beta = \dfrac{\Delta f}{f_m}$. $\Delta f = 75\ \mathrm{kHz}$ (constant for a given transmitter drive level).
\begin{itemize}
\item At $f_m = 15\ \mathrm{kHz}$: $\beta_{\max} = \dfrac{75}{15} = \mathbf{5}$.
\item At $f_m = 50\ \mathrm{Hz}$: $\beta_{\min} = \dfrac{75000}{50} = \mathbf{1500}$.
\end{itemize}
\textit{Note:} $\beta$ varies inversely with $f_m$. Low frequencies produce very high $\beta$, generating many significant sidebands.

\item \textbf{Carson's Rule Bandwidth:}
Use maximum modulating frequency $f_m(\max) = 15\ \mathrm{kHz}$ for bandwidth estimation.
$BW \approx 2(\Delta f + f_m) = 2(75 + 15) = 180\ \mathrm{kHz}$.
(Using Bessel table for $\beta=5$: Significant sidebands up to $n \approx \beta+1 = 6$. $BW = 2 \times 6 \times 15 = 180\ \mathrm{kHz}$. Matches perfectly).

\item \textbf{Receiver IF Bandwidth Check:}
Required BW $\approx 180\ \mathrm{kHz}$. Receiver IF BW = $200\ \mathrm{kHz}$.
\textbf{Yes, it is sufficient.} The IF bandwidth is slightly wider than Carson's bandwidth, allowing passage of $\approx 98\%$ of signal power without excessive distortion, while rejecting adjacent channels (spaced $200\ \mathrm{kHz}$ apart in many regions).

\item \textbf{Comparison with AM:}
AM Bandwidth for $f_m(\max) = 15\ \mathrm{kHz}$: $BW_{\mathrm{AM}} = 2 \times 15 = 30\ \mathrm{kHz}$.
FM Bandwidth (Carson) = $180\ \mathrm{kHz}$.
\textbf{FM requires 6 times the bandwidth of AM} for the same baseband signal. This is the trade-off for FM's superior noise performance and fidelity.
\end{enumerate}
\end{exoresolu}

\begin{aretenir}[Comparison: AM vs FM (GCE Exam Standard)]
\begin{center}
\begin{tabularx}{\linewidth}{|l|X|X|}
\hline
\rowcolor{popblueL}
\textbf{Feature} & \textbf{AM (DSB-FC)} & \textbf{FM (WBFM)} \\ \hline
\textbf{Bandwidth} & Narrow ($2f_m$) & Wide ($2(\Delta f+f_m)$) \\ \hline
\textbf{Noise Immunity} & Poor (amplitude noise adds directly) & Excellent (Limiter rejects AM noise) \\ \hline
\textbf{Fidelity} & Limited by narrow BW & High (wide BW allows 15 kHz audio) \\ \hline
\textbf{Receiver Complexity} & Simple (Envelope Detector) & Complex (Limiter + Discriminator/PLL) \\ \hline
\textbf{Capture Effect} & No (strongest adds) & Yes (stronger signal captures receiver) \\ \hline
\textbf{Power Efficiency} & Low (Carrier wastes $>50\%$ power) & High (Constant amplitude, Class C amps) \\ \hline
\textbf{Propagation} & Ground/Sky wave (MF/HF) & Space wave only (VHF+) \\ \hline
\end{tabularx}
\end{center}
\end{aretenir}

\begin{methode}[Demodulation Techniques]
\begin{enumerate}
\item \textbf{AM Demodulation -- Envelope Detector:}
\begin{itemize}
\item Circuit: Diode (rectifier) + Parallel $R_L C$ load.
\item \textbf{Time Constant Condition:} $\dfrac{1}{\omega_c} \ll R_L C \ll \dfrac{1}{\omega_m(\max)}$.
\item Too small $RC$: Carrier ripple on output. Too large $RC$: Diagonal clipping (failure to follow envelope decay when $m$ high).
\item Output: $v_{\mathrm{out}}(t) \propto E_c(1 + m \sin \omega_m t) \xrightarrow{\text{DC Block}}$ Audio.
\end{itemize}
\item \textbf{FM Demodulation -- Frequency Discriminators:} Convert $\Delta f \to \Delta V \to$ Envelope Detect.
\begin{itemize}
\item \textbf{Slope Detector:} Tuned circuit offset from $f_c$. Operates on slope of resonance curve. Simple but non-linear, poor linearity.
\item \textbf{Foster-Seeley Discriminator:} Centre-tapped transformer + two diodes. Output $\propto$ frequency deviation. Symmetrical S-curve. Good linearity. \textbf{Requires Limiter before it} (sensitive to AM).
\item \textbf{Ratio Detector:} Similar topology but with tertiary winding/long $\tau$ capacitor giving inherent limiting (AM rejection). No separate limiter strictly needed. Wider capture range.
\item \textbf{Phase Locked Loop (PLL) -- Modern Standard:} VCO locks to input frequency. Control voltage $\propto$ frequency deviation $\to$ Demodulated output. Excellent linearity, noise performance, integrable (IC).
\end{itemize}
\item \textbf{S-Curve:} Plot of $V_{\mathrm{out}}$ vs $f_{\mathrm{in}}$. Linear region = usable range. Slope = Sensitivity (V/Hz). Saturation limits = Dynamic Range ($\pm \Delta f$).
\end{enumerate}
\end{methode}

\begin{exoresolu}[AM Envelope Detector Design and FM Discriminator Characteristics]
An AM signal with carrier frequency $1\ \mathrm{MHz}$ and maximum modulating frequency $5\ \mathrm{kHz}$ (modulation index $m=0.8$) is to be demodulated using a simple envelope detector. The load resistance $R_L = 10\ \mathrm{k}\Omega$.
\begin{enumerate}
\item Determine the suitable range for the capacitor $C$ to avoid both carrier ripple and diagonal clipping.
\item Select a standard capacitor value and verify.
\item For FM demodulation, sketch the ideal "S-curve" of a discriminator and label: Linear region, Sensitivity, Dynamic Range. Explain why a Limiter is essential before a Foster-Seeley discriminator but not strictly before a Ratio Detector.
\end{enumerate}
\tcblower
\textbf{Detailed Solution}
\begin{enumerate}
\item \textbf{RC Time Constant Constraints:}
$\omega_c = 2\pi f_c = 2\pi \times 10^6 \approx 6.28 \times 10^6\ \mathrm{rad/s}$.
$\omega_m(\max) = 2\pi \times 5000 \approx 3.14 \times 10^4\ \mathrm{rad/s}$.
Condition: $\dfrac{1}{\omega_c} \ll RC \ll \dfrac{1}{\omega_m}$.
Lower bound (Carrier ripple): $RC \gg \dfrac{1}{6.28 \times 10^6} \approx 0.16\ \mu\mathrm{s}$.
Upper bound (Diagonal clipping): $RC \ll \dfrac{1}{3.14 \times 10^4} \approx 31.8\ \mu\mathrm{s}$.
Safe Range: \textbf{$1\ \mu\mathrm{s}$ to $10\ \mu\mathrm{s}$} (Geometric mean $\approx 2.2\ \mu\mathrm{s}$).

\item \textbf{Component Selection:}
$R_L = 10\ \mathrm{k}\Omega$.
Choose $RC = 5\ \mu\mathrm{s}$ (mid-range).
$C = \dfrac{5 \times 10^{-6}}{10 \times 10^3} = 500\ \mathrm{pF}$.
Standard value: $470\ \mathrm{pF}$ or $560\ \mathrm{pF}$.
Check $C = 470\ \mathrm{pF}$: $RC = 4.7\ \mu\mathrm{s}$.
$\frac{1}{\omega_c} = 0.16\ \mu\mathrm{s} \ll 4.7\ \mu\mathrm{s} \ll 31.8\ \mu\mathrm{s} = \frac{1}{\omega_m}$. \textbf{Condition satisfied.}

\item \textbf{FM Discriminator S-Curve and Limiter:}
\textbf{Sketch Description:} Graph of $V_{\mathrm{out}}$ (Vertical) vs $f_{\mathrm{in}}$ (Horizontal).
\begin{itemize}
\item Shape: Inverted 'S' (or 'S' depending on polarity). Passes through origin ($f_c$, 0V).
\item \textbf{Linear Region:} Central portion where $V_{\mathrm{out}} \propto (f_{\mathrm{in}} - f_c)$. Demodulation occurs here.
\item \textbf{Sensitivity (Slope):} $S = dV/df$ (Volts/Hz or V/kHz) in linear region.
\item \textbf{Dynamic Range / Peak Deviation Limit:} Frequencies $f_c \pm \Delta f$ where curve starts to saturate/flatten. Beyond this, distortion occurs.
\end{itemize}
\textbf{Limiter Requirement:}
\begin{itemize}
\item \textbf{Foster-Seeley:} Output amplitude depends on input signal amplitude (AM component) \textbf{AND} frequency. Any amplitude noise (lightning, ignition, fading) appears directly on the audio output. \textbf{A Limiter (hard limiter/clipper) is mandatory} to remove AM before discrimination.
\item \textbf{Ratio Detector:} Circuit topology (diode charging a large capacitor $C_3$ with long time constant) creates a reference voltage proportional to the \emph{average} signal amplitude. The output is taken as a \textbf{ratio} of voltages across the load resistors. This makes the output inherently independent of input amplitude variations (AM rejection built-in). \textbf{No separate limiter stage is required}, though one is often added for extra quieting.
\end{itemize}
\end{enumerate}
\end{exoresolu}

\courssub{Digital Conversion and Communication Channels}

\begin{propriete}[Analogue-to-Digital Conversion (ADC)]
\begin{itemize}
\item \textbf{Sampling (Nyquist Theorem):} To reconstruct an analogue signal of maximum frequency $f_{\max}$ without aliasing, sampling frequency $f_s \ge 2 f_{\max}$. Practical systems use $f_s \ge 2.2 f_{\max}$ to allow for anti-aliasing filter roll-off.
\item \textbf{Quantisation:} Mapping continuous amplitude to $2^N$ discrete levels. $N$ = number of bits.
\item \textbf{Resolution (Step Size):} $\Delta V = \dfrac{V_{\mathrm{FSR}}}{2^N}$ where $V_{\mathrm{FSR}}$ is Full Scale Range.
\item \textbf{Quantisation Error:} Maximum error $= \pm \Delta V / 2$. RMS Quantisation Noise Voltage $V_{q,\mathrm{rms}} = \dfrac{\Delta V}{\sqrt{12}}$.
\item \textbf{Signal-to-Quantisation-Noise Ratio (SQNR):} For a full-scale sine wave ($V_{\mathrm{peak}} = V_{\mathrm{FSR}}/2$):
\[
\mathrm{SQNR}_{\mathrm{dB}} \approx 6.02 N + 1.76\ \mathrm{dB}
\]
Rule of thumb: \textbf{+1 bit $\approx$ +6 dB SQNR}.
\item \textbf{ADC Types (Principles):}
\begin{itemize}
\item \emph{Flash (Parallel):} Fastest, low resolution ($<8$ bit), high power, $2^N-1$ comparators.
\item \emph{Successive Approximation (SAR):} Medium speed, 12--16 bit, 1 comparator + DAC + logic. Standard general purpose.
\item \emph{Sigma-Delta ($\Sigma\Delta$):} Very high resolution (16--24+ bit), oversampling + noise shaping. Audio, instrumentation.
\item \emph{Dual Slope (Integrating):} High accuracy, slow. Digital Multimeters (DMM).
\end{itemize}
\end{itemize}
\end{propriete}

\begin{propriete}[Digital-to-Analogue Conversion (DAC)]
\begin{itemize}
\item \textbf{Reconstruction:} Digital codes $\to$ Staircase voltage (Zero-Order Hold) $\to$ \textbf{Reconstruction Filter} (Low Pass, cutoff $f_c \approx f_s/2$) $\to$ Smoothed Analogue Output.
\item \textbf{Key Specs:} Resolution ($N$ bits), Settling Time, Glitch Energy (code transitions), Monotonicity.
\item \textbf{Common Architecture:} R-2R Ladder (voltage mode) or Current Steering (high speed).
\end{itemize}
\end{propriete}

\begin{methode}[ADC Specifications and Nyquist Rate]
\begin{enumerate}
\item \textbf{Nyquist Frequency:} $f_{\max} < f_s/2$.
\item \textbf{Resolution:} $\Delta V = V_{\mathrm{FSR}} / 2^N$.
\item \textbf{SQNR:} $6.02N + 1.76$ dB (full-scale sine).
\item \textbf{Anti-Aliasing Filter:} Essential before ADC. Low-pass, cutoff $f_c \approx f_{\max}$, stopband at $f_s - f_{\max}$ with attenuation $> \mathrm{SQNR}$ to prevent aliases folding into baseband.
\item \textbf{DAC Reconstruction:} Zero-order hold $\to$ Sinc roll-off ($\sin x/x$) $\to$ Reconstruction Filter compensates.
\end{enumerate}
\end{methode}

\begin{exoresolu}[ADC Specifications and Nyquist Rate]
A digital audio system samples at $44.1\ \mathrm{kHz}$ (CD standard) using a $16$-bit ADC. The input signal range is $\pm 1\ \mathrm{V}$.
\begin{enumerate}
\item What is the maximum audio frequency that can be recorded without aliasing?
\item Calculate the voltage resolution (step size) of the ADC.
\item Determine the theoretical Signal-to-Quantisation-Noise Ratio (SQNR) in dB.
\item If the system is upgraded to $24$-bit at the same sampling rate, what is the new SQNR? By how many dB does it improve?
\item Explain why an anti-aliasing filter is essential before the ADC, and state its required cutoff frequency characteristic.
\end{enumerate}
\tcblower
\textbf{Detailed Solution}
\begin{enumerate}
\item \textbf{Nyquist Frequency:}
$f_s = 44.1\ \mathrm{kHz}$.
$f_{\max} < \dfrac{f_s}{2} = \dfrac{44.1}{2} = 22.05\ \mathrm{kHz}$.
\textbf{Answer:} $22.05\ \mathrm{kHz}$ (Theoretical limit). Practical limit $\approx 20\ \mathrm{kHz}$ due to filter roll-off.

\item \textbf{Resolution (Step Size $\Delta V$):}
$V_{\mathrm{FSR}} = 1 - (-1) = 2\ \mathrm{V}$.
$N = 16 \implies 2^N = 65536$ levels.
$\Delta V = \dfrac{2\ \mathrm{V}}{65536} \approx 30.5\ \mu\mathrm{V}$.

\item \textbf{SQNR (16-bit):}
$\mathrm{SQNR} \approx 6.02 N + 1.76 = 6.02 \times 16 + 1.76 = 96.32 + 1.76 = 98.08\ \mathrm{dB}$.
(Standard CD dynamic range $\approx 96\ \mathrm{dB}$, close to this theoretical max).

\item \textbf{SQNR (24-bit):}
$\mathrm{SQNR}_{24} \approx 6.02 \times 24 + 1.76 = 144.48 + 1.76 = 146.24\ \mathrm{dB}$.
Improvement $= 146.24 - 98.08 = 48.16\ \mathrm{dB}$.
(Equivalent to $8 \text{ bits} \times 6.02 \approx 48\ \mathrm{dB}$).

\item \textbf{Anti-Aliasing Filter:}
\textbf{Why:} Any frequency component $f > f_s/2$ in the input signal will be "mirrored" (aliased) into the baseband $0 \text{ to } f_s/2$ as a false frequency $f_{\mathrm{alias}} = |f - k f_s|$. This corrupts the digital signal irreversibly.
\textbf{Characteristic:} Low-pass filter. Cutoff $f_c \approx f_{\max}$ (e.g., $20\ \mathrm{kHz}$). Stopband must start at $f_s - f_{\max} = 24.1\ \mathrm{kHz}$ with high attenuation ($> 96\ \mathrm{dB}$ for 16-bit) to prevent aliases folding into the passband. This requires a steep filter (high order) or oversampling techniques.
\end{enumerate}
\end{exoresolu}

\begin{propriete}[Communication Channels: Bandwidth, Sidebands, Satellite, Base Stations]
\begin{itemize}
\item \textbf{Channel Bandwidth:} Range of frequencies allocated for a single communication channel. Must accommodate modulated signal BW (e.g., AM: $10\ \mathrm{kHz}$, FM: $200\ \mathrm{kHz}$, TV: $6\ \mathrm{MHz}$/$8\ \mathrm{MHz}$, GSM: $200\ \mathrm{kHz}$).
\item \textbf{Sidebands:} Frequency components generated by modulation symmetrically around $f_c$.
\begin{itemize}
\item AM: Two sidebands (USB, LSB) + Carrier (DSB-FC). BW = $2f_m$.
\item SSB (Single Sideband): Suppress carrier and one sideband. BW = $f_m$. Power efficient, complex receiver (coherent detection).
\item FM: Infinite sidebands theoretically; Carson's BW ($2(\Delta f+f_m)$) used practically.
\end{itemize}
\item \textbf{Satellite Communication (Geostationary - GEO):}
\begin{itemize}
\item Orbit: $\approx 35786\ \mathrm{km}$ above equator. Period = $24\ \mathrm{h}$. Appears fixed.
\item \textbf{Uplink ($f_u$) / Downlink ($f_d$):} Frequencies differ (e.g., C-band: $6/4\ \mathrm{GHz}$, Ku-band: $14/12\ \mathrm{GHz}$) to avoid interference (Tx leakage into Rx).
\item \textbf{Transponder Block Diagram:} Rx Antenna $\to$ LNA (Low Noise Amp) $\to$ Bandpass Filter $\to$ Mixer + LO (Frequency Conversion $f_u \to f_d$) $\to$ HPA (High Power Amp, TWTA/SSPA) $\to$ Tx Antenna.
\item \textbf{Free Space Path Loss (FSPL):}
\[
L_p = \left( \frac{4\pi d}{\lambda} \right)^2 = \left( \frac{4\pi d f}{c} \right)^2
\]
In dB ($d$ in km, $f$ in GHz): $L_p(\mathrm{dB}) = 20\log_{10}(d) + 20\log_{10}(f) + 92.45$.
\item \textbf{Link Budget:} $P_R (\mathrm{dBW}) = \mathrm{EIRP} + G_R - L_p - L_{\mathrm{atm}} - L_{\mathrm{other}}$.
$\mathrm{EIRP} = P_T + G_T$ (dBW). $G_R$ = Rx Antenna Gain (dBi).
\item \textbf{Latency:} Round trip $\approx 240\ \mathrm{ms}$ (propagation only). Affects voice/TCP.
\end{itemize}
\item \textbf{Base Station (Cellular):} Fixed transceiver serving a \textbf{Cell}. Functions: Radio Resource Management, Handover control, Connection to Mobile Switching Centre (MSC)/Core Network.
\item \textbf{Frequency Reuse:} Same frequency channels reused in geographically separated cells (Co-channel Interference CCI managed). Cluster size $K$ (e.g., 3, 4, 7, 12). Capacity $\propto$ Number of Clusters. Sectoring (120°/60°) reduces CCI $\to$ smaller $K$ $\to$ higher capacity.
\end{itemize}
\end{propriete}

\begin{exoresolu}[Satellite Link Budget and Cellular Concepts]
A Geostationary satellite operates in Ku-band. Uplink $f_u = 14\ \mathrm{GHz}$, Downlink $f_d = 12\ \mathrm{GHz}$. Satellite distance $d = 36000\ \mathrm{km}$. Earth station antenna gain $G_T = 50\ \mathrm{dBi}$, Satellite receive antenna gain $G_R = 30\ \mathrm{dBi}$. Transmitter EIRP = $60\ \mathrm{dBW}$.
\begin{enumerate}
\item Calculate the Free Space Path Loss (FSPL) for the Uplink in dB.
\item Calculate the Received Power at the Satellite ($P_R$) in dBW.
\item If the system uses Frequency Division Multiple Access (FDMA) with a total transponder bandwidth of $500\ \mathrm{MHz}$ and each channel requires $36\ \mathrm{MHz}$ (including guard bands), how many channels can be supported?
\item Explain the concept of "Frequency Reuse" in cellular networks and how it increases capacity.
\end{enumerate}
\tcblower
\textbf{Detailed Solution}
\begin{enumerate}
\item \textbf{Uplink FSPL:}
Formula: $L_p(\mathrm{dB}) = 20\log_{10}(d_{\mathrm{km}}) + 20\log_{10}(f_{\mathrm{GHz}}) + 92.45$.
$d = 36000\ \mathrm{km}$, $f = 14\ \mathrm{GHz}$.
$20\log_{10}(36000) = 20 \times 4.5563 = 91.126$.
$20\log_{10}(14) = 20 \times 1.1461 = 22.922$.
$L_p = 91.126 + 22.922 + 92.45 = 206.50\ \mathrm{dB}$.

\item \textbf{Received Power at Satellite ($P_R$):}
Link Budget: $P_R (\mathrm{dBW}) = \mathrm{EIRP} + G_R - L_p$.
$\mathrm{EIRP} = 60\ \mathrm{dBW}$ (Given as combined $P_T + G_T$).
$G_R = 30\ \mathrm{dBi}$.
$P_R = 60 + 30 - 206.5 = -116.5\ \mathrm{dBW}$.
Convert to Watts: $P_R = 10^{-11.65} \approx 2.24 \times 10^{-12}\ \mathrm{W} = 2.24\ \mathrm{pW}$.
\textit{Note:} This tiny power requires a very Low Noise Amplifier (LNA) at the satellite (Noise Figure $\approx 1\ \mathrm{dB}$).

\item \textbf{FDMA Channel Capacity:}
Total BW = $500\ \mathrm{MHz}$. Channel BW = $36\ \mathrm{MHz}$.
$N = \left\lfloor \dfrac{500}{36} \right\rfloor = \left\lfloor 13.88 \right\rfloor = \mathbf{13\ \text{channels}}$.
(Guard bands consume the remainder).

\item \textbf{Frequency Reuse in Cellular Networks:}
\textbf{Concept:} The same set of frequency channels is reused in different cells that are geographically separated far enough so that co-channel interference (CCI) is below an acceptable threshold (Carrier-to-Interference ratio C/I).
\textbf{Mechanism:} Service area divided into clusters of $K$ cells (e.g., $K=3, 4, 7, 12$). Each cell in a cluster gets a unique subset of the total channels. The cluster pattern repeats spatially.
\textbf{Capacity Increase:} Total system capacity $C = M \times N_c$ where $M$ = number of clusters, $N_c$ = channels per cluster ($N_{\mathrm{total}}/K$). Without reuse ($K=1$), $C = N_{\mathrm{total}}$. With reuse, $C = M \times (N_{\mathrm{total}}/K)$. As area grows, $M$ grows $\implies$ Capacity grows linearly with area. This is the fundamental principle allowing millions of users with limited spectrum.
\textbf{Trade-off:} Smaller $K$ (tighter reuse) $\to$ Higher capacity but higher interference $\to$ Requires sectoring (directional antennas) / power control / advanced coding.
\end{enumerate}
\end{exoresolu}

\begin{methode}[Synthesis: Complete Link Analysis -- Shannon-Hartley \& Link Budget]
\begin{enumerate}
\item \textbf{Shannon-Hartley Theorem:} Maximum theoretical channel capacity (error-free) for a bandlimited AWGN channel:
\[
C = B \log_2 \left(1 + \frac{S}{N}\right) \quad \text{bits/s}
\]
$B$ = Bandwidth (Hz), $S/N$ = Signal-to-Noise Ratio (linear power ratio).
\item \textbf{Thermal Noise Power:} $N = k T B$. $k = 1.38 \times 10^{-23}\ \mathrm{J/K}$ (Boltzmann). $T$ = Noise Temperature (K), typically $290\ \mathrm{K}$ (room temp).
$N(\mathrm{dBm}) = -174 + 10\log_{10}(B/\mathrm{Hz}) + NF$ (Noise Figure in dB).
\item \textbf{Energy per Bit:} $E_b = S / R_b$. $N_0 = kT$. $\frac{E_b}{N_0} = \frac{S}{N} \frac{B}{R_b}$.
Shannon Limit: $\frac{E_b}{N_0} > \ln 2 \approx -1.6\ \mathrm{dB}$ for $R_b \to C$.
\item \textbf{Link Budget Steps:}
1. Tx Power ($P_T$) + Tx Antenna Gain ($G_T$) = EIRP (dBW/dBm).
2. $- \text{Path Loss } (L_p) - \text{Atmospheric/Misc Losses } (L_a)$.
3. $+ \text{Rx Antenna Gain } (G_R)$.
4. $= \text{Received Power } (P_R)$.
5. Noise Power $N = k T_{\mathrm{sys}} B$ (dBm).
6. $\mathrm{SNR} = P_R - N$ (in dB).
7. Check if $\mathrm{SNR} > \mathrm{SNR}_{\mathrm{req}}$ for modulation/coding (e.g., BPSK $\approx 10\ \mathrm{dB}$ $E_b/N_0$ for BER $10^{-6}$).
\item \textbf{Digital Modulation Spectral Efficiency:} $\eta = R_b / B$ (bits/s/Hz). BPSK=1, QPSK=2, 16-QAM=4, 64-QAM=6. Higher $\eta$ requires higher SNR.
\end{enumerate}
\end{methode}

\begin{exoresolu}[Shannon Capacity and Link Budget for a Digital Radio Link]
A point-to-point microwave link operates at $10\ \mathrm{GHz}$ over $30\ \mathrm{km}$. Bandwidth $B = 20\ \mathrm{MHz}$. Tx Power $P_T = 1\ \mathrm{W}$ ($30\ \mathrm{dBm}$). Tx/Rx Antenna Gains = $35\ \mathrm{dBi}$ each. Receiver Noise Figure $NF = 3\ \mathrm{dB}$. System uses QPSK ($R_b = 40\ \mathrm{Mbps}$). Required $E_b/N_0 = 10\ \mathrm{dB}$ for target BER.
\begin{enumerate}
\item Calculate Free Space Path Loss (FSPL).
\item Calculate Received Power $P_R$ (dBm).
\item Calculate Receiver Noise Power $N$ (dBm).
\item Determine achieved $E_b/N_0$ (dB). Is the link viable? What is the Fade Margin?
\item What is the Shannon Capacity of this channel? Is QPSK spectral efficiency close to the limit?
\end{enumerate}
\tcblower
\textbf{Detailed Solution}
\begin{enumerate}
\item \textbf{FSPL:}
$L_p = 20\log_{10}(30) + 20\log_{10}(10) + 92.45 = 29.54 + 20 + 92.45 = 141.99\ \mathrm{dB} \approx 142\ \mathrm{dB}$.

\item \textbf{Received Power $P_R$:}
$\mathrm{EIRP} = P_T + G_T = 30 + 35 = 65\ \mathrm{dBm}$.
$P_R = \mathrm{EIRP} - L_p + G_R = 65 - 142 + 35 = -42\ \mathrm{dBm}$.

\item \textbf{Noise Power $N$:}
$N(\mathrm{dBm}) = -174 + 10\log_{10}(B) + NF$.
$10\log_{10}(20 \times 10^6) = 10 \times 7.3010 = 73.01\ \mathrm{dB}$.
$N = -174 + 73.01 + 3 = -97.99\ \mathrm{dBm} \approx -98\ \mathrm{dBm}$.

\item \textbf{Achieved $E_b/N_0$ and Viability:}
$\mathrm{SNR} = P_R - N = -42 - (-98) = 56\ \mathrm{dB}$.
$R_b = 40\ \mathrm{Mbps} \implies 10\log_{10}(R_b) = 76.02\ \mathrm{dB}$.
$10\log_{10}(B) = 73.01\ \mathrm{dB}$.
$\frac{E_b}{N_0} (\mathrm{dB}) = \mathrm{SNR} + 10\log_{10}(B/R_b) = 56 + (73.01 - 76.02) = 56 - 3.01 = 52.99\ \mathrm{dB}$.
Required $E_b/N_0 = 10\ \mathrm{dB}$.
\textbf{Fade Margin} = Achieved - Required = $53 - 10 = 43\ \mathrm{dB}$.
\textbf{Conclusion:} Link is \textbf{highly viable} with enormous fade margin ($43\ \mathrm{dB}$). Typical design targets $20\ \mathrm{dB}$--$30\ \mathrm{dB}$ margin.

\item \textbf{Shannon Capacity:}
$\mathrm{SNR}_{\mathrm{lin}} = 10^{56/10} \approx 3.98 \times 10^5$.
$C = B \log_2(1 + \mathrm{SNR}) = 20 \times 10^6 \times \log_2(3.98 \times 10^5)$.
$\log_2(3.98 \times 10^5) = \frac{\log_{10}(3.98 \times 10^5)}{\log_{10}2} \approx \frac{5.600}{0.3010} \approx 18.60$.
$C \approx 20 \times 10^6 \times 18.60 = 372\ \mathrm{Mbps}$.
\textbf{Spectral Efficiency:}
Shannon Limit $\eta_{\max} = C/B \approx 18.6\ \mathrm{bit/s/Hz}$.
QPSK Efficiency $\eta_{\mathrm{QPSK}} = 2\ \mathrm{bit/s/Hz}$.
\textbf{Analysis:} QPSK uses only $\approx 10\%$ of the theoretical capacity. Higher order modulation (e.g., 256-QAM, $\eta=8$) or coding could increase throughput significantly, but would require higher SNR (reducing fade margin). The current link is "power limited" (excess power) but "bandwidth limited" (fixed 20 MHz).
\end{enumerate}
\end{exoresolu}

\begin{savaistu}[Cameroon Context: Communication Infrastructure]
\begin{itemize}
\item \textbf{Satellite:} Cameroon uses \textbf{Intelsat}, \textbf{Eutelsat}, \textbf{RASCOM} (African satellite) for international trunking, TV distribution (CRTV), and VSAT networks in remote areas (e.g., Far North, East regions).
\item \textbf{Fibre Optic:} \textbf{CAMTEL} manages the national backbone (WACS, SAT-3 submarine cable landing points in Kribi/Limbe). Major cities (Yaoundé, Douala, Bafoussam, Garoua) linked by fibre.
\item \textbf{Mobile:} \textbf{MTN Cameroon}, \textbf{Orange Cameroon}, \textbf{Nexttel (Viettel)}, \textbf{Camtel Mobile}. Technologies: GSM (900/1800 MHz), UMTS (2100 MHz), LTE (800/1800/2600 MHz). Base stations (BTS/NodeB/eNodeB) ubiquitous.
\item \textbf{Regulation:} \textbf{ART} (Agence de Régulation des Télécommunications) manages spectrum allocation, licensing, numbering.
\item \textbf{Broadcast:} CRTV (AM/MF: 918 kHz Yaoundé; FM/VHF: 88--108 MHz nationwide; TV/UHF). Private radio/TV stations licensed by ART.
\end{itemize}
\end{savaistu}

\begin{attention}[Common Examination Traps]
\begin{itemize}
\item \textbf{AM Modulation Index:} $m = E_m/E_c$. Never exceed 1. If $V_{\max}, V_{\min}$ given on CRO: $m = (V_{\max}-V_{\min})/(V_{\max}+V_{\min})$.
\item \textbf{FM Modulation Index:} $\beta = \Delta f / f_m$. Can be $>1$. Do not confuse with AM $m$.
\item \textbf{Bandwidth:} AM = $2f_m$. FM = $2(\Delta f + f_m)$ (Carson). Don't use $2f_m$ for FM.
\item \textbf{Tuning Circuit Q:} $Q = \omega L / r$ (series $r$). If parallel $R_p$ given, $Q = R_p / (\omega L)$. Check circuit diagram!
\item \textbf{Envelope Detector RC:} $\frac{1}{\omega_c} \ll RC \ll \frac{1}{\omega_m}$. Geometric mean $\sqrt{\frac{1}{\omega_c \omega_m}}$ is a good design choice.
\item \textbf{Satellite Link Budget:} Uplink and Downlink frequencies differ $\to$ Path Loss differs. Use correct $f$ for each leg.
\item \textbf{Shannon Capacity:} $C = B \log_2(1+S/N)$. $S/N$ is \textbf{linear power ratio}, not dB. Convert dB $\to$ linear first.
\item \textbf{Units:} Always write units (Hz, W, dB, dBm, dBW, $\Omega$, F, H). $\mathrm{dBm}$ ref $1\ \mathrm{mW}$, $\mathrm{dBW}$ ref $1\ \mathrm{W}$.
\end{itemize}
\end{attention}

\newpage

\coursec{Exercises}

\courssub{I check my knowledge}

\begin{exercice}[MCQ: Radio wave propagation]
Which of the following statements about radio wave propagation is correct?
\begin{enumerate}
\item Ground waves follow the curvature of the Earth and are most effective at low frequencies.
\item Sky waves are reflected by the ionosphere and are used for very high frequency (VHF) communication.
\item Space waves travel in straight lines and are the only mode used for satellite communication.
\item The critical frequency is the highest frequency that can be reflected by the ionosphere at vertical incidence.
\end{enumerate}
\end{exercice}

\begin{exercice}[True/False with justification]
State whether each statement is true or false and justify your answer briefly.
\begin{enumerate}
\item In amplitude modulation (AM), the bandwidth required is twice the maximum modulating frequency.
\item Frequency modulation (FM) is more susceptible to noise than AM because noise affects frequency variations.
\item A half-wave dipole aerial has a length equal to one wavelength of the transmitted signal.
\item The Nyquist sampling theorem states that the sampling frequency must be at least equal to the highest frequency present in the signal.
\end{enumerate}
\end{exercice}

\begin{exercice}[Definitions]
Define the following terms as used in communication systems:
\begin{enumerate}
\item Modulation index (for AM)
\item Sidebands
\item Quantisation error
\item Geostationary satellite
\end{enumerate}
\end{exercice}

\begin{exercice}[Direct calculation: Wavelength and aerial length]
A radio station broadcasts on a carrier frequency of 95.5 MHz. Calculate:
\begin{enumerate}
\item The wavelength of the carrier wave in metres.
\item The length of a half-wave dipole aerial for this frequency.
\end{enumerate}
\end{exercice}

\courssub{I practise}

\begin{exercice}[AM modulation index and overmodulation]
An AM wave displayed on an oscilloscope has a maximum envelope amplitude of 12 V and a minimum envelope amplitude of 4 V.
\begin{enumerate}
\item Calculate the modulation index $m$.
\item State whether the signal is overmodulated. Justify your answer.
\item If the carrier amplitude is 8 V, determine the peak amplitude of each sideband.
\end{enumerate}
\end{exercice}

\begin{exercice}[Tuning circuit frequency range]
A radio receiver uses a tuning circuit consisting of a fixed inductor of inductance $L = 180\ \mu\mathrm{H}$ and a variable capacitor adjustable from $C_{\min} = 20\ \mathrm{pF}$ to $C_{\max} = 300\ \mathrm{pF}$.
\begin{enumerate}
\item Determine the minimum and maximum resonant frequencies of the circuit.
\item Identify the radio band (LW, MW, SW, VHF/FM) covered by this range.
\item Sketch the resonance curve (current vs frequency) for a given capacitance, labelling the resonant frequency, bandwidth and quality factor $Q$.
\end{enumerate}
\end{exercice}

\begin{exercice}[AM frequency spectrum and bandwidth]
An audio signal containing frequencies up to 4.5 kHz amplitude-modulates a carrier of frequency 1 MHz.
\begin{enumerate}
\item Sketch the frequency spectrum of the resulting AM wave, clearly marking the carrier and the two sidebands.
\item Calculate the channel bandwidth required.
\item If the same audio signal were used for FM with a maximum frequency deviation of 75 kHz, estimate the bandwidth using Carson's rule.
\end{enumerate}
\end{exercice}

\begin{exercice}[Digital audio parameters]
A music signal is sampled at 44.1 kHz with 16-bit quantisation.
\begin{enumerate}
\item Calculate the number of quantisation levels.
\item Determine the bit rate (in kbps) for a stereo signal (two channels).
\item What is the maximum audio frequency that can be faithfully reproduced according to the Nyquist criterion?
\item Explain why a low-pass anti-aliasing filter is needed before sampling.
\end{enumerate}
\end{exercice}

\courssub{I prepare for the GCE}

\begin{exercice}[{Propagation modes and reception behind a hill] [15 marks}]
A car radio receives both MW (medium wave, ~1 MHz) and VHF/FM (~100 MHz) broadcasts. When the car drives behind a steep hill, the FM signal is lost while the MW signal remains audible.
\begin{enumerate}
\item Name and describe the three main modes of radio wave propagation. [6 marks]
\item Using diagrams, explain why the MW signal can still be received behind the hill while the FM signal cannot. [6 marks]
\item State one advantage and one disadvantage of using sky wave propagation for long-distance communication. [3 marks]
\end{enumerate}
\end{exercice}

\begin{exercice}[{Comparison of AM and FM; TV sound transmission] [15 marks}]
\begin{enumerate}
\item Compare amplitude modulation (AM) and frequency modulation (FM) under the following four headings:
\begin{itemize}
\item Susceptibility to noise
\item Bandwidth requirement
\item Range (for a given transmitter power)
\item Complexity of transmitter and receiver circuits
\end{itemize}
[8 marks]
\item The Cameroon Radio Television (CRTV) transmits television sound using FM rather than AM. Justify this choice with reference to the comparison above. [4 marks]
\item An FM transmitter has a carrier frequency of 98.5 MHz and a maximum frequency deviation of 75 kHz. If the modulating signal has a maximum frequency of 15 kHz, calculate the approximate bandwidth using Carson's rule. [3 marks]
\end{enumerate}
\end{exercice}

\begin{exercice}[{Satellite communication delay] [12 marks}]
A geostationary satellite orbits at an altitude of 36,000 km above the Earth's surface. A telephone call is routed from Douala (Cameroon) to London (UK) via this satellite. Assume the uplink and downlink paths are each approximately equal to the satellite altitude (ignore Earth curvature for this calculation).
\begin{enumerate}
\item Calculate the minimum time delay for a signal to travel from Douala to the satellite and then to London (one-way). [4 marks]
\item If the conversation is a two-way telephone call, what is the total round-trip delay experienced by the speakers? [2 marks]
\item Comment on the effect of this delay on natural telephone conversation. [3 marks]
\item Suggest one technical method to mitigate the annoyance caused by this delay in satellite telephony. [3 marks]
\end{enumerate}
\end{exercice}

\newpage

\coursec{Solutions to the exercises}

\coursec{Corrected Exercises — Communication (PHY14)}

\begin{corrige}[Exercise 1 — MCQ: Radio wave propagation]
\textbf{Correct answers: (a) and (d).}

\begin{itemize}
\item \textbf{(a) Correct.} Ground (surface) waves travel along the Earth's surface, following its curvature by diffraction. Diffraction is significant only when the wavelength is comparable to obstacles, i.e. at \cle{low frequencies} (long wavelengths, LW and MW bands).
\item \textbf{(b) False.} Sky waves are indeed reflected (refracted) back by the ionosphere, but this works for HF (short wave, roughly $3$--$30\ \mathrm{MHz}$). VHF waves ($30$--$300\ \mathrm{MHz}$) pass through the ionosphere and are \emph{not} reflected; VHF uses the space wave.
\item \textbf{(c) False.} Space waves do travel in straight lines (line of sight) and are used for satellite links, but they are \emph{not} the only mode used for satellites in general; the word ``only'' makes the statement incorrect.
\item \textbf{(d) Correct.} The \cle{critical frequency} is by definition the highest frequency reflected back to Earth when the wave is sent \emph{vertically} (at vertical incidence) towards the ionosphere. Above it, the wave penetrates the layer and escapes into space.
\end{itemize}
\end{corrige}

\begin{corrige}[Exercise 2 — True/False with justification]
\begin{enumerate}
\item \textbf{True.} An AM wave contains the carrier at $f_c$ and two sidebands extending from $f_c - f_m$ to $f_c + f_m$, where $f_m$ is the maximum modulating frequency. Hence the bandwidth is
\[ B = (f_c + f_m) - (f_c - f_m) = 2f_m. \]
\item \textbf{False.} It is the opposite: FM is \emph{less} susceptible to noise. Noise and interference mainly add unwanted \emph{amplitude} variations; in FM the information is carried by frequency variations, and amplitude limiters in the receiver remove the noise. In AM, noise directly corrupts the amplitude that carries the information.
\item \textbf{False.} A \emph{half-wave} dipole has a total length $L = \lambda/2$, i.e. \textbf{half} a wavelength, not a full wavelength.
\item \textbf{False.} The Nyquist theorem requires the sampling frequency to be \emph{at least twice} the highest frequency present: $f_s \geq 2 f_{\max}$. Sampling at exactly $f_{\max}$ would cause severe loss of information (aliasing).
\end{enumerate}
\end{corrige}

\begin{corrige}[Exercise 3 — Definitions]
\begin{enumerate}
\item \textbf{Modulation index (AM):} the ratio of the amplitude of the modulating signal to the amplitude of the carrier,
\[ m = \frac{A_m}{A_c}, \]
equivalently $m = \dfrac{A_{\max} - A_{\min}}{A_{\max} + A_{\min}}$ from the envelope. It is dimensionless; $m \leq 1$ for distortion-free AM.
\item \textbf{Sidebands:} the bands of new frequencies produced by modulation, symmetrically placed on either side of the carrier: the \emph{upper sideband} (USB) around $f_c + f_m$ and the \emph{lower sideband} (LSB) around $f_c - f_m$. They carry all the information of the modulating signal.
\item \textbf{Quantisation error:} the difference between the actual value of the analogue sample and the nearest quantisation level (digital value) assigned to it. It is at most half a quantisation step, $\pm q/2$, and appears as quantisation noise.
\item \textbf{Geostationary satellite:} a satellite in a circular equatorial orbit at an altitude of about $36\,000\ \mathrm{km}$, with an orbital period equal to the Earth's rotation period (24 h). It therefore appears fixed above a point on the equator, allowing fixed ground aerials to communicate with it continuously.
\end{enumerate}
\end{corrige}

\begin{corrige}[Exercise 4 — Direct calculation: Wavelength and aerial length]
\textbf{Given:} $f = 95.5\ \mathrm{MHz} = 9.55\times 10^{7}\ \mathrm{Hz}$; $c = 3.00\times 10^{8}\ \mathrm{m\,s^{-1}}$.

\begin{enumerate}
\item Wavelength:
\[ \lambda = \frac{c}{f} = \frac{3.00\times 10^{8}}{9.55\times 10^{7}} = 3.14\ \mathrm{m}. \]
\item A half-wave dipole has length $L = \lambda/2$:
\[ L = \frac{3.14}{2} = 1.57\ \mathrm{m}. \]
\end{enumerate}
This is why FM radio aerials (telescopic rods on radios, or rooftop aerials) are of the order of a metre long.
\end{corrige}

\begin{corrige}[Exercise 5 — AM modulation index and overmodulation]
\textbf{Given:} $A_{\max} = 12\ \mathrm{V}$, $A_{\min} = 4\ \mathrm{V}$.

\begin{enumerate}
\item The modulation index is obtained from the envelope:
\[ m = \frac{A_{\max} - A_{\min}}{A_{\max} + A_{\min}} = \frac{12 - 4}{12 + 4} = \frac{8}{16} = 0.50. \]
\item The signal is \textbf{not overmodulated}, since $m = 0.50 < 1$. Overmodulation occurs when $m > 1$; the envelope then crosses zero and the message is distorted on demodulation.
\item The carrier amplitude is $A_c = \dfrac{A_{\max}+A_{\min}}{2} = 8\ \mathrm{V}$ (consistent with the data). Each sideband has peak amplitude
\[ A_{\text{SB}} = \frac{m A_c}{2} = \frac{0.50 \times 8}{2} = 2.0\ \mathrm{V}. \]
\end{enumerate}
\begin{attention}[Common mistake]
Do not write $A_{\text{SB}} = mA_c$: the modulating signal's energy is shared equally between the \emph{two} sidebands, so each gets $mA_c/2$.
\end{attention}
\end{corrige}

\begin{corrige}[Exercise 6 — Tuning circuit frequency range]
The resonant frequency of an $LC$ circuit is
\[ f_0 = \frac{1}{2\pi\sqrt{LC}}. \]
Since $f_0 \propto 1/\sqrt{C}$, the \emph{maximum} frequency corresponds to $C_{\min}$ and vice versa.

\begin{enumerate}
\item Maximum frequency ($C_{\min} = 20\ \mathrm{pF}$):
\[ f_{\max} = \frac{1}{2\pi\sqrt{180\times 10^{-6}\times 20\times 10^{-12}}}
= \frac{1}{2\pi\sqrt{3.6\times 10^{-15}}}
\approx 2.65\ \mathrm{MHz}. \]
Minimum frequency ($C_{\max} = 300\ \mathrm{pF}$):
\[ f_{\min} = \frac{1}{2\pi\sqrt{180\times 10^{-6}\times 300\times 10^{-12}}}
= \frac{1}{2\pi\sqrt{5.4\times 10^{-14}}}
\approx 0.685\ \mathrm{MHz} = 685\ \mathrm{kHz}. \]
\item The range $0.685$--$2.65\ \mathrm{MHz}$ covers the \textbf{medium wave (MW)} band (roughly $0.53$--$1.6\ \mathrm{MHz}$) and extends into the lower \textbf{short wave (SW)} band. It does not reach LW ($< 300\ \mathrm{kHz}$) or VHF/FM ($88$--$108\ \mathrm{MHz}$).
\item Resonance curve: the current $I$ is maximum at $f_0$; the bandwidth $\Delta f$ is measured between the half-power points ($I = I_{\max}/\sqrt{2}$), and $Q = f_0/\Delta f$.
\end{enumerate}

\begin{popfigure}
\begin{center}
\begin{tikzpicture}
\begin{axis}[
width=10cm, height=6cm,
xlabel={$f$}, ylabel={$I/I_{\max}$},
xmin=0, xmax=10, ymin=0, ymax=1.1,
xtick={5}, xticklabels={$f_0$},
ytick={1, 0.707}, yticklabels={$1$, $1/\sqrt{2}$},
axis lines=left, thick,
domain=0:10, samples=200]
\addplot[very thick, popblue] {1/sqrt(max(0,1 + 4*(x-5)^2))};
\addplot[dashed, popmuted] coordinates {(0,0.707) (4.5,0.707)};
\addplot[dashed, popmuted] coordinates {(0,0.707) (5.5,0.707)};
\addplot[dashed, popmuted] coordinates {(5,0) (5,1)};
\draw[poporange, thick, <->] (axis cs:4.5,0.85) -- (axis cs:5.5,0.85)
node[midway, above, popdark]{$\Delta f$};
\node[popdark, align=center] at (axis cs:7.5,0.5) {$Q = \dfrac{f_0}{\Delta f}$};
\end{axis}
\end{tikzpicture}
\end{center}
\legende{Resonance curve of the tuning circuit: sharpness set by the quality factor $Q$. Half-power points at $f_0 \pm \Delta f/2$.}
\end{popfigure}
\end{corrige}

\begin{corrige}[Exercise 7 — AM frequency spectrum and bandwidth]
\textbf{Given:} $f_c = 1\ \mathrm{MHz} = 1000\ \mathrm{kHz}$, $f_m(\max) = 4.5\ \mathrm{kHz}$.

\begin{enumerate}
\item The spectrum contains the carrier at $1000\ \mathrm{kHz}$, a lower sideband from $1000 - 4.5 = 995.5\ \mathrm{kHz}$ up to the carrier, and an upper sideband from the carrier to $1000 + 4.5 = 1004.5\ \mathrm{kHz}$.
\end{enumerate}

\begin{popfigure}
\begin{center}
\begin{tikzpicture}
\begin{axis}[
width=11cm, height=5.5cm,
xlabel={$f$ (kHz)}, ylabel={amplitude},
xmin=990, xmax=1010, ymin=0, ymax=1.3,
xtick={995.5, 1000, 1004.5},
xticklabels={$995.5$, $1000$, $1004.5$},
ytick=\empty, axis lines=left, thick]
\addplot[very thick, popblue] coordinates {(1000,0) (1000,1)};
\addplot[very thick, poporange, fill=poporange, fill opacity=0.3]
coordinates {(995.5,0) (997.75,0.45) (1000,0) };
\addplot[very thick, poporange, fill=poporange, fill opacity=0.3]
coordinates {(1000,0) (1002.25,0.45) (1004.5,0)};
\node[popblue] at (axis cs:1000,1.1) {carrier};
\node[popdark] at (axis cs:996.6,0.62) {LSB};
\node[popdark] at (axis cs:1003.4,0.62) {USB};
\draw[popdark, thick, <->] (axis cs:995.5,0.9) -- (axis cs:1004.5,0.9)
node[midway, above]{$B = 9\ \mathrm{kHz}$};
\end{axis}
\end{tikzpicture}
\end{center}
\legende{Frequency spectrum of the AM wave: carrier and two sidebands.}
\end{popfigure}

\begin{enumerate}
\setcounter{enumi}{1}
\item Channel bandwidth:
\[ B_{\text{AM}} = 2 f_m = 2 \times 4.5 = 9.0\ \mathrm{kHz}. \]
\item Carson's rule for FM (assuming standard broadcast deviation $\Delta f = 75\ \mathrm{kHz}$):
\[ B_{\text{FM}} \approx 2(\Delta f + f_m) = 2(75 + 4.5) = 159\ \mathrm{kHz}. \]
\end{enumerate}
\textbf{Comment:} FM needs nearly 18 times more bandwidth than AM for the same audio signal — the price paid for its superior noise immunity.
\end{corrige}

\begin{corrige}[Exercise 8 — Digital audio parameters]
\textbf{Given:} $f_s = 44.1\ \mathrm{kHz}$, $n = 16$ bits, 2 channels (stereo).

\begin{enumerate}
\item Number of quantisation levels:
\[ N = 2^n = 2^{16} = 65\,536 \text{ levels}. \]
\item Bit rate for stereo:
\[ R = n \times f_s \times 2 = 16 \times 44.1\times 10^{3} \times 2
= 1.4112\times 10^{6}\ \mathrm{bit\,s^{-1}} = 1411\ \mathrm{kbps}. \]
(This is exactly the bit rate of audio CD.)
\item Nyquist criterion: $f_{\max} = f_s / 2 = 44.1/2 = 22.05\ \mathrm{kHz}$. This comfortably covers the human hearing range (up to about $20\ \mathrm{kHz}$).
\item Any frequency component \emph{above} $f_s/2$ present at the sampler input is ``folded back'' into the audio band as a false low frequency (\cle{aliasing}); once sampled, this corruption cannot be removed. A \textbf{low-pass anti-aliasing filter} with cut-off at about $22\ \mathrm{kHz}$ is therefore placed before the ADC to eliminate all components that could alias, guaranteeing faithful reconstruction.
\end{enumerate}
\end{corrige}

\begin{corrige}[Exercise 9 — Propagation modes and reception behind a hill (15 marks)]
\textbf{(a) The three propagation modes [6 marks — 2 per mode: name 1, description 1]}

\begin{itemize}
\item \textbf{Ground (surface) wave:} the wave travels along the Earth's surface, following its curvature by diffraction around obstacles. Effective at low frequencies (LW, MW); range limited to a few hundred km by ground absorption.
\item \textbf{Sky wave:} the wave is directed upwards and refracted/reflected back to Earth by the ionosphere. Used in HF (short wave, $3$--$30\ \mathrm{MHz}$); enables intercontinental communication by successive hops, but is variable with time of day and solar activity.
\item \textbf{Space (line-of-sight) wave:} the wave travels in a straight line from transmitting to receiving aerial, including reflection from the ground. Used at VHF/UHF and above (FM radio, TV, satellite links); range limited by the horizon and obstacles.
\end{itemize}

\textbf{(b) Why MW survives behind the hill but FM does not [6 marks]}

MW ($\sim 1\ \mathrm{MHz}$) has $\lambda = c/f \approx 300\ \mathrm{m}$, comparable to the size of the hill: the wave \textbf{diffracts} strongly around and over the hill and also propagates as a ground wave hugging the terrain, so it reaches the car.

VHF/FM ($\sim 100\ \mathrm{MHz}$) has $\lambda \approx 3\ \mathrm{m}$, far smaller than the hill: diffraction is negligible, the wave travels essentially in a straight line, and the hill casts a \textbf{radio shadow} in which no signal is received.

\begin{popfigure}
\begin{center}
\begin{tikzpicture}[scale=0.9]
\draw[thick, popdark] (-7,0) -- (7,0);
\fill[popgreenL] (1,0) .. controls (1.6,2.6) and (2.6,2.8) .. (3.4,0) -- cycle;
\draw[popdark] (1,0) .. controls (1.6,2.6) and (2.6,2.8) .. (3.4,0);
\node at (2.2,0.7) {hill};
\draw[very thick, popblue] (-6,1.8) .. controls (-2,1.6) and (-0.5,3.4) .. (2.2,3.2)
.. controls (4,3.0) and (5,1.2) .. (6,0.6);
\node[popblue] at (-4.6,2.3) {MW: diffracted};
\draw[very thick, poporange, dashed] (-6,3.4) -- (1.35,3.4);
\draw[very thick, poporange, dashed] (2.9,3.4) -- (6,3.4);
\node[poporange] at (-4.2,3.8) {FM: straight line (blocked)};
\fill[poppurple] (-6,0) rectangle (-5.7,1.9);
\node[below] at (-5.85,0) {TX};
\draw[popmuted, thick] (5.6,0) rectangle (6.4,0.5);
\node[below] at (6,0) {car};
\end{tikzpicture}
\end{center}
\legende{Long-wave MW signal diffracts over the hill; the short-wavelength FM signal is blocked.}
\end{popfigure}

\emph{Examiner expectations:} a labelled diagram, the comparison of $\lambda$ with obstacle size, and the key word \textbf{diffraction}.

\textbf{(c) Sky wave: one advantage, one disadvantage [3 marks]}
\begin{itemize}
\item \emph{Advantage:} very long-distance (intercontinental) communication possible without satellites or repeaters, using ionospheric reflection.
\item \emph{Disadvantage:} unreliable and variable — the ionosphere changes with time of day, season and solar activity, causing fading and sometimes total loss of signal.
\end{itemize}
\end{corrige}

\begin{corrige}[Exercise 10 — Comparison of AM and FM; TV sound transmission (15 marks)]
\textbf{(a) Comparison [8 marks — 2 per heading]}

\begin{center}
\begin{tabularx}{\linewidth}{|l|X|X|}
\hline
\rowcolor{popblueL} \textbf{Criterion} & \textbf{AM} & \textbf{FM} \\
\hline
Noise susceptibility & High: noise adds amplitude variations which directly corrupt the message & Low: information is in frequency; amplitude limiters remove noise \\
\hline
Bandwidth & Narrow: $B = 2f_m$ (e.g. $9\ \mathrm{kHz}$) & Wide: $B \approx 2(\Delta f + f_m)$ (e.g. $180\ \mathrm{kHz}$) \\
\hline
Range (same power) & Greater: ground/sky waves at LW--MW--HF cover long distances & Smaller: VHF space wave limited roughly to line of sight \\
\hline
Circuit complexity & Simple, cheap transmitter and receiver (diode detector) & More complex, costlier (limiter, discriminator/PLL) \\
\hline
\end{tabularx}
\end{center}

\textbf{(b) Why CRTV uses FM for TV sound [4 marks]}
\begin{itemize}
\item Television sound must be of \textbf{high fidelity}; FM's superior \textbf{noise immunity} gives clear, hiss-free audio, whereas AM would crackle with every interference (motors, lightning, electrical appliances).
\item TV already operates at VHF/UHF where wide channels are available, so FM's \textbf{large bandwidth} is affordable.
\item The wide FM bandwidth accommodates the full audio range ($15\ \mathrm{kHz}$) with a large deviation, giving a high signal-to-noise ratio.
\item AM's only real advantages (long range, simplicity) are irrelevant here since TV is broadcast from local transmitters within line of sight.
\end{itemize}

\textbf{(c) Carson's rule [3 marks]}
\[ B \approx 2(\Delta f + f_m) = 2(75 + 15) = 2 \times 90 = 180\ \mathrm{kHz}. \]
\emph{Mark scheme:} correct formula [1], correct substitution [1], answer $180\ \mathrm{kHz}$ with unit [1].
\end{corrige}

\begin{corrige}[Exercise 11 — Satellite communication delay (12 marks)]
\textbf{Given:} altitude $h = 36\,000\ \mathrm{km} = 3.6\times 10^{7}\ \mathrm{m}$; $c = 3.00\times 10^{8}\ \mathrm{m\,s^{-1}}$.

\textbf{(a) One-way delay [4 marks]}

The path is Douala $\to$ satellite $\to$ London, i.e. two hops each of length approximately $h$ (standard approximation for geostationary links):
\[ d = 2h = 7.2\times 10^{7}\ \mathrm{m}. \]
\[ t = \frac{d}{c} = \frac{7.2\times 10^{7}}{3.00\times 10^{8}} = 0.24\ \mathrm{s}. \]
\emph{Mark scheme:} identifying two hops [1], $d = 2h$ [1], formula $t = d/c$ [1], answer $0.24\ \mathrm{s}$ [1].

\textbf{(b) Round-trip delay [2 marks]}

A reply must make the same journey back:
\[ t_{\text{RT}} = 2 \times 0.24 = 0.48\ \mathrm{s} \approx 0.5\ \mathrm{s}. \]

\textbf{(c) Effect on conversation [3 marks]}

A pause of about half a second between speaking and hearing the reply is clearly perceptible. Speakers begin to talk over each other, hesitate, or repeat themselves, because the natural rhythm of conversation expects replies within a small fraction of a second. The conversation feels unnatural and tiring.

\textbf{(d) Mitigation [3 marks]}

Any one of:
\begin{itemize}
\item \textbf{Echo cancellers} and digital signal processing at the exchanges, so that the delayed echo of one's own voice (the most disturbing effect) is suppressed.
\item Routing the call via \textbf{undersea fibre-optic cables} (or low-Earth-orbit satellite constellations) instead of a geostationary satellite, drastically shortening the path and the delay.
\item Disciplined ``push-to-talk'' style protocols where each speaker waits for a clear end of transmission.
\end{itemize}
\emph{Examiner expectation:} one well-explained technical method suffices for full marks.
\end{corrige}

\newpage

\coursec{Summary sheet}

\begin{popfigure}
\begin{tikzpicture}[font=\small,
  root/.style={rectangle, rounded corners, draw=popdark, fill=popdark, text=white, align=center, font=\bfseries\small, minimum width=3.4cm, minimum height=1cm},
  br/.style={rectangle, rounded corners, draw=#1, fill=#1!15, text=popink, align=center, font=\small, minimum width=3.2cm, minimum height=0.8cm},
  leaf/.style={rectangle, rounded corners, draw=#1, fill=white, text=popink, align=center, font=\scriptsize, minimum width=2.9cm, minimum height=0.7cm}
]
\node[root] (R) at (0,0) {Communication};
\node[br=popblue] (A) at (-6.4,2.6) {Representing information};
\node[br=popgreen] (B) at (0,3.4) {Radio wave propagation};
\node[br=poporange] (C) at (6.4,2.6) {Aerials \& tuning};
\node[br=popgold] (D) at (-6.4,-2.6) {Modulation};
\node[br=poppurple] (E) at (0,-3.4) {Digital conversion};
\node[br=poppink] (F) at (6.4,-2.6) {Channels};
\node[leaf=popblue] (A1) at (-9.8,4.4) {Analogue: continuous};
\node[leaf=popblue] (A2) at (-6.4,4.8) {Digital: 0 and 1};
\node[leaf=popgreen] (B1) at (-2.6,5.2) {Ground wave};
\node[leaf=popgreen] (B2) at (0,5.6) {Sky wave};
\node[leaf=popgreen] (B3) at (2.6,5.2) {Space wave};
\node[leaf=poporange] (C1) at (6.4,4.8) {Transmitting, receiving};
\node[leaf=poporange] (C2) at (9.8,4.4) {$f_0=\dfrac{1}{2\pi\sqrt{LC}}$};
\node[leaf=popgold] (D1) at (-9.8,-4.4) {AM};
\node[leaf=popgold] (D2) at (-6.4,-4.8) {FM};
\node[leaf=popgold] (D3) at (-3.2,-4.4) {Demodulation};
\node[leaf=poppurple] (E1) at (-1.6,-5.2) {ADC};
\node[leaf=poppurple] (E2) at (1.6,-5.2) {DAC};
\node[leaf=poppink] (F1) at (6.4,-4.8) {Bandwidth, sidebands};
\node[leaf=poppink] (F2) at (9.8,-4.4) {Satellites, base stations};
\draw[popline, thick] (R) -- (A) (R) -- (B) (R) -- (C) (R) -- (D) (R) -- (E) (R) -- (F);
\draw[popblue] (A) -- (A1) (A) -- (A2);
\draw[popgreen] (B) -- (B1) (B) -- (B2) (B) -- (B3);
\draw[poporange] (C) -- (C1) (C) -- (C2);
\draw[popgold] (D) -- (D1) (D) -- (D2) (D) -- (D3);
\draw[poppurple] (E) -- (E1) (E) -- (E2);
\draw[poppink] (F) -- (F1) (F) -- (F2);
\end{tikzpicture}
\legende{Concept map of the chapter Communication: the six themes and their key ideas.}
\end{popfigure}

\begin{aretenir}[Communication --- Summary Sheet]
\textbf{1. Key definitions}
\begin{itemize}
  \item \cle{Analogue signal}: varies continuously with time, taking any value within a range (e.g.\ the voltage from a microphone).
  \item \cle{Digital signal}: takes only discrete values, usually two levels (0 and 1); information is coded in binary.
  \item \cle{Carrier wave}: high-frequency radio wave that ``carries'' the low-frequency information signal.
  \item \cle{Modulation}: process of superimposing the information signal onto the carrier by varying one of its characteristics (amplitude or frequency).
  \item \cle{Demodulation} (detection): recovering the original information signal from the modulated carrier at the receiver.
  \item \cle{Bandwidth}: range of frequencies occupied by a signal or passed by a channel, $\Delta f = f_{\max}-f_{\min}$.
  \item \cle{Sidebands}: bands of frequencies above and below the carrier frequency $f_c$, produced by modulation.
  \item \cle{Base station}: fixed transmitter--receiver linking mobile phones to the telephone network within its cell.
\end{itemize}

\textbf{2. Laws and formulas to know}
\begin{itemize}
  \item Wave equation: $c = f\lambda$ with $c = \SI{3.00e8}{m.s^{-1}}$ for all electromagnetic waves in air or vacuum.
  \item Resonance of the $LC$ tuning circuit: $f_0 = \dfrac{1}{2\pi\sqrt{LC}}$; at resonance the response is maximum (peak of the resonance curve). The sharper the peak, the better the \cle{selectivity}.
  \item AM: carrier amplitude varies with the signal; modulation index $m = \dfrac{A_{\text{signal}}}{A_{\text{carrier}}}$; keep $m \le 1$ to avoid distortion. Bandwidth $= 2f_m$ (two sidebands $f_c \pm f_m$).
  \item FM: carrier frequency varies with the signal; amplitude constant; bandwidth larger than AM.
  \item Sampling: to digitise correctly, sample at $f_s \ge 2f_{\max}$ (sampling criterion).
\end{itemize}

\textbf{3. Propagation of radio waves}
\begin{itemize}
  \item \cle{Ground (surface) wave}: follows the Earth's surface; low frequencies; long range over land/sea (LW, MW).
  \item \cle{Sky wave}: reflected by the ionosphere back to Earth; short waves; very long range, but varies with time of day.
  \item \cle{Space wave}: straight line-of-sight between aerials; VHF/UHF (FM radio, TV, satellites); limited by the horizon, needs relay masts or satellites.
\end{itemize}

\textbf{4. AM versus FM; analogue versus digital}
\begin{itemize}
  \item AM: simpler, cheaper, longer range, smaller bandwidth; but noisy and lower sound quality.
  \item FM: noise-resistant, better quality; but shorter range (space wave), wider bandwidth, costlier.
  \item Digital advantages: immune to noise, easy regeneration, storage and encryption; analogue is simpler and needs no conversion.
\end{itemize}

\textbf{5. Methods}
\begin{itemize}
  \item To find the station received by a tuning circuit: compute $f_0 = \dfrac{1}{2\pi\sqrt{LC}}$ and compare with the carrier frequencies.
  \item To sketch AM: draw the carrier, then vary its amplitude following the signal envelope; for FM keep amplitude constant and vary the spacing of the oscillations.
  \item ADC chain: sampling $\to$ quantisation $\to$ binary coding; DAC reverses it.
\end{itemize}

\textbf{6. Common mistakes}
\begin{itemize}
  \item Confusing carrier frequency $f_c$ with signal frequency $f_m$; the sidebands are at $f_c \pm f_m$, not at $f_m$.
  \item Forgetting the $2\pi$ in $f_0 = \dfrac{1}{2\pi\sqrt{LC}}$.
  \item Saying FM changes the amplitude --- it changes the \emph{frequency}.
  \item Mixing up sky wave (ionosphere) and space wave (line of sight).
\end{itemize}

\textbf{7. GCE tips}
\begin{itemize}
  \item Always give units (Hz, m) and use $c = \SI{3.00e8}{m.s^{-1}}$; quote 3 significant figures.
  \item In structured questions, define terms before comparing (e.g.\ state what modulation is before contrasting AM and FM).
  \item Be ready to draw and label a resonance curve and an AM waveform --- frequent exam items.
\end{itemize}
\end{aretenir}

\findecours

\end{document}
